\begin{python}
from scripts.calc_9 import *
\end{python}

\section{Приближенный расчет реакции по спектру}

Связь между спектральной плотностью сигнала и его временным представлением устанавливается с помощью обратного преобразования Фурье^
\begin{equation}
    i_2(t) = \frac{1}{\pi} \int_{0}^{\infty} A_2(\omega) \cos(\omega t + \Phi_2(\omega)) d\omega.
\end{equation}

Аналитическое вычисление данного интеграла затруднительно. Приближенный расчет осуществляется путем замены бесконечного верхнего предела интегрирования на конечное значение $\omega_{\text{max}}$ и применения методов численного интегрирования.

Спектр реакции на меандр ${I}_{2, \text{меандр}}(j\omega)$ вычисляется через спектр одиночного импульса ${I}_{1, \text{имп}}(j\omega)$ и передаточную функцию цепи ${H}(j\omega)$:
\begin{equation}
    {I}_{2, \text{меандр}}(j\omega) = {H}(j\omega) \cdot {I}_{1, \text{имп}}(j\omega) \cdot \left(1 - e^{-j\omega t_{\text{и}}}\right)
\end{equation}

Продемонстрируем вычисление для двух длительностей меандра в контрольных точках:
\begin{itemize}
    \item Для $t_{\text{и}1} = \pv[.2]{n['ti1']}$ в момент $t = \pv[.2]{n['ti1']*1.5}$:
    \begin{align*}
        i_{2\text{ (прибл)}} &= \pv[.4]{ex9['i2_approx_demo1']}  \\
        i_{2\text{ (точн)}} &= \pv[.4]{ex9['i2_exact_demo1']} 
    \end{align*}
    \item Для $t_{\text{и}2} = \pv[.2]{n['ti2']}$ в момент $t = \pv[.2]{n['ti2']*1.5}$:
    \begin{align*}
        i_{2\text{ (прибл)}} &= \pv[.4]{ex9['i2_approx_demo2']}  \\
        i_{2\text{ (точн)}} &= \pv[.4]{ex9['i2_exact_demo2']} 
    \end{align*}
\end{itemize}

Сравнительные графики реакций цепи, полученные точным операторным и приближенным спектральным методами, представлены на \cref{fig:plot_approx_reaction_1,fig:plot_approx_reaction_2}.

\begin{figure}[H]
    \centering
    %% Creator: Matplotlib, PGF backend
%%
%% To include the figure in your LaTeX document, write
%%   \input{<filename>.pgf}
%%
%% Make sure the required packages are loaded in your preamble
%%   \usepackage{pgf}
%%
%% Also ensure that all the required font packages are loaded; for instance,
%% the lmodern package is sometimes necessary when using math font.
%%   \usepackage{lmodern}
%%
%% Figures using additional raster images can only be included by \input if
%% they are in the same directory as the main LaTeX file. For loading figures
%% from other directories you can use the `import` package
%%   \usepackage{import}
%%
%% and then include the figures with
%%   \import{<path to file>}{<filename>.pgf}
%%
%% Matplotlib used the following preamble
%%   \def\mathdefault#1{#1}
%%   \everymath=\expandafter{\the\everymath\displaystyle}
%%   \IfFileExists{scrextend.sty}{
%%     \usepackage[fontsize=12.000000pt]{scrextend}
%%   }{
%%     \renewcommand{\normalsize}{\fontsize{12.000000}{14.400000}\selectfont}
%%     \normalsize
%%   }
%%   \usepackage{fontspec}\setmainfont{Times New Roman}\usepackage[english,russian]{babel}\usepackage{amsmath}
%%   \ifdefined\pdftexversion\else  % non-pdftex case.
%%     \usepackage{fontspec}
%%     \setmainfont{times.ttf}[Path=\detokenize{/usr/local/share/fonts/truetype/times-new-roman/}]
%%     \setsansfont{DejaVuSans.ttf}[Path=\detokenize{/usr/share/matplotlib/mpl-data/fonts/ttf/}]
%%     \setmonofont{DejaVuSansMono.ttf}[Path=\detokenize{/usr/share/matplotlib/mpl-data/fonts/ttf/}]
%%   \fi
%%   \makeatletter\@ifpackageloaded{underscore}{}{\usepackage[strings]{underscore}}\makeatother
%%
\begingroup%
\makeatletter%
\begin{pgfpicture}%
\pgfpathrectangle{\pgfpointorigin}{\pgfqpoint{6.642382in}{3.740000in}}%
\pgfusepath{use as bounding box, clip}%
\begin{pgfscope}%
\pgfsetbuttcap%
\pgfsetmiterjoin%
\definecolor{currentfill}{rgb}{1.000000,1.000000,1.000000}%
\pgfsetfillcolor{currentfill}%
\pgfsetlinewidth{0.000000pt}%
\definecolor{currentstroke}{rgb}{1.000000,1.000000,1.000000}%
\pgfsetstrokecolor{currentstroke}%
\pgfsetdash{}{0pt}%
\pgfpathmoveto{\pgfqpoint{0.000000in}{0.000000in}}%
\pgfpathlineto{\pgfqpoint{6.642382in}{0.000000in}}%
\pgfpathlineto{\pgfqpoint{6.642382in}{3.740000in}}%
\pgfpathlineto{\pgfqpoint{0.000000in}{3.740000in}}%
\pgfpathlineto{\pgfqpoint{0.000000in}{0.000000in}}%
\pgfpathclose%
\pgfusepath{fill}%
\end{pgfscope}%
\begin{pgfscope}%
\pgfsetbuttcap%
\pgfsetmiterjoin%
\definecolor{currentfill}{rgb}{1.000000,1.000000,1.000000}%
\pgfsetfillcolor{currentfill}%
\pgfsetlinewidth{0.000000pt}%
\definecolor{currentstroke}{rgb}{0.000000,0.000000,0.000000}%
\pgfsetstrokecolor{currentstroke}%
\pgfsetstrokeopacity{0.000000}%
\pgfsetdash{}{0pt}%
\pgfpathmoveto{\pgfqpoint{0.781757in}{0.456901in}}%
\pgfpathlineto{\pgfqpoint{5.860625in}{0.456901in}}%
\pgfpathlineto{\pgfqpoint{5.860625in}{3.273333in}}%
\pgfpathlineto{\pgfqpoint{0.781757in}{3.273333in}}%
\pgfpathlineto{\pgfqpoint{0.781757in}{0.456901in}}%
\pgfpathclose%
\pgfusepath{fill}%
\end{pgfscope}%
\begin{pgfscope}%
\pgfpathrectangle{\pgfqpoint{0.781757in}{0.456901in}}{\pgfqpoint{5.078868in}{2.816432in}}%
\pgfusepath{clip}%
\pgfsetbuttcap%
\pgfsetroundjoin%
\pgfsetlinewidth{0.501875pt}%
\definecolor{currentstroke}{rgb}{0.501961,0.501961,0.501961}%
\pgfsetstrokecolor{currentstroke}%
\pgfsetstrokeopacity{0.400000}%
\pgfsetdash{{1.850000pt}{0.800000pt}}{0.000000pt}%
\pgfpathmoveto{\pgfqpoint{0.781757in}{0.456901in}}%
\pgfpathlineto{\pgfqpoint{0.781757in}{3.273333in}}%
\pgfusepath{stroke}%
\end{pgfscope}%
\begin{pgfscope}%
\pgfsetbuttcap%
\pgfsetroundjoin%
\definecolor{currentfill}{rgb}{1.000000,1.000000,1.000000}%
\pgfsetfillcolor{currentfill}%
\pgfsetlinewidth{0.803000pt}%
\definecolor{currentstroke}{rgb}{0.000000,0.000000,0.000000}%
\pgfsetstrokecolor{currentstroke}%
\pgfsetdash{}{0pt}%
\pgfsys@defobject{currentmarker}{\pgfqpoint{0.000000in}{0.000000in}}{\pgfqpoint{0.000000in}{0.048611in}}{%
\pgfpathmoveto{\pgfqpoint{0.000000in}{0.000000in}}%
\pgfpathlineto{\pgfqpoint{0.000000in}{0.048611in}}%
\pgfusepath{stroke,fill}%
}%
\begin{pgfscope}%
\pgfsys@transformshift{0.781757in}{0.456901in}%
\pgfsys@useobject{currentmarker}{}%
\end{pgfscope}%
\end{pgfscope}%
\begin{pgfscope}%
\definecolor{textcolor}{rgb}{0.000000,0.000000,0.000000}%
\pgfsetstrokecolor{textcolor}%
\pgfsetfillcolor{textcolor}%
\pgftext[x=0.781757in,y=0.408290in,,top]{\color{textcolor}{\rmfamily\fontsize{12.000000}{14.400000}\selectfont\catcode`\^=\active\def^{\ifmmode\sp\else\^{}\fi}\catcode`\%=\active\def%{\%}$\mathdefault{0}$}}%
\end{pgfscope}%
\begin{pgfscope}%
\pgfpathrectangle{\pgfqpoint{0.781757in}{0.456901in}}{\pgfqpoint{5.078868in}{2.816432in}}%
\pgfusepath{clip}%
\pgfsetbuttcap%
\pgfsetroundjoin%
\pgfsetlinewidth{0.501875pt}%
\definecolor{currentstroke}{rgb}{0.501961,0.501961,0.501961}%
\pgfsetstrokecolor{currentstroke}%
\pgfsetstrokeopacity{0.400000}%
\pgfsetdash{{1.850000pt}{0.800000pt}}{0.000000pt}%
\pgfpathmoveto{\pgfqpoint{1.808699in}{0.456901in}}%
\pgfpathlineto{\pgfqpoint{1.808699in}{3.273333in}}%
\pgfusepath{stroke}%
\end{pgfscope}%
\begin{pgfscope}%
\pgfsetbuttcap%
\pgfsetroundjoin%
\definecolor{currentfill}{rgb}{1.000000,1.000000,1.000000}%
\pgfsetfillcolor{currentfill}%
\pgfsetlinewidth{0.803000pt}%
\definecolor{currentstroke}{rgb}{0.000000,0.000000,0.000000}%
\pgfsetstrokecolor{currentstroke}%
\pgfsetdash{}{0pt}%
\pgfsys@defobject{currentmarker}{\pgfqpoint{0.000000in}{0.000000in}}{\pgfqpoint{0.000000in}{0.048611in}}{%
\pgfpathmoveto{\pgfqpoint{0.000000in}{0.000000in}}%
\pgfpathlineto{\pgfqpoint{0.000000in}{0.048611in}}%
\pgfusepath{stroke,fill}%
}%
\begin{pgfscope}%
\pgfsys@transformshift{1.808699in}{0.456901in}%
\pgfsys@useobject{currentmarker}{}%
\end{pgfscope}%
\end{pgfscope}%
\begin{pgfscope}%
\definecolor{textcolor}{rgb}{0.000000,0.000000,0.000000}%
\pgfsetstrokecolor{textcolor}%
\pgfsetfillcolor{textcolor}%
\pgftext[x=1.808699in,y=0.408290in,,top]{\color{textcolor}{\rmfamily\fontsize{12.000000}{14.400000}\selectfont\catcode`\^=\active\def^{\ifmmode\sp\else\^{}\fi}\catcode`\%=\active\def%{\%}$\mathdefault{20}$}}%
\end{pgfscope}%
\begin{pgfscope}%
\pgfpathrectangle{\pgfqpoint{0.781757in}{0.456901in}}{\pgfqpoint{5.078868in}{2.816432in}}%
\pgfusepath{clip}%
\pgfsetbuttcap%
\pgfsetroundjoin%
\pgfsetlinewidth{0.501875pt}%
\definecolor{currentstroke}{rgb}{0.501961,0.501961,0.501961}%
\pgfsetstrokecolor{currentstroke}%
\pgfsetstrokeopacity{0.400000}%
\pgfsetdash{{1.850000pt}{0.800000pt}}{0.000000pt}%
\pgfpathmoveto{\pgfqpoint{2.835640in}{0.456901in}}%
\pgfpathlineto{\pgfqpoint{2.835640in}{3.273333in}}%
\pgfusepath{stroke}%
\end{pgfscope}%
\begin{pgfscope}%
\pgfsetbuttcap%
\pgfsetroundjoin%
\definecolor{currentfill}{rgb}{1.000000,1.000000,1.000000}%
\pgfsetfillcolor{currentfill}%
\pgfsetlinewidth{0.803000pt}%
\definecolor{currentstroke}{rgb}{0.000000,0.000000,0.000000}%
\pgfsetstrokecolor{currentstroke}%
\pgfsetdash{}{0pt}%
\pgfsys@defobject{currentmarker}{\pgfqpoint{0.000000in}{0.000000in}}{\pgfqpoint{0.000000in}{0.048611in}}{%
\pgfpathmoveto{\pgfqpoint{0.000000in}{0.000000in}}%
\pgfpathlineto{\pgfqpoint{0.000000in}{0.048611in}}%
\pgfusepath{stroke,fill}%
}%
\begin{pgfscope}%
\pgfsys@transformshift{2.835640in}{0.456901in}%
\pgfsys@useobject{currentmarker}{}%
\end{pgfscope}%
\end{pgfscope}%
\begin{pgfscope}%
\definecolor{textcolor}{rgb}{0.000000,0.000000,0.000000}%
\pgfsetstrokecolor{textcolor}%
\pgfsetfillcolor{textcolor}%
\pgftext[x=2.835640in,y=0.408290in,,top]{\color{textcolor}{\rmfamily\fontsize{12.000000}{14.400000}\selectfont\catcode`\^=\active\def^{\ifmmode\sp\else\^{}\fi}\catcode`\%=\active\def%{\%}$\mathdefault{40}$}}%
\end{pgfscope}%
\begin{pgfscope}%
\pgfpathrectangle{\pgfqpoint{0.781757in}{0.456901in}}{\pgfqpoint{5.078868in}{2.816432in}}%
\pgfusepath{clip}%
\pgfsetbuttcap%
\pgfsetroundjoin%
\pgfsetlinewidth{0.501875pt}%
\definecolor{currentstroke}{rgb}{0.501961,0.501961,0.501961}%
\pgfsetstrokecolor{currentstroke}%
\pgfsetstrokeopacity{0.400000}%
\pgfsetdash{{1.850000pt}{0.800000pt}}{0.000000pt}%
\pgfpathmoveto{\pgfqpoint{3.862582in}{0.456901in}}%
\pgfpathlineto{\pgfqpoint{3.862582in}{3.273333in}}%
\pgfusepath{stroke}%
\end{pgfscope}%
\begin{pgfscope}%
\pgfsetbuttcap%
\pgfsetroundjoin%
\definecolor{currentfill}{rgb}{1.000000,1.000000,1.000000}%
\pgfsetfillcolor{currentfill}%
\pgfsetlinewidth{0.803000pt}%
\definecolor{currentstroke}{rgb}{0.000000,0.000000,0.000000}%
\pgfsetstrokecolor{currentstroke}%
\pgfsetdash{}{0pt}%
\pgfsys@defobject{currentmarker}{\pgfqpoint{0.000000in}{0.000000in}}{\pgfqpoint{0.000000in}{0.048611in}}{%
\pgfpathmoveto{\pgfqpoint{0.000000in}{0.000000in}}%
\pgfpathlineto{\pgfqpoint{0.000000in}{0.048611in}}%
\pgfusepath{stroke,fill}%
}%
\begin{pgfscope}%
\pgfsys@transformshift{3.862582in}{0.456901in}%
\pgfsys@useobject{currentmarker}{}%
\end{pgfscope}%
\end{pgfscope}%
\begin{pgfscope}%
\definecolor{textcolor}{rgb}{0.000000,0.000000,0.000000}%
\pgfsetstrokecolor{textcolor}%
\pgfsetfillcolor{textcolor}%
\pgftext[x=3.862582in,y=0.408290in,,top]{\color{textcolor}{\rmfamily\fontsize{12.000000}{14.400000}\selectfont\catcode`\^=\active\def^{\ifmmode\sp\else\^{}\fi}\catcode`\%=\active\def%{\%}$\mathdefault{60}$}}%
\end{pgfscope}%
\begin{pgfscope}%
\pgfpathrectangle{\pgfqpoint{0.781757in}{0.456901in}}{\pgfqpoint{5.078868in}{2.816432in}}%
\pgfusepath{clip}%
\pgfsetbuttcap%
\pgfsetroundjoin%
\pgfsetlinewidth{0.501875pt}%
\definecolor{currentstroke}{rgb}{0.501961,0.501961,0.501961}%
\pgfsetstrokecolor{currentstroke}%
\pgfsetstrokeopacity{0.400000}%
\pgfsetdash{{1.850000pt}{0.800000pt}}{0.000000pt}%
\pgfpathmoveto{\pgfqpoint{4.889523in}{0.456901in}}%
\pgfpathlineto{\pgfqpoint{4.889523in}{3.273333in}}%
\pgfusepath{stroke}%
\end{pgfscope}%
\begin{pgfscope}%
\pgfsetbuttcap%
\pgfsetroundjoin%
\definecolor{currentfill}{rgb}{1.000000,1.000000,1.000000}%
\pgfsetfillcolor{currentfill}%
\pgfsetlinewidth{0.803000pt}%
\definecolor{currentstroke}{rgb}{0.000000,0.000000,0.000000}%
\pgfsetstrokecolor{currentstroke}%
\pgfsetdash{}{0pt}%
\pgfsys@defobject{currentmarker}{\pgfqpoint{0.000000in}{0.000000in}}{\pgfqpoint{0.000000in}{0.048611in}}{%
\pgfpathmoveto{\pgfqpoint{0.000000in}{0.000000in}}%
\pgfpathlineto{\pgfqpoint{0.000000in}{0.048611in}}%
\pgfusepath{stroke,fill}%
}%
\begin{pgfscope}%
\pgfsys@transformshift{4.889523in}{0.456901in}%
\pgfsys@useobject{currentmarker}{}%
\end{pgfscope}%
\end{pgfscope}%
\begin{pgfscope}%
\definecolor{textcolor}{rgb}{0.000000,0.000000,0.000000}%
\pgfsetstrokecolor{textcolor}%
\pgfsetfillcolor{textcolor}%
\pgftext[x=4.889523in,y=0.408290in,,top]{\color{textcolor}{\rmfamily\fontsize{12.000000}{14.400000}\selectfont\catcode`\^=\active\def^{\ifmmode\sp\else\^{}\fi}\catcode`\%=\active\def%{\%}$\mathdefault{80}$}}%
\end{pgfscope}%
\begin{pgfscope}%
\pgfpathrectangle{\pgfqpoint{0.781757in}{0.456901in}}{\pgfqpoint{5.078868in}{2.816432in}}%
\pgfusepath{clip}%
\pgfsetbuttcap%
\pgfsetroundjoin%
\pgfsetlinewidth{0.501875pt}%
\definecolor{currentstroke}{rgb}{0.501961,0.501961,0.501961}%
\pgfsetstrokecolor{currentstroke}%
\pgfsetstrokeopacity{0.400000}%
\pgfsetdash{{1.850000pt}{0.800000pt}}{0.000000pt}%
\pgfpathmoveto{\pgfqpoint{1.038492in}{0.456901in}}%
\pgfpathlineto{\pgfqpoint{1.038492in}{3.273333in}}%
\pgfusepath{stroke}%
\end{pgfscope}%
\begin{pgfscope}%
\pgfsetbuttcap%
\pgfsetroundjoin%
\definecolor{currentfill}{rgb}{1.000000,1.000000,1.000000}%
\pgfsetfillcolor{currentfill}%
\pgfsetlinewidth{0.602250pt}%
\definecolor{currentstroke}{rgb}{0.000000,0.000000,0.000000}%
\pgfsetstrokecolor{currentstroke}%
\pgfsetdash{}{0pt}%
\pgfsys@defobject{currentmarker}{\pgfqpoint{0.000000in}{0.000000in}}{\pgfqpoint{0.000000in}{0.027778in}}{%
\pgfpathmoveto{\pgfqpoint{0.000000in}{0.000000in}}%
\pgfpathlineto{\pgfqpoint{0.000000in}{0.027778in}}%
\pgfusepath{stroke,fill}%
}%
\begin{pgfscope}%
\pgfsys@transformshift{1.038492in}{0.456901in}%
\pgfsys@useobject{currentmarker}{}%
\end{pgfscope}%
\end{pgfscope}%
\begin{pgfscope}%
\pgfpathrectangle{\pgfqpoint{0.781757in}{0.456901in}}{\pgfqpoint{5.078868in}{2.816432in}}%
\pgfusepath{clip}%
\pgfsetbuttcap%
\pgfsetroundjoin%
\pgfsetlinewidth{0.501875pt}%
\definecolor{currentstroke}{rgb}{0.501961,0.501961,0.501961}%
\pgfsetstrokecolor{currentstroke}%
\pgfsetstrokeopacity{0.400000}%
\pgfsetdash{{1.850000pt}{0.800000pt}}{0.000000pt}%
\pgfpathmoveto{\pgfqpoint{1.295228in}{0.456901in}}%
\pgfpathlineto{\pgfqpoint{1.295228in}{3.273333in}}%
\pgfusepath{stroke}%
\end{pgfscope}%
\begin{pgfscope}%
\pgfsetbuttcap%
\pgfsetroundjoin%
\definecolor{currentfill}{rgb}{1.000000,1.000000,1.000000}%
\pgfsetfillcolor{currentfill}%
\pgfsetlinewidth{0.602250pt}%
\definecolor{currentstroke}{rgb}{0.000000,0.000000,0.000000}%
\pgfsetstrokecolor{currentstroke}%
\pgfsetdash{}{0pt}%
\pgfsys@defobject{currentmarker}{\pgfqpoint{0.000000in}{0.000000in}}{\pgfqpoint{0.000000in}{0.027778in}}{%
\pgfpathmoveto{\pgfqpoint{0.000000in}{0.000000in}}%
\pgfpathlineto{\pgfqpoint{0.000000in}{0.027778in}}%
\pgfusepath{stroke,fill}%
}%
\begin{pgfscope}%
\pgfsys@transformshift{1.295228in}{0.456901in}%
\pgfsys@useobject{currentmarker}{}%
\end{pgfscope}%
\end{pgfscope}%
\begin{pgfscope}%
\pgfpathrectangle{\pgfqpoint{0.781757in}{0.456901in}}{\pgfqpoint{5.078868in}{2.816432in}}%
\pgfusepath{clip}%
\pgfsetbuttcap%
\pgfsetroundjoin%
\pgfsetlinewidth{0.501875pt}%
\definecolor{currentstroke}{rgb}{0.501961,0.501961,0.501961}%
\pgfsetstrokecolor{currentstroke}%
\pgfsetstrokeopacity{0.400000}%
\pgfsetdash{{1.850000pt}{0.800000pt}}{0.000000pt}%
\pgfpathmoveto{\pgfqpoint{1.551963in}{0.456901in}}%
\pgfpathlineto{\pgfqpoint{1.551963in}{3.273333in}}%
\pgfusepath{stroke}%
\end{pgfscope}%
\begin{pgfscope}%
\pgfsetbuttcap%
\pgfsetroundjoin%
\definecolor{currentfill}{rgb}{1.000000,1.000000,1.000000}%
\pgfsetfillcolor{currentfill}%
\pgfsetlinewidth{0.602250pt}%
\definecolor{currentstroke}{rgb}{0.000000,0.000000,0.000000}%
\pgfsetstrokecolor{currentstroke}%
\pgfsetdash{}{0pt}%
\pgfsys@defobject{currentmarker}{\pgfqpoint{0.000000in}{0.000000in}}{\pgfqpoint{0.000000in}{0.027778in}}{%
\pgfpathmoveto{\pgfqpoint{0.000000in}{0.000000in}}%
\pgfpathlineto{\pgfqpoint{0.000000in}{0.027778in}}%
\pgfusepath{stroke,fill}%
}%
\begin{pgfscope}%
\pgfsys@transformshift{1.551963in}{0.456901in}%
\pgfsys@useobject{currentmarker}{}%
\end{pgfscope}%
\end{pgfscope}%
\begin{pgfscope}%
\pgfpathrectangle{\pgfqpoint{0.781757in}{0.456901in}}{\pgfqpoint{5.078868in}{2.816432in}}%
\pgfusepath{clip}%
\pgfsetbuttcap%
\pgfsetroundjoin%
\pgfsetlinewidth{0.501875pt}%
\definecolor{currentstroke}{rgb}{0.501961,0.501961,0.501961}%
\pgfsetstrokecolor{currentstroke}%
\pgfsetstrokeopacity{0.400000}%
\pgfsetdash{{1.850000pt}{0.800000pt}}{0.000000pt}%
\pgfpathmoveto{\pgfqpoint{2.065434in}{0.456901in}}%
\pgfpathlineto{\pgfqpoint{2.065434in}{3.273333in}}%
\pgfusepath{stroke}%
\end{pgfscope}%
\begin{pgfscope}%
\pgfsetbuttcap%
\pgfsetroundjoin%
\definecolor{currentfill}{rgb}{1.000000,1.000000,1.000000}%
\pgfsetfillcolor{currentfill}%
\pgfsetlinewidth{0.602250pt}%
\definecolor{currentstroke}{rgb}{0.000000,0.000000,0.000000}%
\pgfsetstrokecolor{currentstroke}%
\pgfsetdash{}{0pt}%
\pgfsys@defobject{currentmarker}{\pgfqpoint{0.000000in}{0.000000in}}{\pgfqpoint{0.000000in}{0.027778in}}{%
\pgfpathmoveto{\pgfqpoint{0.000000in}{0.000000in}}%
\pgfpathlineto{\pgfqpoint{0.000000in}{0.027778in}}%
\pgfusepath{stroke,fill}%
}%
\begin{pgfscope}%
\pgfsys@transformshift{2.065434in}{0.456901in}%
\pgfsys@useobject{currentmarker}{}%
\end{pgfscope}%
\end{pgfscope}%
\begin{pgfscope}%
\pgfpathrectangle{\pgfqpoint{0.781757in}{0.456901in}}{\pgfqpoint{5.078868in}{2.816432in}}%
\pgfusepath{clip}%
\pgfsetbuttcap%
\pgfsetroundjoin%
\pgfsetlinewidth{0.501875pt}%
\definecolor{currentstroke}{rgb}{0.501961,0.501961,0.501961}%
\pgfsetstrokecolor{currentstroke}%
\pgfsetstrokeopacity{0.400000}%
\pgfsetdash{{1.850000pt}{0.800000pt}}{0.000000pt}%
\pgfpathmoveto{\pgfqpoint{2.322169in}{0.456901in}}%
\pgfpathlineto{\pgfqpoint{2.322169in}{3.273333in}}%
\pgfusepath{stroke}%
\end{pgfscope}%
\begin{pgfscope}%
\pgfsetbuttcap%
\pgfsetroundjoin%
\definecolor{currentfill}{rgb}{1.000000,1.000000,1.000000}%
\pgfsetfillcolor{currentfill}%
\pgfsetlinewidth{0.602250pt}%
\definecolor{currentstroke}{rgb}{0.000000,0.000000,0.000000}%
\pgfsetstrokecolor{currentstroke}%
\pgfsetdash{}{0pt}%
\pgfsys@defobject{currentmarker}{\pgfqpoint{0.000000in}{0.000000in}}{\pgfqpoint{0.000000in}{0.027778in}}{%
\pgfpathmoveto{\pgfqpoint{0.000000in}{0.000000in}}%
\pgfpathlineto{\pgfqpoint{0.000000in}{0.027778in}}%
\pgfusepath{stroke,fill}%
}%
\begin{pgfscope}%
\pgfsys@transformshift{2.322169in}{0.456901in}%
\pgfsys@useobject{currentmarker}{}%
\end{pgfscope}%
\end{pgfscope}%
\begin{pgfscope}%
\pgfpathrectangle{\pgfqpoint{0.781757in}{0.456901in}}{\pgfqpoint{5.078868in}{2.816432in}}%
\pgfusepath{clip}%
\pgfsetbuttcap%
\pgfsetroundjoin%
\pgfsetlinewidth{0.501875pt}%
\definecolor{currentstroke}{rgb}{0.501961,0.501961,0.501961}%
\pgfsetstrokecolor{currentstroke}%
\pgfsetstrokeopacity{0.400000}%
\pgfsetdash{{1.850000pt}{0.800000pt}}{0.000000pt}%
\pgfpathmoveto{\pgfqpoint{2.578905in}{0.456901in}}%
\pgfpathlineto{\pgfqpoint{2.578905in}{3.273333in}}%
\pgfusepath{stroke}%
\end{pgfscope}%
\begin{pgfscope}%
\pgfsetbuttcap%
\pgfsetroundjoin%
\definecolor{currentfill}{rgb}{1.000000,1.000000,1.000000}%
\pgfsetfillcolor{currentfill}%
\pgfsetlinewidth{0.602250pt}%
\definecolor{currentstroke}{rgb}{0.000000,0.000000,0.000000}%
\pgfsetstrokecolor{currentstroke}%
\pgfsetdash{}{0pt}%
\pgfsys@defobject{currentmarker}{\pgfqpoint{0.000000in}{0.000000in}}{\pgfqpoint{0.000000in}{0.027778in}}{%
\pgfpathmoveto{\pgfqpoint{0.000000in}{0.000000in}}%
\pgfpathlineto{\pgfqpoint{0.000000in}{0.027778in}}%
\pgfusepath{stroke,fill}%
}%
\begin{pgfscope}%
\pgfsys@transformshift{2.578905in}{0.456901in}%
\pgfsys@useobject{currentmarker}{}%
\end{pgfscope}%
\end{pgfscope}%
\begin{pgfscope}%
\pgfpathrectangle{\pgfqpoint{0.781757in}{0.456901in}}{\pgfqpoint{5.078868in}{2.816432in}}%
\pgfusepath{clip}%
\pgfsetbuttcap%
\pgfsetroundjoin%
\pgfsetlinewidth{0.501875pt}%
\definecolor{currentstroke}{rgb}{0.501961,0.501961,0.501961}%
\pgfsetstrokecolor{currentstroke}%
\pgfsetstrokeopacity{0.400000}%
\pgfsetdash{{1.850000pt}{0.800000pt}}{0.000000pt}%
\pgfpathmoveto{\pgfqpoint{3.092376in}{0.456901in}}%
\pgfpathlineto{\pgfqpoint{3.092376in}{3.273333in}}%
\pgfusepath{stroke}%
\end{pgfscope}%
\begin{pgfscope}%
\pgfsetbuttcap%
\pgfsetroundjoin%
\definecolor{currentfill}{rgb}{1.000000,1.000000,1.000000}%
\pgfsetfillcolor{currentfill}%
\pgfsetlinewidth{0.602250pt}%
\definecolor{currentstroke}{rgb}{0.000000,0.000000,0.000000}%
\pgfsetstrokecolor{currentstroke}%
\pgfsetdash{}{0pt}%
\pgfsys@defobject{currentmarker}{\pgfqpoint{0.000000in}{0.000000in}}{\pgfqpoint{0.000000in}{0.027778in}}{%
\pgfpathmoveto{\pgfqpoint{0.000000in}{0.000000in}}%
\pgfpathlineto{\pgfqpoint{0.000000in}{0.027778in}}%
\pgfusepath{stroke,fill}%
}%
\begin{pgfscope}%
\pgfsys@transformshift{3.092376in}{0.456901in}%
\pgfsys@useobject{currentmarker}{}%
\end{pgfscope}%
\end{pgfscope}%
\begin{pgfscope}%
\pgfpathrectangle{\pgfqpoint{0.781757in}{0.456901in}}{\pgfqpoint{5.078868in}{2.816432in}}%
\pgfusepath{clip}%
\pgfsetbuttcap%
\pgfsetroundjoin%
\pgfsetlinewidth{0.501875pt}%
\definecolor{currentstroke}{rgb}{0.501961,0.501961,0.501961}%
\pgfsetstrokecolor{currentstroke}%
\pgfsetstrokeopacity{0.400000}%
\pgfsetdash{{1.850000pt}{0.800000pt}}{0.000000pt}%
\pgfpathmoveto{\pgfqpoint{3.349111in}{0.456901in}}%
\pgfpathlineto{\pgfqpoint{3.349111in}{3.273333in}}%
\pgfusepath{stroke}%
\end{pgfscope}%
\begin{pgfscope}%
\pgfsetbuttcap%
\pgfsetroundjoin%
\definecolor{currentfill}{rgb}{1.000000,1.000000,1.000000}%
\pgfsetfillcolor{currentfill}%
\pgfsetlinewidth{0.602250pt}%
\definecolor{currentstroke}{rgb}{0.000000,0.000000,0.000000}%
\pgfsetstrokecolor{currentstroke}%
\pgfsetdash{}{0pt}%
\pgfsys@defobject{currentmarker}{\pgfqpoint{0.000000in}{0.000000in}}{\pgfqpoint{0.000000in}{0.027778in}}{%
\pgfpathmoveto{\pgfqpoint{0.000000in}{0.000000in}}%
\pgfpathlineto{\pgfqpoint{0.000000in}{0.027778in}}%
\pgfusepath{stroke,fill}%
}%
\begin{pgfscope}%
\pgfsys@transformshift{3.349111in}{0.456901in}%
\pgfsys@useobject{currentmarker}{}%
\end{pgfscope}%
\end{pgfscope}%
\begin{pgfscope}%
\pgfpathrectangle{\pgfqpoint{0.781757in}{0.456901in}}{\pgfqpoint{5.078868in}{2.816432in}}%
\pgfusepath{clip}%
\pgfsetbuttcap%
\pgfsetroundjoin%
\pgfsetlinewidth{0.501875pt}%
\definecolor{currentstroke}{rgb}{0.501961,0.501961,0.501961}%
\pgfsetstrokecolor{currentstroke}%
\pgfsetstrokeopacity{0.400000}%
\pgfsetdash{{1.850000pt}{0.800000pt}}{0.000000pt}%
\pgfpathmoveto{\pgfqpoint{3.605846in}{0.456901in}}%
\pgfpathlineto{\pgfqpoint{3.605846in}{3.273333in}}%
\pgfusepath{stroke}%
\end{pgfscope}%
\begin{pgfscope}%
\pgfsetbuttcap%
\pgfsetroundjoin%
\definecolor{currentfill}{rgb}{1.000000,1.000000,1.000000}%
\pgfsetfillcolor{currentfill}%
\pgfsetlinewidth{0.602250pt}%
\definecolor{currentstroke}{rgb}{0.000000,0.000000,0.000000}%
\pgfsetstrokecolor{currentstroke}%
\pgfsetdash{}{0pt}%
\pgfsys@defobject{currentmarker}{\pgfqpoint{0.000000in}{0.000000in}}{\pgfqpoint{0.000000in}{0.027778in}}{%
\pgfpathmoveto{\pgfqpoint{0.000000in}{0.000000in}}%
\pgfpathlineto{\pgfqpoint{0.000000in}{0.027778in}}%
\pgfusepath{stroke,fill}%
}%
\begin{pgfscope}%
\pgfsys@transformshift{3.605846in}{0.456901in}%
\pgfsys@useobject{currentmarker}{}%
\end{pgfscope}%
\end{pgfscope}%
\begin{pgfscope}%
\pgfpathrectangle{\pgfqpoint{0.781757in}{0.456901in}}{\pgfqpoint{5.078868in}{2.816432in}}%
\pgfusepath{clip}%
\pgfsetbuttcap%
\pgfsetroundjoin%
\pgfsetlinewidth{0.501875pt}%
\definecolor{currentstroke}{rgb}{0.501961,0.501961,0.501961}%
\pgfsetstrokecolor{currentstroke}%
\pgfsetstrokeopacity{0.400000}%
\pgfsetdash{{1.850000pt}{0.800000pt}}{0.000000pt}%
\pgfpathmoveto{\pgfqpoint{4.119317in}{0.456901in}}%
\pgfpathlineto{\pgfqpoint{4.119317in}{3.273333in}}%
\pgfusepath{stroke}%
\end{pgfscope}%
\begin{pgfscope}%
\pgfsetbuttcap%
\pgfsetroundjoin%
\definecolor{currentfill}{rgb}{1.000000,1.000000,1.000000}%
\pgfsetfillcolor{currentfill}%
\pgfsetlinewidth{0.602250pt}%
\definecolor{currentstroke}{rgb}{0.000000,0.000000,0.000000}%
\pgfsetstrokecolor{currentstroke}%
\pgfsetdash{}{0pt}%
\pgfsys@defobject{currentmarker}{\pgfqpoint{0.000000in}{0.000000in}}{\pgfqpoint{0.000000in}{0.027778in}}{%
\pgfpathmoveto{\pgfqpoint{0.000000in}{0.000000in}}%
\pgfpathlineto{\pgfqpoint{0.000000in}{0.027778in}}%
\pgfusepath{stroke,fill}%
}%
\begin{pgfscope}%
\pgfsys@transformshift{4.119317in}{0.456901in}%
\pgfsys@useobject{currentmarker}{}%
\end{pgfscope}%
\end{pgfscope}%
\begin{pgfscope}%
\pgfpathrectangle{\pgfqpoint{0.781757in}{0.456901in}}{\pgfqpoint{5.078868in}{2.816432in}}%
\pgfusepath{clip}%
\pgfsetbuttcap%
\pgfsetroundjoin%
\pgfsetlinewidth{0.501875pt}%
\definecolor{currentstroke}{rgb}{0.501961,0.501961,0.501961}%
\pgfsetstrokecolor{currentstroke}%
\pgfsetstrokeopacity{0.400000}%
\pgfsetdash{{1.850000pt}{0.800000pt}}{0.000000pt}%
\pgfpathmoveto{\pgfqpoint{4.376053in}{0.456901in}}%
\pgfpathlineto{\pgfqpoint{4.376053in}{3.273333in}}%
\pgfusepath{stroke}%
\end{pgfscope}%
\begin{pgfscope}%
\pgfsetbuttcap%
\pgfsetroundjoin%
\definecolor{currentfill}{rgb}{1.000000,1.000000,1.000000}%
\pgfsetfillcolor{currentfill}%
\pgfsetlinewidth{0.602250pt}%
\definecolor{currentstroke}{rgb}{0.000000,0.000000,0.000000}%
\pgfsetstrokecolor{currentstroke}%
\pgfsetdash{}{0pt}%
\pgfsys@defobject{currentmarker}{\pgfqpoint{0.000000in}{0.000000in}}{\pgfqpoint{0.000000in}{0.027778in}}{%
\pgfpathmoveto{\pgfqpoint{0.000000in}{0.000000in}}%
\pgfpathlineto{\pgfqpoint{0.000000in}{0.027778in}}%
\pgfusepath{stroke,fill}%
}%
\begin{pgfscope}%
\pgfsys@transformshift{4.376053in}{0.456901in}%
\pgfsys@useobject{currentmarker}{}%
\end{pgfscope}%
\end{pgfscope}%
\begin{pgfscope}%
\pgfpathrectangle{\pgfqpoint{0.781757in}{0.456901in}}{\pgfqpoint{5.078868in}{2.816432in}}%
\pgfusepath{clip}%
\pgfsetbuttcap%
\pgfsetroundjoin%
\pgfsetlinewidth{0.501875pt}%
\definecolor{currentstroke}{rgb}{0.501961,0.501961,0.501961}%
\pgfsetstrokecolor{currentstroke}%
\pgfsetstrokeopacity{0.400000}%
\pgfsetdash{{1.850000pt}{0.800000pt}}{0.000000pt}%
\pgfpathmoveto{\pgfqpoint{4.632788in}{0.456901in}}%
\pgfpathlineto{\pgfqpoint{4.632788in}{3.273333in}}%
\pgfusepath{stroke}%
\end{pgfscope}%
\begin{pgfscope}%
\pgfsetbuttcap%
\pgfsetroundjoin%
\definecolor{currentfill}{rgb}{1.000000,1.000000,1.000000}%
\pgfsetfillcolor{currentfill}%
\pgfsetlinewidth{0.602250pt}%
\definecolor{currentstroke}{rgb}{0.000000,0.000000,0.000000}%
\pgfsetstrokecolor{currentstroke}%
\pgfsetdash{}{0pt}%
\pgfsys@defobject{currentmarker}{\pgfqpoint{0.000000in}{0.000000in}}{\pgfqpoint{0.000000in}{0.027778in}}{%
\pgfpathmoveto{\pgfqpoint{0.000000in}{0.000000in}}%
\pgfpathlineto{\pgfqpoint{0.000000in}{0.027778in}}%
\pgfusepath{stroke,fill}%
}%
\begin{pgfscope}%
\pgfsys@transformshift{4.632788in}{0.456901in}%
\pgfsys@useobject{currentmarker}{}%
\end{pgfscope}%
\end{pgfscope}%
\begin{pgfscope}%
\pgfpathrectangle{\pgfqpoint{0.781757in}{0.456901in}}{\pgfqpoint{5.078868in}{2.816432in}}%
\pgfusepath{clip}%
\pgfsetbuttcap%
\pgfsetroundjoin%
\pgfsetlinewidth{0.501875pt}%
\definecolor{currentstroke}{rgb}{0.501961,0.501961,0.501961}%
\pgfsetstrokecolor{currentstroke}%
\pgfsetstrokeopacity{0.400000}%
\pgfsetdash{{1.850000pt}{0.800000pt}}{0.000000pt}%
\pgfpathmoveto{\pgfqpoint{5.146259in}{0.456901in}}%
\pgfpathlineto{\pgfqpoint{5.146259in}{3.273333in}}%
\pgfusepath{stroke}%
\end{pgfscope}%
\begin{pgfscope}%
\pgfsetbuttcap%
\pgfsetroundjoin%
\definecolor{currentfill}{rgb}{1.000000,1.000000,1.000000}%
\pgfsetfillcolor{currentfill}%
\pgfsetlinewidth{0.602250pt}%
\definecolor{currentstroke}{rgb}{0.000000,0.000000,0.000000}%
\pgfsetstrokecolor{currentstroke}%
\pgfsetdash{}{0pt}%
\pgfsys@defobject{currentmarker}{\pgfqpoint{0.000000in}{0.000000in}}{\pgfqpoint{0.000000in}{0.027778in}}{%
\pgfpathmoveto{\pgfqpoint{0.000000in}{0.000000in}}%
\pgfpathlineto{\pgfqpoint{0.000000in}{0.027778in}}%
\pgfusepath{stroke,fill}%
}%
\begin{pgfscope}%
\pgfsys@transformshift{5.146259in}{0.456901in}%
\pgfsys@useobject{currentmarker}{}%
\end{pgfscope}%
\end{pgfscope}%
\begin{pgfscope}%
\pgfpathrectangle{\pgfqpoint{0.781757in}{0.456901in}}{\pgfqpoint{5.078868in}{2.816432in}}%
\pgfusepath{clip}%
\pgfsetbuttcap%
\pgfsetroundjoin%
\pgfsetlinewidth{0.501875pt}%
\definecolor{currentstroke}{rgb}{0.501961,0.501961,0.501961}%
\pgfsetstrokecolor{currentstroke}%
\pgfsetstrokeopacity{0.400000}%
\pgfsetdash{{1.850000pt}{0.800000pt}}{0.000000pt}%
\pgfpathmoveto{\pgfqpoint{5.402994in}{0.456901in}}%
\pgfpathlineto{\pgfqpoint{5.402994in}{3.273333in}}%
\pgfusepath{stroke}%
\end{pgfscope}%
\begin{pgfscope}%
\pgfsetbuttcap%
\pgfsetroundjoin%
\definecolor{currentfill}{rgb}{1.000000,1.000000,1.000000}%
\pgfsetfillcolor{currentfill}%
\pgfsetlinewidth{0.602250pt}%
\definecolor{currentstroke}{rgb}{0.000000,0.000000,0.000000}%
\pgfsetstrokecolor{currentstroke}%
\pgfsetdash{}{0pt}%
\pgfsys@defobject{currentmarker}{\pgfqpoint{0.000000in}{0.000000in}}{\pgfqpoint{0.000000in}{0.027778in}}{%
\pgfpathmoveto{\pgfqpoint{0.000000in}{0.000000in}}%
\pgfpathlineto{\pgfqpoint{0.000000in}{0.027778in}}%
\pgfusepath{stroke,fill}%
}%
\begin{pgfscope}%
\pgfsys@transformshift{5.402994in}{0.456901in}%
\pgfsys@useobject{currentmarker}{}%
\end{pgfscope}%
\end{pgfscope}%
\begin{pgfscope}%
\pgfpathrectangle{\pgfqpoint{0.781757in}{0.456901in}}{\pgfqpoint{5.078868in}{2.816432in}}%
\pgfusepath{clip}%
\pgfsetbuttcap%
\pgfsetroundjoin%
\pgfsetlinewidth{0.501875pt}%
\definecolor{currentstroke}{rgb}{0.501961,0.501961,0.501961}%
\pgfsetstrokecolor{currentstroke}%
\pgfsetstrokeopacity{0.400000}%
\pgfsetdash{{1.850000pt}{0.800000pt}}{0.000000pt}%
\pgfpathmoveto{\pgfqpoint{5.659730in}{0.456901in}}%
\pgfpathlineto{\pgfqpoint{5.659730in}{3.273333in}}%
\pgfusepath{stroke}%
\end{pgfscope}%
\begin{pgfscope}%
\pgfsetbuttcap%
\pgfsetroundjoin%
\definecolor{currentfill}{rgb}{1.000000,1.000000,1.000000}%
\pgfsetfillcolor{currentfill}%
\pgfsetlinewidth{0.602250pt}%
\definecolor{currentstroke}{rgb}{0.000000,0.000000,0.000000}%
\pgfsetstrokecolor{currentstroke}%
\pgfsetdash{}{0pt}%
\pgfsys@defobject{currentmarker}{\pgfqpoint{0.000000in}{0.000000in}}{\pgfqpoint{0.000000in}{0.027778in}}{%
\pgfpathmoveto{\pgfqpoint{0.000000in}{0.000000in}}%
\pgfpathlineto{\pgfqpoint{0.000000in}{0.027778in}}%
\pgfusepath{stroke,fill}%
}%
\begin{pgfscope}%
\pgfsys@transformshift{5.659730in}{0.456901in}%
\pgfsys@useobject{currentmarker}{}%
\end{pgfscope}%
\end{pgfscope}%
\begin{pgfscope}%
\definecolor{textcolor}{rgb}{0.000000,0.000000,0.000000}%
\pgfsetstrokecolor{textcolor}%
\pgfsetfillcolor{textcolor}%
\pgftext[x=3.321191in,y=0.201367in,,top]{\color{textcolor}{\rmfamily\fontsize{12.000000}{14.400000}\selectfont\catcode`\^=\active\def^{\ifmmode\sp\else\^{}\fi}\catcode`\%=\active\def%{\%}$t$}}%
\end{pgfscope}%
\begin{pgfscope}%
\pgfpathrectangle{\pgfqpoint{0.781757in}{0.456901in}}{\pgfqpoint{5.078868in}{2.816432in}}%
\pgfusepath{clip}%
\pgfsetbuttcap%
\pgfsetroundjoin%
\pgfsetlinewidth{0.501875pt}%
\definecolor{currentstroke}{rgb}{0.501961,0.501961,0.501961}%
\pgfsetstrokecolor{currentstroke}%
\pgfsetstrokeopacity{0.400000}%
\pgfsetdash{{1.850000pt}{0.800000pt}}{0.000000pt}%
\pgfpathmoveto{\pgfqpoint{0.781757in}{0.883462in}}%
\pgfpathlineto{\pgfqpoint{5.860625in}{0.883462in}}%
\pgfusepath{stroke}%
\end{pgfscope}%
\begin{pgfscope}%
\pgfsetbuttcap%
\pgfsetroundjoin%
\definecolor{currentfill}{rgb}{1.000000,1.000000,1.000000}%
\pgfsetfillcolor{currentfill}%
\pgfsetlinewidth{0.803000pt}%
\definecolor{currentstroke}{rgb}{0.000000,0.000000,0.000000}%
\pgfsetstrokecolor{currentstroke}%
\pgfsetdash{}{0pt}%
\pgfsys@defobject{currentmarker}{\pgfqpoint{0.000000in}{0.000000in}}{\pgfqpoint{0.048611in}{0.000000in}}{%
\pgfpathmoveto{\pgfqpoint{0.000000in}{0.000000in}}%
\pgfpathlineto{\pgfqpoint{0.048611in}{0.000000in}}%
\pgfusepath{stroke,fill}%
}%
\begin{pgfscope}%
\pgfsys@transformshift{0.781757in}{0.883462in}%
\pgfsys@useobject{currentmarker}{}%
\end{pgfscope}%
\end{pgfscope}%
\begin{pgfscope}%
\definecolor{textcolor}{rgb}{0.000000,0.000000,0.000000}%
\pgfsetstrokecolor{textcolor}%
\pgfsetfillcolor{textcolor}%
\pgftext[x=0.353325in, y=0.825601in, left, base]{\color{textcolor}{\rmfamily\fontsize{12.000000}{14.400000}\selectfont\catcode`\^=\active\def^{\ifmmode\sp\else\^{}\fi}\catcode`\%=\active\def%{\%}$\mathdefault{\ensuremath{-}1.0}$}}%
\end{pgfscope}%
\begin{pgfscope}%
\pgfpathrectangle{\pgfqpoint{0.781757in}{0.456901in}}{\pgfqpoint{5.078868in}{2.816432in}}%
\pgfusepath{clip}%
\pgfsetbuttcap%
\pgfsetroundjoin%
\pgfsetlinewidth{0.501875pt}%
\definecolor{currentstroke}{rgb}{0.501961,0.501961,0.501961}%
\pgfsetstrokecolor{currentstroke}%
\pgfsetstrokeopacity{0.400000}%
\pgfsetdash{{1.850000pt}{0.800000pt}}{0.000000pt}%
\pgfpathmoveto{\pgfqpoint{0.781757in}{1.410952in}}%
\pgfpathlineto{\pgfqpoint{5.860625in}{1.410952in}}%
\pgfusepath{stroke}%
\end{pgfscope}%
\begin{pgfscope}%
\pgfsetbuttcap%
\pgfsetroundjoin%
\definecolor{currentfill}{rgb}{1.000000,1.000000,1.000000}%
\pgfsetfillcolor{currentfill}%
\pgfsetlinewidth{0.803000pt}%
\definecolor{currentstroke}{rgb}{0.000000,0.000000,0.000000}%
\pgfsetstrokecolor{currentstroke}%
\pgfsetdash{}{0pt}%
\pgfsys@defobject{currentmarker}{\pgfqpoint{0.000000in}{0.000000in}}{\pgfqpoint{0.048611in}{0.000000in}}{%
\pgfpathmoveto{\pgfqpoint{0.000000in}{0.000000in}}%
\pgfpathlineto{\pgfqpoint{0.048611in}{0.000000in}}%
\pgfusepath{stroke,fill}%
}%
\begin{pgfscope}%
\pgfsys@transformshift{0.781757in}{1.410952in}%
\pgfsys@useobject{currentmarker}{}%
\end{pgfscope}%
\end{pgfscope}%
\begin{pgfscope}%
\definecolor{textcolor}{rgb}{0.000000,0.000000,0.000000}%
\pgfsetstrokecolor{textcolor}%
\pgfsetfillcolor{textcolor}%
\pgftext[x=0.353325in, y=1.353091in, left, base]{\color{textcolor}{\rmfamily\fontsize{12.000000}{14.400000}\selectfont\catcode`\^=\active\def^{\ifmmode\sp\else\^{}\fi}\catcode`\%=\active\def%{\%}$\mathdefault{\ensuremath{-}0.5}$}}%
\end{pgfscope}%
\begin{pgfscope}%
\pgfpathrectangle{\pgfqpoint{0.781757in}{0.456901in}}{\pgfqpoint{5.078868in}{2.816432in}}%
\pgfusepath{clip}%
\pgfsetbuttcap%
\pgfsetroundjoin%
\pgfsetlinewidth{0.501875pt}%
\definecolor{currentstroke}{rgb}{0.501961,0.501961,0.501961}%
\pgfsetstrokecolor{currentstroke}%
\pgfsetstrokeopacity{0.400000}%
\pgfsetdash{{1.850000pt}{0.800000pt}}{0.000000pt}%
\pgfpathmoveto{\pgfqpoint{0.781757in}{1.938442in}}%
\pgfpathlineto{\pgfqpoint{5.860625in}{1.938442in}}%
\pgfusepath{stroke}%
\end{pgfscope}%
\begin{pgfscope}%
\pgfsetbuttcap%
\pgfsetroundjoin%
\definecolor{currentfill}{rgb}{1.000000,1.000000,1.000000}%
\pgfsetfillcolor{currentfill}%
\pgfsetlinewidth{0.803000pt}%
\definecolor{currentstroke}{rgb}{0.000000,0.000000,0.000000}%
\pgfsetstrokecolor{currentstroke}%
\pgfsetdash{}{0pt}%
\pgfsys@defobject{currentmarker}{\pgfqpoint{0.000000in}{0.000000in}}{\pgfqpoint{0.048611in}{0.000000in}}{%
\pgfpathmoveto{\pgfqpoint{0.000000in}{0.000000in}}%
\pgfpathlineto{\pgfqpoint{0.048611in}{0.000000in}}%
\pgfusepath{stroke,fill}%
}%
\begin{pgfscope}%
\pgfsys@transformshift{0.781757in}{1.938442in}%
\pgfsys@useobject{currentmarker}{}%
\end{pgfscope}%
\end{pgfscope}%
\begin{pgfscope}%
\definecolor{textcolor}{rgb}{0.000000,0.000000,0.000000}%
\pgfsetstrokecolor{textcolor}%
\pgfsetfillcolor{textcolor}%
\pgftext[x=0.482955in, y=1.880581in, left, base]{\color{textcolor}{\rmfamily\fontsize{12.000000}{14.400000}\selectfont\catcode`\^=\active\def^{\ifmmode\sp\else\^{}\fi}\catcode`\%=\active\def%{\%}$\mathdefault{0.0}$}}%
\end{pgfscope}%
\begin{pgfscope}%
\pgfpathrectangle{\pgfqpoint{0.781757in}{0.456901in}}{\pgfqpoint{5.078868in}{2.816432in}}%
\pgfusepath{clip}%
\pgfsetbuttcap%
\pgfsetroundjoin%
\pgfsetlinewidth{0.501875pt}%
\definecolor{currentstroke}{rgb}{0.501961,0.501961,0.501961}%
\pgfsetstrokecolor{currentstroke}%
\pgfsetstrokeopacity{0.400000}%
\pgfsetdash{{1.850000pt}{0.800000pt}}{0.000000pt}%
\pgfpathmoveto{\pgfqpoint{0.781757in}{2.465933in}}%
\pgfpathlineto{\pgfqpoint{5.860625in}{2.465933in}}%
\pgfusepath{stroke}%
\end{pgfscope}%
\begin{pgfscope}%
\pgfsetbuttcap%
\pgfsetroundjoin%
\definecolor{currentfill}{rgb}{1.000000,1.000000,1.000000}%
\pgfsetfillcolor{currentfill}%
\pgfsetlinewidth{0.803000pt}%
\definecolor{currentstroke}{rgb}{0.000000,0.000000,0.000000}%
\pgfsetstrokecolor{currentstroke}%
\pgfsetdash{}{0pt}%
\pgfsys@defobject{currentmarker}{\pgfqpoint{0.000000in}{0.000000in}}{\pgfqpoint{0.048611in}{0.000000in}}{%
\pgfpathmoveto{\pgfqpoint{0.000000in}{0.000000in}}%
\pgfpathlineto{\pgfqpoint{0.048611in}{0.000000in}}%
\pgfusepath{stroke,fill}%
}%
\begin{pgfscope}%
\pgfsys@transformshift{0.781757in}{2.465933in}%
\pgfsys@useobject{currentmarker}{}%
\end{pgfscope}%
\end{pgfscope}%
\begin{pgfscope}%
\definecolor{textcolor}{rgb}{0.000000,0.000000,0.000000}%
\pgfsetstrokecolor{textcolor}%
\pgfsetfillcolor{textcolor}%
\pgftext[x=0.482955in, y=2.408071in, left, base]{\color{textcolor}{\rmfamily\fontsize{12.000000}{14.400000}\selectfont\catcode`\^=\active\def^{\ifmmode\sp\else\^{}\fi}\catcode`\%=\active\def%{\%}$\mathdefault{0.5}$}}%
\end{pgfscope}%
\begin{pgfscope}%
\pgfpathrectangle{\pgfqpoint{0.781757in}{0.456901in}}{\pgfqpoint{5.078868in}{2.816432in}}%
\pgfusepath{clip}%
\pgfsetbuttcap%
\pgfsetroundjoin%
\pgfsetlinewidth{0.501875pt}%
\definecolor{currentstroke}{rgb}{0.501961,0.501961,0.501961}%
\pgfsetstrokecolor{currentstroke}%
\pgfsetstrokeopacity{0.400000}%
\pgfsetdash{{1.850000pt}{0.800000pt}}{0.000000pt}%
\pgfpathmoveto{\pgfqpoint{0.781757in}{2.993423in}}%
\pgfpathlineto{\pgfqpoint{5.860625in}{2.993423in}}%
\pgfusepath{stroke}%
\end{pgfscope}%
\begin{pgfscope}%
\pgfsetbuttcap%
\pgfsetroundjoin%
\definecolor{currentfill}{rgb}{1.000000,1.000000,1.000000}%
\pgfsetfillcolor{currentfill}%
\pgfsetlinewidth{0.803000pt}%
\definecolor{currentstroke}{rgb}{0.000000,0.000000,0.000000}%
\pgfsetstrokecolor{currentstroke}%
\pgfsetdash{}{0pt}%
\pgfsys@defobject{currentmarker}{\pgfqpoint{0.000000in}{0.000000in}}{\pgfqpoint{0.048611in}{0.000000in}}{%
\pgfpathmoveto{\pgfqpoint{0.000000in}{0.000000in}}%
\pgfpathlineto{\pgfqpoint{0.048611in}{0.000000in}}%
\pgfusepath{stroke,fill}%
}%
\begin{pgfscope}%
\pgfsys@transformshift{0.781757in}{2.993423in}%
\pgfsys@useobject{currentmarker}{}%
\end{pgfscope}%
\end{pgfscope}%
\begin{pgfscope}%
\definecolor{textcolor}{rgb}{0.000000,0.000000,0.000000}%
\pgfsetstrokecolor{textcolor}%
\pgfsetfillcolor{textcolor}%
\pgftext[x=0.482955in, y=2.935562in, left, base]{\color{textcolor}{\rmfamily\fontsize{12.000000}{14.400000}\selectfont\catcode`\^=\active\def^{\ifmmode\sp\else\^{}\fi}\catcode`\%=\active\def%{\%}$\mathdefault{1.0}$}}%
\end{pgfscope}%
\begin{pgfscope}%
\definecolor{textcolor}{rgb}{0.000000,0.000000,0.000000}%
\pgfsetstrokecolor{textcolor}%
\pgfsetfillcolor{textcolor}%
\pgftext[x=0.297770in,y=1.865117in,,bottom,rotate=90.000000]{\color{textcolor}{\rmfamily\fontsize{12.000000}{14.400000}\selectfont\catcode`\^=\active\def^{\ifmmode\sp\else\^{}\fi}\catcode`\%=\active\def%{\%}$i_2(t)$}}%
\end{pgfscope}%
\begin{pgfscope}%
\pgfpathrectangle{\pgfqpoint{0.781757in}{0.456901in}}{\pgfqpoint{5.078868in}{2.816432in}}%
\pgfusepath{clip}%
\pgfsetrectcap%
\pgfsetroundjoin%
\pgfsetlinewidth{2.007500pt}%
\definecolor{currentstroke}{rgb}{0.000000,0.000000,0.000000}%
\pgfsetstrokecolor{currentstroke}%
\pgfsetdash{}{0pt}%
\pgfpathmoveto{\pgfqpoint{0.779190in}{2.092290in}}%
\pgfpathlineto{\pgfqpoint{0.798586in}{2.212448in}}%
\pgfpathlineto{\pgfqpoint{0.876173in}{2.735747in}}%
\pgfpathlineto{\pgfqpoint{0.895569in}{2.848837in}}%
\pgfpathlineto{\pgfqpoint{0.914966in}{2.951610in}}%
\pgfpathlineto{\pgfqpoint{0.934363in}{3.039641in}}%
\pgfpathlineto{\pgfqpoint{0.944061in}{3.075804in}}%
\pgfpathlineto{\pgfqpoint{0.953759in}{3.105293in}}%
\pgfpathlineto{\pgfqpoint{0.963458in}{3.127128in}}%
\pgfpathlineto{\pgfqpoint{0.973156in}{3.140581in}}%
\pgfpathlineto{\pgfqpoint{0.982854in}{3.145314in}}%
\pgfpathlineto{\pgfqpoint{0.992553in}{3.141473in}}%
\pgfpathlineto{\pgfqpoint{1.002251in}{3.129740in}}%
\pgfpathlineto{\pgfqpoint{1.011949in}{3.111319in}}%
\pgfpathlineto{\pgfqpoint{1.021648in}{3.087857in}}%
\pgfpathlineto{\pgfqpoint{1.060441in}{2.984020in}}%
\pgfpathlineto{\pgfqpoint{1.070139in}{2.964977in}}%
\pgfpathlineto{\pgfqpoint{1.079838in}{2.951243in}}%
\pgfpathlineto{\pgfqpoint{1.089536in}{2.943165in}}%
\pgfpathlineto{\pgfqpoint{1.099234in}{2.940518in}}%
\pgfpathlineto{\pgfqpoint{1.108933in}{2.942579in}}%
\pgfpathlineto{\pgfqpoint{1.118631in}{2.948247in}}%
\pgfpathlineto{\pgfqpoint{1.138027in}{2.965107in}}%
\pgfpathlineto{\pgfqpoint{1.157424in}{2.981085in}}%
\pgfpathlineto{\pgfqpoint{1.167122in}{2.986602in}}%
\pgfpathlineto{\pgfqpoint{1.176821in}{2.990066in}}%
\pgfpathlineto{\pgfqpoint{1.186519in}{2.991643in}}%
\pgfpathlineto{\pgfqpoint{1.205916in}{2.991197in}}%
\pgfpathlineto{\pgfqpoint{1.225312in}{2.990465in}}%
\pgfpathlineto{\pgfqpoint{1.235011in}{2.991422in}}%
\pgfpathlineto{\pgfqpoint{1.254407in}{2.996884in}}%
\pgfpathlineto{\pgfqpoint{1.283502in}{3.008821in}}%
\pgfpathlineto{\pgfqpoint{1.293201in}{3.011368in}}%
\pgfpathlineto{\pgfqpoint{1.302899in}{3.012235in}}%
\pgfpathlineto{\pgfqpoint{1.312597in}{3.011130in}}%
\pgfpathlineto{\pgfqpoint{1.322296in}{3.008041in}}%
\pgfpathlineto{\pgfqpoint{1.341692in}{2.997346in}}%
\pgfpathlineto{\pgfqpoint{1.361089in}{2.985176in}}%
\pgfpathlineto{\pgfqpoint{1.370787in}{2.980527in}}%
\pgfpathlineto{\pgfqpoint{1.380486in}{2.977718in}}%
\pgfpathlineto{\pgfqpoint{1.390184in}{2.977109in}}%
\pgfpathlineto{\pgfqpoint{1.399882in}{2.978747in}}%
\pgfpathlineto{\pgfqpoint{1.409580in}{2.982349in}}%
\pgfpathlineto{\pgfqpoint{1.448374in}{3.002664in}}%
\pgfpathlineto{\pgfqpoint{1.458072in}{3.005269in}}%
\pgfpathlineto{\pgfqpoint{1.467770in}{3.005776in}}%
\pgfpathlineto{\pgfqpoint{1.477469in}{3.004129in}}%
\pgfpathlineto{\pgfqpoint{1.487167in}{3.000615in}}%
\pgfpathlineto{\pgfqpoint{1.525960in}{2.981932in}}%
\pgfpathlineto{\pgfqpoint{1.535659in}{2.980109in}}%
\pgfpathlineto{\pgfqpoint{1.545357in}{2.980506in}}%
\pgfpathlineto{\pgfqpoint{1.555055in}{2.983128in}}%
\pgfpathlineto{\pgfqpoint{1.574452in}{2.993301in}}%
\pgfpathlineto{\pgfqpoint{1.593849in}{3.004631in}}%
\pgfpathlineto{\pgfqpoint{1.603547in}{3.008411in}}%
\pgfpathlineto{\pgfqpoint{1.613245in}{3.009955in}}%
\pgfpathlineto{\pgfqpoint{1.622944in}{3.008913in}}%
\pgfpathlineto{\pgfqpoint{1.632642in}{3.005336in}}%
\pgfpathlineto{\pgfqpoint{1.652039in}{2.992761in}}%
\pgfpathlineto{\pgfqpoint{1.671435in}{2.979399in}}%
\pgfpathlineto{\pgfqpoint{1.681133in}{2.975135in}}%
\pgfpathlineto{\pgfqpoint{1.690832in}{2.973612in}}%
\pgfpathlineto{\pgfqpoint{1.700530in}{2.975221in}}%
\pgfpathlineto{\pgfqpoint{1.710228in}{2.979874in}}%
\pgfpathlineto{\pgfqpoint{1.719927in}{2.986996in}}%
\pgfpathlineto{\pgfqpoint{1.749022in}{3.011870in}}%
\pgfpathlineto{\pgfqpoint{1.758720in}{3.016859in}}%
\pgfpathlineto{\pgfqpoint{1.768418in}{3.018317in}}%
\pgfpathlineto{\pgfqpoint{1.778117in}{3.015736in}}%
\pgfpathlineto{\pgfqpoint{1.787815in}{3.009211in}}%
\pgfpathlineto{\pgfqpoint{1.797513in}{2.999465in}}%
\pgfpathlineto{\pgfqpoint{1.826608in}{2.965782in}}%
\pgfpathlineto{\pgfqpoint{1.836307in}{2.959217in}}%
\pgfpathlineto{\pgfqpoint{1.846005in}{2.957737in}}%
\pgfpathlineto{\pgfqpoint{1.855703in}{2.962240in}}%
\pgfpathlineto{\pgfqpoint{1.865402in}{2.972812in}}%
\pgfpathlineto{\pgfqpoint{1.875100in}{2.988613in}}%
\pgfpathlineto{\pgfqpoint{1.904195in}{3.045762in}}%
\pgfpathlineto{\pgfqpoint{1.913893in}{3.057613in}}%
\pgfpathlineto{\pgfqpoint{1.923591in}{3.059991in}}%
\pgfpathlineto{\pgfqpoint{1.933290in}{3.049671in}}%
\pgfpathlineto{\pgfqpoint{1.942988in}{3.024075in}}%
\pgfpathlineto{\pgfqpoint{1.952686in}{2.981508in}}%
\pgfpathlineto{\pgfqpoint{1.962385in}{2.921301in}}%
\pgfpathlineto{\pgfqpoint{1.972083in}{2.843845in}}%
\pgfpathlineto{\pgfqpoint{1.981781in}{2.750506in}}%
\pgfpathlineto{\pgfqpoint{2.001178in}{2.525394in}}%
\pgfpathlineto{\pgfqpoint{2.030273in}{2.134478in}}%
\pgfpathlineto{\pgfqpoint{2.069066in}{1.610726in}}%
\pgfpathlineto{\pgfqpoint{2.098161in}{1.249860in}}%
\pgfpathlineto{\pgfqpoint{2.117558in}{1.030783in}}%
\pgfpathlineto{\pgfqpoint{2.136955in}{0.839595in}}%
\pgfpathlineto{\pgfqpoint{2.146653in}{0.759271in}}%
\pgfpathlineto{\pgfqpoint{2.156351in}{0.692031in}}%
\pgfpathlineto{\pgfqpoint{2.166050in}{0.639853in}}%
\pgfpathlineto{\pgfqpoint{2.175748in}{0.604238in}}%
\pgfpathlineto{\pgfqpoint{2.185446in}{0.585966in}}%
\pgfpathlineto{\pgfqpoint{2.195144in}{0.584920in}}%
\pgfpathlineto{\pgfqpoint{2.204843in}{0.600002in}}%
\pgfpathlineto{\pgfqpoint{2.214541in}{0.629144in}}%
\pgfpathlineto{\pgfqpoint{2.224239in}{0.669438in}}%
\pgfpathlineto{\pgfqpoint{2.243636in}{0.769030in}}%
\pgfpathlineto{\pgfqpoint{2.263033in}{0.868540in}}%
\pgfpathlineto{\pgfqpoint{2.272731in}{0.909914in}}%
\pgfpathlineto{\pgfqpoint{2.282429in}{0.942643in}}%
\pgfpathlineto{\pgfqpoint{2.292128in}{0.965597in}}%
\pgfpathlineto{\pgfqpoint{2.301826in}{0.978629in}}%
\pgfpathlineto{\pgfqpoint{2.311524in}{0.982480in}}%
\pgfpathlineto{\pgfqpoint{2.321223in}{0.978612in}}%
\pgfpathlineto{\pgfqpoint{2.330921in}{0.968964in}}%
\pgfpathlineto{\pgfqpoint{2.350318in}{0.940876in}}%
\pgfpathlineto{\pgfqpoint{2.369714in}{0.913432in}}%
\pgfpathlineto{\pgfqpoint{2.379413in}{0.902960in}}%
\pgfpathlineto{\pgfqpoint{2.389111in}{0.895177in}}%
\pgfpathlineto{\pgfqpoint{2.398809in}{0.889852in}}%
\pgfpathlineto{\pgfqpoint{2.418206in}{0.884034in}}%
\pgfpathlineto{\pgfqpoint{2.437603in}{0.879545in}}%
\pgfpathlineto{\pgfqpoint{2.456999in}{0.872369in}}%
\pgfpathlineto{\pgfqpoint{2.486094in}{0.859063in}}%
\pgfpathlineto{\pgfqpoint{2.495792in}{0.856290in}}%
\pgfpathlineto{\pgfqpoint{2.505491in}{0.855402in}}%
\pgfpathlineto{\pgfqpoint{2.515189in}{0.856775in}}%
\pgfpathlineto{\pgfqpoint{2.524887in}{0.860484in}}%
\pgfpathlineto{\pgfqpoint{2.534586in}{0.866276in}}%
\pgfpathlineto{\pgfqpoint{2.573379in}{0.896235in}}%
\pgfpathlineto{\pgfqpoint{2.583077in}{0.901118in}}%
\pgfpathlineto{\pgfqpoint{2.592776in}{0.903612in}}%
\pgfpathlineto{\pgfqpoint{2.602474in}{0.903523in}}%
\pgfpathlineto{\pgfqpoint{2.612172in}{0.901006in}}%
\pgfpathlineto{\pgfqpoint{2.631569in}{0.890870in}}%
\pgfpathlineto{\pgfqpoint{2.650966in}{0.879394in}}%
\pgfpathlineto{\pgfqpoint{2.660664in}{0.875245in}}%
\pgfpathlineto{\pgfqpoint{2.670362in}{0.872948in}}%
\pgfpathlineto{\pgfqpoint{2.680061in}{0.872739in}}%
\pgfpathlineto{\pgfqpoint{2.689759in}{0.874523in}}%
\pgfpathlineto{\pgfqpoint{2.709156in}{0.882211in}}%
\pgfpathlineto{\pgfqpoint{2.728552in}{0.890562in}}%
\pgfpathlineto{\pgfqpoint{2.738250in}{0.893136in}}%
\pgfpathlineto{\pgfqpoint{2.747949in}{0.893965in}}%
\pgfpathlineto{\pgfqpoint{2.757647in}{0.892886in}}%
\pgfpathlineto{\pgfqpoint{2.767345in}{0.890050in}}%
\pgfpathlineto{\pgfqpoint{2.815837in}{0.870435in}}%
\pgfpathlineto{\pgfqpoint{2.825535in}{0.870162in}}%
\pgfpathlineto{\pgfqpoint{2.835234in}{0.871949in}}%
\pgfpathlineto{\pgfqpoint{2.844932in}{0.875576in}}%
\pgfpathlineto{\pgfqpoint{2.883725in}{0.895546in}}%
\pgfpathlineto{\pgfqpoint{2.893424in}{0.897961in}}%
\pgfpathlineto{\pgfqpoint{2.903122in}{0.898171in}}%
\pgfpathlineto{\pgfqpoint{2.912820in}{0.896082in}}%
\pgfpathlineto{\pgfqpoint{2.922519in}{0.891950in}}%
\pgfpathlineto{\pgfqpoint{2.961312in}{0.869649in}}%
\pgfpathlineto{\pgfqpoint{2.971010in}{0.867104in}}%
\pgfpathlineto{\pgfqpoint{2.980709in}{0.867122in}}%
\pgfpathlineto{\pgfqpoint{2.990407in}{0.869830in}}%
\pgfpathlineto{\pgfqpoint{3.000105in}{0.874940in}}%
\pgfpathlineto{\pgfqpoint{3.038898in}{0.902071in}}%
\pgfpathlineto{\pgfqpoint{3.048597in}{0.904980in}}%
\pgfpathlineto{\pgfqpoint{3.058295in}{0.904529in}}%
\pgfpathlineto{\pgfqpoint{3.067993in}{0.900442in}}%
\pgfpathlineto{\pgfqpoint{3.077692in}{0.892964in}}%
\pgfpathlineto{\pgfqpoint{3.097088in}{0.871453in}}%
\pgfpathlineto{\pgfqpoint{3.106787in}{0.860342in}}%
\pgfpathlineto{\pgfqpoint{3.116485in}{0.851420in}}%
\pgfpathlineto{\pgfqpoint{3.126183in}{0.846605in}}%
\pgfpathlineto{\pgfqpoint{3.135882in}{0.847667in}}%
\pgfpathlineto{\pgfqpoint{3.145580in}{0.856064in}}%
\pgfpathlineto{\pgfqpoint{3.155278in}{0.872800in}}%
\pgfpathlineto{\pgfqpoint{3.164977in}{0.898354in}}%
\pgfpathlineto{\pgfqpoint{3.174675in}{0.932652in}}%
\pgfpathlineto{\pgfqpoint{3.184373in}{0.975113in}}%
\pgfpathlineto{\pgfqpoint{3.194072in}{1.024739in}}%
\pgfpathlineto{\pgfqpoint{3.213468in}{1.140209in}}%
\pgfpathlineto{\pgfqpoint{3.252262in}{1.399000in}}%
\pgfpathlineto{\pgfqpoint{3.291055in}{1.653073in}}%
\pgfpathlineto{\pgfqpoint{3.320150in}{1.829445in}}%
\pgfpathlineto{\pgfqpoint{3.339546in}{1.932742in}}%
\pgfpathlineto{\pgfqpoint{3.349245in}{1.977165in}}%
\pgfpathlineto{\pgfqpoint{3.358943in}{2.015238in}}%
\pgfpathlineto{\pgfqpoint{3.368641in}{2.045908in}}%
\pgfpathlineto{\pgfqpoint{3.378340in}{2.068357in}}%
\pgfpathlineto{\pgfqpoint{3.388038in}{2.082111in}}%
\pgfpathlineto{\pgfqpoint{3.397736in}{2.087130in}}%
\pgfpathlineto{\pgfqpoint{3.407435in}{2.083841in}}%
\pgfpathlineto{\pgfqpoint{3.417133in}{2.073139in}}%
\pgfpathlineto{\pgfqpoint{3.426831in}{2.056328in}}%
\pgfpathlineto{\pgfqpoint{3.436530in}{2.035020in}}%
\pgfpathlineto{\pgfqpoint{3.475323in}{1.940349in}}%
\pgfpathlineto{\pgfqpoint{3.485021in}{1.922049in}}%
\pgfpathlineto{\pgfqpoint{3.494720in}{1.907903in}}%
\pgfpathlineto{\pgfqpoint{3.504418in}{1.898203in}}%
\pgfpathlineto{\pgfqpoint{3.514116in}{1.892846in}}%
\pgfpathlineto{\pgfqpoint{3.523814in}{1.891391in}}%
\pgfpathlineto{\pgfqpoint{3.533513in}{1.893152in}}%
\pgfpathlineto{\pgfqpoint{3.543211in}{1.897302in}}%
\pgfpathlineto{\pgfqpoint{3.562608in}{1.909381in}}%
\pgfpathlineto{\pgfqpoint{3.582004in}{1.921845in}}%
\pgfpathlineto{\pgfqpoint{3.601401in}{1.931531in}}%
\pgfpathlineto{\pgfqpoint{3.620798in}{1.938020in}}%
\pgfpathlineto{\pgfqpoint{3.649893in}{1.944301in}}%
\pgfpathlineto{\pgfqpoint{3.678988in}{1.948846in}}%
\pgfpathlineto{\pgfqpoint{3.698384in}{1.950230in}}%
\pgfpathlineto{\pgfqpoint{3.717781in}{1.949228in}}%
\pgfpathlineto{\pgfqpoint{3.737178in}{1.945662in}}%
\pgfpathlineto{\pgfqpoint{3.785669in}{1.933714in}}%
\pgfpathlineto{\pgfqpoint{3.805066in}{1.932019in}}%
\pgfpathlineto{\pgfqpoint{3.824462in}{1.932974in}}%
\pgfpathlineto{\pgfqpoint{3.882652in}{1.939829in}}%
\pgfpathlineto{\pgfqpoint{3.911747in}{1.939382in}}%
\pgfpathlineto{\pgfqpoint{3.960239in}{1.937602in}}%
\pgfpathlineto{\pgfqpoint{4.047524in}{1.939671in}}%
\pgfpathlineto{\pgfqpoint{4.105714in}{1.936694in}}%
\pgfpathlineto{\pgfqpoint{4.144507in}{1.938245in}}%
\pgfpathlineto{\pgfqpoint{4.183300in}{1.939681in}}%
\pgfpathlineto{\pgfqpoint{4.222094in}{1.938446in}}%
\pgfpathlineto{\pgfqpoint{4.260887in}{1.937446in}}%
\pgfpathlineto{\pgfqpoint{4.377267in}{1.938202in}}%
\pgfpathlineto{\pgfqpoint{4.416060in}{1.937612in}}%
\pgfpathlineto{\pgfqpoint{4.522742in}{1.938444in}}%
\pgfpathlineto{\pgfqpoint{4.571233in}{1.937842in}}%
\pgfpathlineto{\pgfqpoint{4.668216in}{1.938571in}}%
\pgfpathlineto{\pgfqpoint{4.726406in}{1.937999in}}%
\pgfpathlineto{\pgfqpoint{4.813691in}{1.938659in}}%
\pgfpathlineto{\pgfqpoint{4.881579in}{1.938127in}}%
\pgfpathlineto{\pgfqpoint{4.968864in}{1.938557in}}%
\pgfpathlineto{\pgfqpoint{5.036753in}{1.938228in}}%
\pgfpathlineto{\pgfqpoint{5.124038in}{1.938481in}}%
\pgfpathlineto{\pgfqpoint{5.201624in}{1.938435in}}%
\pgfpathlineto{\pgfqpoint{5.279211in}{1.938424in}}%
\pgfpathlineto{\pgfqpoint{5.356797in}{1.938483in}}%
\pgfpathlineto{\pgfqpoint{5.444082in}{1.938293in}}%
\pgfpathlineto{\pgfqpoint{5.618652in}{1.938228in}}%
\pgfpathlineto{\pgfqpoint{5.618652in}{1.938228in}}%
\pgfusepath{stroke}%
\end{pgfscope}%
\begin{pgfscope}%
\pgfpathrectangle{\pgfqpoint{0.781757in}{0.456901in}}{\pgfqpoint{5.078868in}{2.816432in}}%
\pgfusepath{clip}%
\pgfsetbuttcap%
\pgfsetroundjoin%
\pgfsetlinewidth{2.007500pt}%
\definecolor{currentstroke}{rgb}{0.835294,0.368627,0.000000}%
\pgfsetstrokecolor{currentstroke}%
\pgfsetdash{{7.400000pt}{3.200000pt}}{0.000000pt}%
\pgfpathmoveto{\pgfqpoint{0.779190in}{1.938442in}}%
\pgfpathlineto{\pgfqpoint{0.788888in}{2.182964in}}%
\pgfpathlineto{\pgfqpoint{0.798586in}{2.341602in}}%
\pgfpathlineto{\pgfqpoint{0.808285in}{2.402137in}}%
\pgfpathlineto{\pgfqpoint{0.817983in}{2.426045in}}%
\pgfpathlineto{\pgfqpoint{0.827681in}{2.444092in}}%
\pgfpathlineto{\pgfqpoint{0.837380in}{2.470306in}}%
\pgfpathlineto{\pgfqpoint{0.847078in}{2.509605in}}%
\pgfpathlineto{\pgfqpoint{0.856776in}{2.562014in}}%
\pgfpathlineto{\pgfqpoint{0.876173in}{2.695103in}}%
\pgfpathlineto{\pgfqpoint{0.905268in}{2.909284in}}%
\pgfpathlineto{\pgfqpoint{0.914966in}{2.971224in}}%
\pgfpathlineto{\pgfqpoint{0.924664in}{3.024591in}}%
\pgfpathlineto{\pgfqpoint{0.934363in}{3.068118in}}%
\pgfpathlineto{\pgfqpoint{0.944061in}{3.101210in}}%
\pgfpathlineto{\pgfqpoint{0.953759in}{3.123866in}}%
\pgfpathlineto{\pgfqpoint{0.963458in}{3.136597in}}%
\pgfpathlineto{\pgfqpoint{0.973156in}{3.140319in}}%
\pgfpathlineto{\pgfqpoint{0.982854in}{3.136249in}}%
\pgfpathlineto{\pgfqpoint{0.992553in}{3.125795in}}%
\pgfpathlineto{\pgfqpoint{1.002251in}{3.110461in}}%
\pgfpathlineto{\pgfqpoint{1.021648in}{3.071127in}}%
\pgfpathlineto{\pgfqpoint{1.050743in}{3.009892in}}%
\pgfpathlineto{\pgfqpoint{1.060441in}{2.992807in}}%
\pgfpathlineto{\pgfqpoint{1.070139in}{2.978403in}}%
\pgfpathlineto{\pgfqpoint{1.079838in}{2.966971in}}%
\pgfpathlineto{\pgfqpoint{1.089536in}{2.958611in}}%
\pgfpathlineto{\pgfqpoint{1.099234in}{2.953258in}}%
\pgfpathlineto{\pgfqpoint{1.108933in}{2.950707in}}%
\pgfpathlineto{\pgfqpoint{1.118631in}{2.950651in}}%
\pgfpathlineto{\pgfqpoint{1.128329in}{2.952703in}}%
\pgfpathlineto{\pgfqpoint{1.138027in}{2.956437in}}%
\pgfpathlineto{\pgfqpoint{1.157424in}{2.967180in}}%
\pgfpathlineto{\pgfqpoint{1.196217in}{2.990847in}}%
\pgfpathlineto{\pgfqpoint{1.215614in}{2.999377in}}%
\pgfpathlineto{\pgfqpoint{1.235011in}{3.004402in}}%
\pgfpathlineto{\pgfqpoint{1.254407in}{3.006026in}}%
\pgfpathlineto{\pgfqpoint{1.273804in}{3.004931in}}%
\pgfpathlineto{\pgfqpoint{1.302899in}{3.000338in}}%
\pgfpathlineto{\pgfqpoint{1.351391in}{2.992289in}}%
\pgfpathlineto{\pgfqpoint{1.380486in}{2.990025in}}%
\pgfpathlineto{\pgfqpoint{1.409580in}{2.989895in}}%
\pgfpathlineto{\pgfqpoint{1.467770in}{2.992672in}}%
\pgfpathlineto{\pgfqpoint{1.525960in}{2.994458in}}%
\pgfpathlineto{\pgfqpoint{1.603547in}{2.993727in}}%
\pgfpathlineto{\pgfqpoint{1.710228in}{2.993176in}}%
\pgfpathlineto{\pgfqpoint{1.981781in}{2.993398in}}%
\pgfpathlineto{\pgfqpoint{1.991480in}{2.948797in}}%
\pgfpathlineto{\pgfqpoint{2.001178in}{2.371554in}}%
\pgfpathlineto{\pgfqpoint{2.010876in}{2.135079in}}%
\pgfpathlineto{\pgfqpoint{2.020575in}{2.046802in}}%
\pgfpathlineto{\pgfqpoint{2.039971in}{1.968042in}}%
\pgfpathlineto{\pgfqpoint{2.049670in}{1.907816in}}%
\pgfpathlineto{\pgfqpoint{2.059368in}{1.820627in}}%
\pgfpathlineto{\pgfqpoint{2.069066in}{1.708406in}}%
\pgfpathlineto{\pgfqpoint{2.088463in}{1.434155in}}%
\pgfpathlineto{\pgfqpoint{2.117558in}{1.010936in}}%
\pgfpathlineto{\pgfqpoint{2.127256in}{0.892118in}}%
\pgfpathlineto{\pgfqpoint{2.136955in}{0.791398in}}%
\pgfpathlineto{\pgfqpoint{2.146653in}{0.710872in}}%
\pgfpathlineto{\pgfqpoint{2.156351in}{0.651339in}}%
\pgfpathlineto{\pgfqpoint{2.166050in}{0.612455in}}%
\pgfpathlineto{\pgfqpoint{2.175748in}{0.592919in}}%
\pgfpathlineto{\pgfqpoint{2.185446in}{0.590687in}}%
\pgfpathlineto{\pgfqpoint{2.195144in}{0.603181in}}%
\pgfpathlineto{\pgfqpoint{2.204843in}{0.627503in}}%
\pgfpathlineto{\pgfqpoint{2.214541in}{0.660624in}}%
\pgfpathlineto{\pgfqpoint{2.233938in}{0.741471in}}%
\pgfpathlineto{\pgfqpoint{2.253334in}{0.824569in}}%
\pgfpathlineto{\pgfqpoint{2.263033in}{0.861854in}}%
\pgfpathlineto{\pgfqpoint{2.272731in}{0.894416in}}%
\pgfpathlineto{\pgfqpoint{2.282429in}{0.921389in}}%
\pgfpathlineto{\pgfqpoint{2.292128in}{0.942316in}}%
\pgfpathlineto{\pgfqpoint{2.301826in}{0.957106in}}%
\pgfpathlineto{\pgfqpoint{2.311524in}{0.965985in}}%
\pgfpathlineto{\pgfqpoint{2.321223in}{0.969434in}}%
\pgfpathlineto{\pgfqpoint{2.330921in}{0.968126in}}%
\pgfpathlineto{\pgfqpoint{2.340619in}{0.962862in}}%
\pgfpathlineto{\pgfqpoint{2.350318in}{0.954516in}}%
\pgfpathlineto{\pgfqpoint{2.369714in}{0.932098in}}%
\pgfpathlineto{\pgfqpoint{2.398809in}{0.895853in}}%
\pgfpathlineto{\pgfqpoint{2.418206in}{0.876633in}}%
\pgfpathlineto{\pgfqpoint{2.427904in}{0.869486in}}%
\pgfpathlineto{\pgfqpoint{2.437603in}{0.864135in}}%
\pgfpathlineto{\pgfqpoint{2.447301in}{0.860562in}}%
\pgfpathlineto{\pgfqpoint{2.456999in}{0.858669in}}%
\pgfpathlineto{\pgfqpoint{2.466697in}{0.858287in}}%
\pgfpathlineto{\pgfqpoint{2.486094in}{0.861163in}}%
\pgfpathlineto{\pgfqpoint{2.505491in}{0.867204in}}%
\pgfpathlineto{\pgfqpoint{2.553982in}{0.884046in}}%
\pgfpathlineto{\pgfqpoint{2.573379in}{0.888297in}}%
\pgfpathlineto{\pgfqpoint{2.592776in}{0.890467in}}%
\pgfpathlineto{\pgfqpoint{2.612172in}{0.890764in}}%
\pgfpathlineto{\pgfqpoint{2.641267in}{0.888787in}}%
\pgfpathlineto{\pgfqpoint{2.709156in}{0.882338in}}%
\pgfpathlineto{\pgfqpoint{2.747949in}{0.881299in}}%
\pgfpathlineto{\pgfqpoint{2.806139in}{0.882644in}}%
\pgfpathlineto{\pgfqpoint{2.883725in}{0.884092in}}%
\pgfpathlineto{\pgfqpoint{3.194072in}{0.883509in}}%
\pgfpathlineto{\pgfqpoint{3.203770in}{1.024230in}}%
\pgfpathlineto{\pgfqpoint{3.213468in}{1.244407in}}%
\pgfpathlineto{\pgfqpoint{3.223167in}{1.331583in}}%
\pgfpathlineto{\pgfqpoint{3.232865in}{1.364137in}}%
\pgfpathlineto{\pgfqpoint{3.242563in}{1.382026in}}%
\pgfpathlineto{\pgfqpoint{3.252262in}{1.404294in}}%
\pgfpathlineto{\pgfqpoint{3.261960in}{1.438580in}}%
\pgfpathlineto{\pgfqpoint{3.271658in}{1.486357in}}%
\pgfpathlineto{\pgfqpoint{3.281356in}{1.545864in}}%
\pgfpathlineto{\pgfqpoint{3.300753in}{1.686247in}}%
\pgfpathlineto{\pgfqpoint{3.320150in}{1.829729in}}%
\pgfpathlineto{\pgfqpoint{3.329848in}{1.894365in}}%
\pgfpathlineto{\pgfqpoint{3.339546in}{1.951069in}}%
\pgfpathlineto{\pgfqpoint{3.349245in}{1.998320in}}%
\pgfpathlineto{\pgfqpoint{3.358943in}{2.035280in}}%
\pgfpathlineto{\pgfqpoint{3.368641in}{2.061740in}}%
\pgfpathlineto{\pgfqpoint{3.378340in}{2.078035in}}%
\pgfpathlineto{\pgfqpoint{3.388038in}{2.084944in}}%
\pgfpathlineto{\pgfqpoint{3.397736in}{2.083584in}}%
\pgfpathlineto{\pgfqpoint{3.407435in}{2.075307in}}%
\pgfpathlineto{\pgfqpoint{3.417133in}{2.061591in}}%
\pgfpathlineto{\pgfqpoint{3.426831in}{2.043955in}}%
\pgfpathlineto{\pgfqpoint{3.475323in}{1.943822in}}%
\pgfpathlineto{\pgfqpoint{3.485021in}{1.928387in}}%
\pgfpathlineto{\pgfqpoint{3.494720in}{1.915841in}}%
\pgfpathlineto{\pgfqpoint{3.504418in}{1.906350in}}%
\pgfpathlineto{\pgfqpoint{3.514116in}{1.899908in}}%
\pgfpathlineto{\pgfqpoint{3.523814in}{1.896358in}}%
\pgfpathlineto{\pgfqpoint{3.533513in}{1.895425in}}%
\pgfpathlineto{\pgfqpoint{3.543211in}{1.896749in}}%
\pgfpathlineto{\pgfqpoint{3.552909in}{1.899916in}}%
\pgfpathlineto{\pgfqpoint{3.572306in}{1.910010in}}%
\pgfpathlineto{\pgfqpoint{3.611099in}{1.933958in}}%
\pgfpathlineto{\pgfqpoint{3.630496in}{1.943082in}}%
\pgfpathlineto{\pgfqpoint{3.649893in}{1.948765in}}%
\pgfpathlineto{\pgfqpoint{3.669289in}{1.950976in}}%
\pgfpathlineto{\pgfqpoint{3.688686in}{1.950314in}}%
\pgfpathlineto{\pgfqpoint{3.717781in}{1.946017in}}%
\pgfpathlineto{\pgfqpoint{3.766273in}{1.937747in}}%
\pgfpathlineto{\pgfqpoint{3.795367in}{1.935197in}}%
\pgfpathlineto{\pgfqpoint{3.824462in}{1.934839in}}%
\pgfpathlineto{\pgfqpoint{3.872954in}{1.936984in}}%
\pgfpathlineto{\pgfqpoint{3.931144in}{1.939305in}}%
\pgfpathlineto{\pgfqpoint{3.989334in}{1.939257in}}%
\pgfpathlineto{\pgfqpoint{4.125110in}{1.938182in}}%
\pgfpathlineto{\pgfqpoint{4.464552in}{1.938442in}}%
\pgfpathlineto{\pgfqpoint{5.618652in}{1.938442in}}%
\pgfpathlineto{\pgfqpoint{5.618652in}{1.938442in}}%
\pgfusepath{stroke}%
\end{pgfscope}%
\begin{pgfscope}%
\pgfsetrectcap%
\pgfsetmiterjoin%
\pgfsetlinewidth{0.602250pt}%
\definecolor{currentstroke}{rgb}{0.000000,0.000000,0.000000}%
\pgfsetstrokecolor{currentstroke}%
\pgfsetdash{}{0pt}%
\pgfpathmoveto{\pgfqpoint{0.781757in}{0.456901in}}%
\pgfpathlineto{\pgfqpoint{0.781757in}{3.273333in}}%
\pgfusepath{stroke}%
\end{pgfscope}%
\begin{pgfscope}%
\pgfsetrectcap%
\pgfsetmiterjoin%
\pgfsetlinewidth{0.602250pt}%
\definecolor{currentstroke}{rgb}{0.000000,0.000000,0.000000}%
\pgfsetstrokecolor{currentstroke}%
\pgfsetdash{}{0pt}%
\pgfpathmoveto{\pgfqpoint{0.781757in}{0.456901in}}%
\pgfpathlineto{\pgfqpoint{5.860625in}{0.456901in}}%
\pgfusepath{stroke}%
\end{pgfscope}%
\begin{pgfscope}%
\pgfsetrectcap%
\pgfsetroundjoin%
\pgfsetlinewidth{2.007500pt}%
\definecolor{currentstroke}{rgb}{0.000000,0.000000,0.000000}%
\pgfsetstrokecolor{currentstroke}%
\pgfsetdash{}{0pt}%
\pgfpathmoveto{\pgfqpoint{0.116667in}{3.565000in}}%
\pgfpathlineto{\pgfqpoint{0.366667in}{3.565000in}}%
\pgfpathlineto{\pgfqpoint{0.616667in}{3.565000in}}%
\pgfusepath{stroke}%
\end{pgfscope}%
\begin{pgfscope}%
\definecolor{textcolor}{rgb}{0.000000,0.000000,0.000000}%
\pgfsetstrokecolor{textcolor}%
\pgfsetfillcolor{textcolor}%
\pgftext[x=0.750000in,y=3.506667in,left,base]{\color{textcolor}{\rmfamily\fontsize{12.000000}{14.400000}\selectfont\catcode`\^=\active\def^{\ifmmode\sp\else\^{}\fi}\catcode`\%=\active\def%{\%}Приближенный расчет по спектру}}%
\end{pgfscope}%
\begin{pgfscope}%
\pgfsetbuttcap%
\pgfsetroundjoin%
\pgfsetlinewidth{2.007500pt}%
\definecolor{currentstroke}{rgb}{0.835294,0.368627,0.000000}%
\pgfsetstrokecolor{currentstroke}%
\pgfsetdash{{7.400000pt}{3.200000pt}}{0.000000pt}%
\pgfpathmoveto{\pgfqpoint{3.507080in}{3.565000in}}%
\pgfpathlineto{\pgfqpoint{3.757080in}{3.565000in}}%
\pgfpathlineto{\pgfqpoint{4.007080in}{3.565000in}}%
\pgfusepath{stroke}%
\end{pgfscope}%
\begin{pgfscope}%
\definecolor{textcolor}{rgb}{0.000000,0.000000,0.000000}%
\pgfsetstrokecolor{textcolor}%
\pgfsetfillcolor{textcolor}%
\pgftext[x=4.140413in,y=3.506667in,left,base]{\color{textcolor}{\rmfamily\fontsize{12.000000}{14.400000}\selectfont\catcode`\^=\active\def^{\ifmmode\sp\else\^{}\fi}\catcode`\%=\active\def%{\%}Точный расчет (разд. 6). $t_\mathrm{и}=47.1$}}%
\end{pgfscope}%
\end{pgfpicture}%
\makeatother%
\endgroup%

    \caption{Сравнение точного и приближенного расчета реакции цепи для $t_{\text{и}} = \pv[.2]{n['ti1']}$}
    \label{fig:plot_approx_reaction_1}
\end{figure}

\begin{figure}[H]
    \centering
    %% Creator: Matplotlib, PGF backend
%%
%% To include the figure in your LaTeX document, write
%%   \input{<filename>.pgf}
%%
%% Make sure the required packages are loaded in your preamble
%%   \usepackage{pgf}
%%
%% Also ensure that all the required font packages are loaded; for instance,
%% the lmodern package is sometimes necessary when using math font.
%%   \usepackage{lmodern}
%%
%% Figures using additional raster images can only be included by \input if
%% they are in the same directory as the main LaTeX file. For loading figures
%% from other directories you can use the `import` package
%%   \usepackage{import}
%%
%% and then include the figures with
%%   \import{<path to file>}{<filename>.pgf}
%%
%% Matplotlib used the following preamble
%%   \def\mathdefault#1{#1}
%%   \everymath=\expandafter{\the\everymath\displaystyle}
%%   \IfFileExists{scrextend.sty}{
%%     \usepackage[fontsize=12.000000pt]{scrextend}
%%   }{
%%     \renewcommand{\normalsize}{\fontsize{12.000000}{14.400000}\selectfont}
%%     \normalsize
%%   }
%%   \usepackage{fontspec}\setmainfont{Times New Roman}\usepackage[english,russian]{babel}\usepackage{amsmath}
%%   \ifdefined\pdftexversion\else  % non-pdftex case.
%%     \usepackage{fontspec}
%%     \setmainfont{times.ttf}[Path=\detokenize{/usr/local/share/fonts/truetype/times-new-roman/}]
%%     \setsansfont{DejaVuSans.ttf}[Path=\detokenize{/usr/share/matplotlib/mpl-data/fonts/ttf/}]
%%     \setmonofont{DejaVuSansMono.ttf}[Path=\detokenize{/usr/share/matplotlib/mpl-data/fonts/ttf/}]
%%   \fi
%%   \makeatletter\@ifpackageloaded{underscore}{}{\usepackage[strings]{underscore}}\makeatother
%%
\begingroup%
\makeatletter%
\begin{pgfpicture}%
\pgfpathrectangle{\pgfpointorigin}{\pgfqpoint{6.642382in}{3.740000in}}%
\pgfusepath{use as bounding box, clip}%
\begin{pgfscope}%
\pgfsetbuttcap%
\pgfsetmiterjoin%
\definecolor{currentfill}{rgb}{1.000000,1.000000,1.000000}%
\pgfsetfillcolor{currentfill}%
\pgfsetlinewidth{0.000000pt}%
\definecolor{currentstroke}{rgb}{1.000000,1.000000,1.000000}%
\pgfsetstrokecolor{currentstroke}%
\pgfsetdash{}{0pt}%
\pgfpathmoveto{\pgfqpoint{0.000000in}{0.000000in}}%
\pgfpathlineto{\pgfqpoint{6.642382in}{0.000000in}}%
\pgfpathlineto{\pgfqpoint{6.642382in}{3.740000in}}%
\pgfpathlineto{\pgfqpoint{0.000000in}{3.740000in}}%
\pgfpathlineto{\pgfqpoint{0.000000in}{0.000000in}}%
\pgfpathclose%
\pgfusepath{fill}%
\end{pgfscope}%
\begin{pgfscope}%
\pgfsetbuttcap%
\pgfsetmiterjoin%
\definecolor{currentfill}{rgb}{1.000000,1.000000,1.000000}%
\pgfsetfillcolor{currentfill}%
\pgfsetlinewidth{0.000000pt}%
\definecolor{currentstroke}{rgb}{0.000000,0.000000,0.000000}%
\pgfsetstrokecolor{currentstroke}%
\pgfsetstrokeopacity{0.000000}%
\pgfsetdash{}{0pt}%
\pgfpathmoveto{\pgfqpoint{0.781757in}{0.456901in}}%
\pgfpathlineto{\pgfqpoint{5.860625in}{0.456901in}}%
\pgfpathlineto{\pgfqpoint{5.860625in}{3.273333in}}%
\pgfpathlineto{\pgfqpoint{0.781757in}{3.273333in}}%
\pgfpathlineto{\pgfqpoint{0.781757in}{0.456901in}}%
\pgfpathclose%
\pgfusepath{fill}%
\end{pgfscope}%
\begin{pgfscope}%
\pgfpathrectangle{\pgfqpoint{0.781757in}{0.456901in}}{\pgfqpoint{5.078868in}{2.816432in}}%
\pgfusepath{clip}%
\pgfsetbuttcap%
\pgfsetroundjoin%
\pgfsetlinewidth{0.501875pt}%
\definecolor{currentstroke}{rgb}{0.501961,0.501961,0.501961}%
\pgfsetstrokecolor{currentstroke}%
\pgfsetstrokeopacity{0.400000}%
\pgfsetdash{{1.850000pt}{0.800000pt}}{0.000000pt}%
\pgfpathmoveto{\pgfqpoint{0.781757in}{0.456901in}}%
\pgfpathlineto{\pgfqpoint{0.781757in}{3.273333in}}%
\pgfusepath{stroke}%
\end{pgfscope}%
\begin{pgfscope}%
\pgfsetbuttcap%
\pgfsetroundjoin%
\definecolor{currentfill}{rgb}{1.000000,1.000000,1.000000}%
\pgfsetfillcolor{currentfill}%
\pgfsetlinewidth{0.803000pt}%
\definecolor{currentstroke}{rgb}{0.000000,0.000000,0.000000}%
\pgfsetstrokecolor{currentstroke}%
\pgfsetdash{}{0pt}%
\pgfsys@defobject{currentmarker}{\pgfqpoint{0.000000in}{0.000000in}}{\pgfqpoint{0.000000in}{0.048611in}}{%
\pgfpathmoveto{\pgfqpoint{0.000000in}{0.000000in}}%
\pgfpathlineto{\pgfqpoint{0.000000in}{0.048611in}}%
\pgfusepath{stroke,fill}%
}%
\begin{pgfscope}%
\pgfsys@transformshift{0.781757in}{0.456901in}%
\pgfsys@useobject{currentmarker}{}%
\end{pgfscope}%
\end{pgfscope}%
\begin{pgfscope}%
\definecolor{textcolor}{rgb}{0.000000,0.000000,0.000000}%
\pgfsetstrokecolor{textcolor}%
\pgfsetfillcolor{textcolor}%
\pgftext[x=0.781757in,y=0.408290in,,top]{\color{textcolor}{\rmfamily\fontsize{12.000000}{14.400000}\selectfont\catcode`\^=\active\def^{\ifmmode\sp\else\^{}\fi}\catcode`\%=\active\def%{\%}$\mathdefault{0}$}}%
\end{pgfscope}%
\begin{pgfscope}%
\pgfpathrectangle{\pgfqpoint{0.781757in}{0.456901in}}{\pgfqpoint{5.078868in}{2.816432in}}%
\pgfusepath{clip}%
\pgfsetbuttcap%
\pgfsetroundjoin%
\pgfsetlinewidth{0.501875pt}%
\definecolor{currentstroke}{rgb}{0.501961,0.501961,0.501961}%
\pgfsetstrokecolor{currentstroke}%
\pgfsetstrokeopacity{0.400000}%
\pgfsetdash{{1.850000pt}{0.800000pt}}{0.000000pt}%
\pgfpathmoveto{\pgfqpoint{1.637491in}{0.456901in}}%
\pgfpathlineto{\pgfqpoint{1.637491in}{3.273333in}}%
\pgfusepath{stroke}%
\end{pgfscope}%
\begin{pgfscope}%
\pgfsetbuttcap%
\pgfsetroundjoin%
\definecolor{currentfill}{rgb}{1.000000,1.000000,1.000000}%
\pgfsetfillcolor{currentfill}%
\pgfsetlinewidth{0.803000pt}%
\definecolor{currentstroke}{rgb}{0.000000,0.000000,0.000000}%
\pgfsetstrokecolor{currentstroke}%
\pgfsetdash{}{0pt}%
\pgfsys@defobject{currentmarker}{\pgfqpoint{0.000000in}{0.000000in}}{\pgfqpoint{0.000000in}{0.048611in}}{%
\pgfpathmoveto{\pgfqpoint{0.000000in}{0.000000in}}%
\pgfpathlineto{\pgfqpoint{0.000000in}{0.048611in}}%
\pgfusepath{stroke,fill}%
}%
\begin{pgfscope}%
\pgfsys@transformshift{1.637491in}{0.456901in}%
\pgfsys@useobject{currentmarker}{}%
\end{pgfscope}%
\end{pgfscope}%
\begin{pgfscope}%
\definecolor{textcolor}{rgb}{0.000000,0.000000,0.000000}%
\pgfsetstrokecolor{textcolor}%
\pgfsetfillcolor{textcolor}%
\pgftext[x=1.637491in,y=0.408290in,,top]{\color{textcolor}{\rmfamily\fontsize{12.000000}{14.400000}\selectfont\catcode`\^=\active\def^{\ifmmode\sp\else\^{}\fi}\catcode`\%=\active\def%{\%}$\mathdefault{5}$}}%
\end{pgfscope}%
\begin{pgfscope}%
\pgfpathrectangle{\pgfqpoint{0.781757in}{0.456901in}}{\pgfqpoint{5.078868in}{2.816432in}}%
\pgfusepath{clip}%
\pgfsetbuttcap%
\pgfsetroundjoin%
\pgfsetlinewidth{0.501875pt}%
\definecolor{currentstroke}{rgb}{0.501961,0.501961,0.501961}%
\pgfsetstrokecolor{currentstroke}%
\pgfsetstrokeopacity{0.400000}%
\pgfsetdash{{1.850000pt}{0.800000pt}}{0.000000pt}%
\pgfpathmoveto{\pgfqpoint{2.493225in}{0.456901in}}%
\pgfpathlineto{\pgfqpoint{2.493225in}{3.273333in}}%
\pgfusepath{stroke}%
\end{pgfscope}%
\begin{pgfscope}%
\pgfsetbuttcap%
\pgfsetroundjoin%
\definecolor{currentfill}{rgb}{1.000000,1.000000,1.000000}%
\pgfsetfillcolor{currentfill}%
\pgfsetlinewidth{0.803000pt}%
\definecolor{currentstroke}{rgb}{0.000000,0.000000,0.000000}%
\pgfsetstrokecolor{currentstroke}%
\pgfsetdash{}{0pt}%
\pgfsys@defobject{currentmarker}{\pgfqpoint{0.000000in}{0.000000in}}{\pgfqpoint{0.000000in}{0.048611in}}{%
\pgfpathmoveto{\pgfqpoint{0.000000in}{0.000000in}}%
\pgfpathlineto{\pgfqpoint{0.000000in}{0.048611in}}%
\pgfusepath{stroke,fill}%
}%
\begin{pgfscope}%
\pgfsys@transformshift{2.493225in}{0.456901in}%
\pgfsys@useobject{currentmarker}{}%
\end{pgfscope}%
\end{pgfscope}%
\begin{pgfscope}%
\definecolor{textcolor}{rgb}{0.000000,0.000000,0.000000}%
\pgfsetstrokecolor{textcolor}%
\pgfsetfillcolor{textcolor}%
\pgftext[x=2.493225in,y=0.408290in,,top]{\color{textcolor}{\rmfamily\fontsize{12.000000}{14.400000}\selectfont\catcode`\^=\active\def^{\ifmmode\sp\else\^{}\fi}\catcode`\%=\active\def%{\%}$\mathdefault{10}$}}%
\end{pgfscope}%
\begin{pgfscope}%
\pgfpathrectangle{\pgfqpoint{0.781757in}{0.456901in}}{\pgfqpoint{5.078868in}{2.816432in}}%
\pgfusepath{clip}%
\pgfsetbuttcap%
\pgfsetroundjoin%
\pgfsetlinewidth{0.501875pt}%
\definecolor{currentstroke}{rgb}{0.501961,0.501961,0.501961}%
\pgfsetstrokecolor{currentstroke}%
\pgfsetstrokeopacity{0.400000}%
\pgfsetdash{{1.850000pt}{0.800000pt}}{0.000000pt}%
\pgfpathmoveto{\pgfqpoint{3.348960in}{0.456901in}}%
\pgfpathlineto{\pgfqpoint{3.348960in}{3.273333in}}%
\pgfusepath{stroke}%
\end{pgfscope}%
\begin{pgfscope}%
\pgfsetbuttcap%
\pgfsetroundjoin%
\definecolor{currentfill}{rgb}{1.000000,1.000000,1.000000}%
\pgfsetfillcolor{currentfill}%
\pgfsetlinewidth{0.803000pt}%
\definecolor{currentstroke}{rgb}{0.000000,0.000000,0.000000}%
\pgfsetstrokecolor{currentstroke}%
\pgfsetdash{}{0pt}%
\pgfsys@defobject{currentmarker}{\pgfqpoint{0.000000in}{0.000000in}}{\pgfqpoint{0.000000in}{0.048611in}}{%
\pgfpathmoveto{\pgfqpoint{0.000000in}{0.000000in}}%
\pgfpathlineto{\pgfqpoint{0.000000in}{0.048611in}}%
\pgfusepath{stroke,fill}%
}%
\begin{pgfscope}%
\pgfsys@transformshift{3.348960in}{0.456901in}%
\pgfsys@useobject{currentmarker}{}%
\end{pgfscope}%
\end{pgfscope}%
\begin{pgfscope}%
\definecolor{textcolor}{rgb}{0.000000,0.000000,0.000000}%
\pgfsetstrokecolor{textcolor}%
\pgfsetfillcolor{textcolor}%
\pgftext[x=3.348960in,y=0.408290in,,top]{\color{textcolor}{\rmfamily\fontsize{12.000000}{14.400000}\selectfont\catcode`\^=\active\def^{\ifmmode\sp\else\^{}\fi}\catcode`\%=\active\def%{\%}$\mathdefault{15}$}}%
\end{pgfscope}%
\begin{pgfscope}%
\pgfpathrectangle{\pgfqpoint{0.781757in}{0.456901in}}{\pgfqpoint{5.078868in}{2.816432in}}%
\pgfusepath{clip}%
\pgfsetbuttcap%
\pgfsetroundjoin%
\pgfsetlinewidth{0.501875pt}%
\definecolor{currentstroke}{rgb}{0.501961,0.501961,0.501961}%
\pgfsetstrokecolor{currentstroke}%
\pgfsetstrokeopacity{0.400000}%
\pgfsetdash{{1.850000pt}{0.800000pt}}{0.000000pt}%
\pgfpathmoveto{\pgfqpoint{4.204694in}{0.456901in}}%
\pgfpathlineto{\pgfqpoint{4.204694in}{3.273333in}}%
\pgfusepath{stroke}%
\end{pgfscope}%
\begin{pgfscope}%
\pgfsetbuttcap%
\pgfsetroundjoin%
\definecolor{currentfill}{rgb}{1.000000,1.000000,1.000000}%
\pgfsetfillcolor{currentfill}%
\pgfsetlinewidth{0.803000pt}%
\definecolor{currentstroke}{rgb}{0.000000,0.000000,0.000000}%
\pgfsetstrokecolor{currentstroke}%
\pgfsetdash{}{0pt}%
\pgfsys@defobject{currentmarker}{\pgfqpoint{0.000000in}{0.000000in}}{\pgfqpoint{0.000000in}{0.048611in}}{%
\pgfpathmoveto{\pgfqpoint{0.000000in}{0.000000in}}%
\pgfpathlineto{\pgfqpoint{0.000000in}{0.048611in}}%
\pgfusepath{stroke,fill}%
}%
\begin{pgfscope}%
\pgfsys@transformshift{4.204694in}{0.456901in}%
\pgfsys@useobject{currentmarker}{}%
\end{pgfscope}%
\end{pgfscope}%
\begin{pgfscope}%
\definecolor{textcolor}{rgb}{0.000000,0.000000,0.000000}%
\pgfsetstrokecolor{textcolor}%
\pgfsetfillcolor{textcolor}%
\pgftext[x=4.204694in,y=0.408290in,,top]{\color{textcolor}{\rmfamily\fontsize{12.000000}{14.400000}\selectfont\catcode`\^=\active\def^{\ifmmode\sp\else\^{}\fi}\catcode`\%=\active\def%{\%}$\mathdefault{20}$}}%
\end{pgfscope}%
\begin{pgfscope}%
\pgfpathrectangle{\pgfqpoint{0.781757in}{0.456901in}}{\pgfqpoint{5.078868in}{2.816432in}}%
\pgfusepath{clip}%
\pgfsetbuttcap%
\pgfsetroundjoin%
\pgfsetlinewidth{0.501875pt}%
\definecolor{currentstroke}{rgb}{0.501961,0.501961,0.501961}%
\pgfsetstrokecolor{currentstroke}%
\pgfsetstrokeopacity{0.400000}%
\pgfsetdash{{1.850000pt}{0.800000pt}}{0.000000pt}%
\pgfpathmoveto{\pgfqpoint{5.060428in}{0.456901in}}%
\pgfpathlineto{\pgfqpoint{5.060428in}{3.273333in}}%
\pgfusepath{stroke}%
\end{pgfscope}%
\begin{pgfscope}%
\pgfsetbuttcap%
\pgfsetroundjoin%
\definecolor{currentfill}{rgb}{1.000000,1.000000,1.000000}%
\pgfsetfillcolor{currentfill}%
\pgfsetlinewidth{0.803000pt}%
\definecolor{currentstroke}{rgb}{0.000000,0.000000,0.000000}%
\pgfsetstrokecolor{currentstroke}%
\pgfsetdash{}{0pt}%
\pgfsys@defobject{currentmarker}{\pgfqpoint{0.000000in}{0.000000in}}{\pgfqpoint{0.000000in}{0.048611in}}{%
\pgfpathmoveto{\pgfqpoint{0.000000in}{0.000000in}}%
\pgfpathlineto{\pgfqpoint{0.000000in}{0.048611in}}%
\pgfusepath{stroke,fill}%
}%
\begin{pgfscope}%
\pgfsys@transformshift{5.060428in}{0.456901in}%
\pgfsys@useobject{currentmarker}{}%
\end{pgfscope}%
\end{pgfscope}%
\begin{pgfscope}%
\definecolor{textcolor}{rgb}{0.000000,0.000000,0.000000}%
\pgfsetstrokecolor{textcolor}%
\pgfsetfillcolor{textcolor}%
\pgftext[x=5.060428in,y=0.408290in,,top]{\color{textcolor}{\rmfamily\fontsize{12.000000}{14.400000}\selectfont\catcode`\^=\active\def^{\ifmmode\sp\else\^{}\fi}\catcode`\%=\active\def%{\%}$\mathdefault{25}$}}%
\end{pgfscope}%
\begin{pgfscope}%
\pgfpathrectangle{\pgfqpoint{0.781757in}{0.456901in}}{\pgfqpoint{5.078868in}{2.816432in}}%
\pgfusepath{clip}%
\pgfsetbuttcap%
\pgfsetroundjoin%
\pgfsetlinewidth{0.501875pt}%
\definecolor{currentstroke}{rgb}{0.501961,0.501961,0.501961}%
\pgfsetstrokecolor{currentstroke}%
\pgfsetstrokeopacity{0.400000}%
\pgfsetdash{{1.850000pt}{0.800000pt}}{0.000000pt}%
\pgfpathmoveto{\pgfqpoint{0.952904in}{0.456901in}}%
\pgfpathlineto{\pgfqpoint{0.952904in}{3.273333in}}%
\pgfusepath{stroke}%
\end{pgfscope}%
\begin{pgfscope}%
\pgfsetbuttcap%
\pgfsetroundjoin%
\definecolor{currentfill}{rgb}{1.000000,1.000000,1.000000}%
\pgfsetfillcolor{currentfill}%
\pgfsetlinewidth{0.602250pt}%
\definecolor{currentstroke}{rgb}{0.000000,0.000000,0.000000}%
\pgfsetstrokecolor{currentstroke}%
\pgfsetdash{}{0pt}%
\pgfsys@defobject{currentmarker}{\pgfqpoint{0.000000in}{0.000000in}}{\pgfqpoint{0.000000in}{0.027778in}}{%
\pgfpathmoveto{\pgfqpoint{0.000000in}{0.000000in}}%
\pgfpathlineto{\pgfqpoint{0.000000in}{0.027778in}}%
\pgfusepath{stroke,fill}%
}%
\begin{pgfscope}%
\pgfsys@transformshift{0.952904in}{0.456901in}%
\pgfsys@useobject{currentmarker}{}%
\end{pgfscope}%
\end{pgfscope}%
\begin{pgfscope}%
\pgfpathrectangle{\pgfqpoint{0.781757in}{0.456901in}}{\pgfqpoint{5.078868in}{2.816432in}}%
\pgfusepath{clip}%
\pgfsetbuttcap%
\pgfsetroundjoin%
\pgfsetlinewidth{0.501875pt}%
\definecolor{currentstroke}{rgb}{0.501961,0.501961,0.501961}%
\pgfsetstrokecolor{currentstroke}%
\pgfsetstrokeopacity{0.400000}%
\pgfsetdash{{1.850000pt}{0.800000pt}}{0.000000pt}%
\pgfpathmoveto{\pgfqpoint{1.124051in}{0.456901in}}%
\pgfpathlineto{\pgfqpoint{1.124051in}{3.273333in}}%
\pgfusepath{stroke}%
\end{pgfscope}%
\begin{pgfscope}%
\pgfsetbuttcap%
\pgfsetroundjoin%
\definecolor{currentfill}{rgb}{1.000000,1.000000,1.000000}%
\pgfsetfillcolor{currentfill}%
\pgfsetlinewidth{0.602250pt}%
\definecolor{currentstroke}{rgb}{0.000000,0.000000,0.000000}%
\pgfsetstrokecolor{currentstroke}%
\pgfsetdash{}{0pt}%
\pgfsys@defobject{currentmarker}{\pgfqpoint{0.000000in}{0.000000in}}{\pgfqpoint{0.000000in}{0.027778in}}{%
\pgfpathmoveto{\pgfqpoint{0.000000in}{0.000000in}}%
\pgfpathlineto{\pgfqpoint{0.000000in}{0.027778in}}%
\pgfusepath{stroke,fill}%
}%
\begin{pgfscope}%
\pgfsys@transformshift{1.124051in}{0.456901in}%
\pgfsys@useobject{currentmarker}{}%
\end{pgfscope}%
\end{pgfscope}%
\begin{pgfscope}%
\pgfpathrectangle{\pgfqpoint{0.781757in}{0.456901in}}{\pgfqpoint{5.078868in}{2.816432in}}%
\pgfusepath{clip}%
\pgfsetbuttcap%
\pgfsetroundjoin%
\pgfsetlinewidth{0.501875pt}%
\definecolor{currentstroke}{rgb}{0.501961,0.501961,0.501961}%
\pgfsetstrokecolor{currentstroke}%
\pgfsetstrokeopacity{0.400000}%
\pgfsetdash{{1.850000pt}{0.800000pt}}{0.000000pt}%
\pgfpathmoveto{\pgfqpoint{1.295197in}{0.456901in}}%
\pgfpathlineto{\pgfqpoint{1.295197in}{3.273333in}}%
\pgfusepath{stroke}%
\end{pgfscope}%
\begin{pgfscope}%
\pgfsetbuttcap%
\pgfsetroundjoin%
\definecolor{currentfill}{rgb}{1.000000,1.000000,1.000000}%
\pgfsetfillcolor{currentfill}%
\pgfsetlinewidth{0.602250pt}%
\definecolor{currentstroke}{rgb}{0.000000,0.000000,0.000000}%
\pgfsetstrokecolor{currentstroke}%
\pgfsetdash{}{0pt}%
\pgfsys@defobject{currentmarker}{\pgfqpoint{0.000000in}{0.000000in}}{\pgfqpoint{0.000000in}{0.027778in}}{%
\pgfpathmoveto{\pgfqpoint{0.000000in}{0.000000in}}%
\pgfpathlineto{\pgfqpoint{0.000000in}{0.027778in}}%
\pgfusepath{stroke,fill}%
}%
\begin{pgfscope}%
\pgfsys@transformshift{1.295197in}{0.456901in}%
\pgfsys@useobject{currentmarker}{}%
\end{pgfscope}%
\end{pgfscope}%
\begin{pgfscope}%
\pgfpathrectangle{\pgfqpoint{0.781757in}{0.456901in}}{\pgfqpoint{5.078868in}{2.816432in}}%
\pgfusepath{clip}%
\pgfsetbuttcap%
\pgfsetroundjoin%
\pgfsetlinewidth{0.501875pt}%
\definecolor{currentstroke}{rgb}{0.501961,0.501961,0.501961}%
\pgfsetstrokecolor{currentstroke}%
\pgfsetstrokeopacity{0.400000}%
\pgfsetdash{{1.850000pt}{0.800000pt}}{0.000000pt}%
\pgfpathmoveto{\pgfqpoint{1.466344in}{0.456901in}}%
\pgfpathlineto{\pgfqpoint{1.466344in}{3.273333in}}%
\pgfusepath{stroke}%
\end{pgfscope}%
\begin{pgfscope}%
\pgfsetbuttcap%
\pgfsetroundjoin%
\definecolor{currentfill}{rgb}{1.000000,1.000000,1.000000}%
\pgfsetfillcolor{currentfill}%
\pgfsetlinewidth{0.602250pt}%
\definecolor{currentstroke}{rgb}{0.000000,0.000000,0.000000}%
\pgfsetstrokecolor{currentstroke}%
\pgfsetdash{}{0pt}%
\pgfsys@defobject{currentmarker}{\pgfqpoint{0.000000in}{0.000000in}}{\pgfqpoint{0.000000in}{0.027778in}}{%
\pgfpathmoveto{\pgfqpoint{0.000000in}{0.000000in}}%
\pgfpathlineto{\pgfqpoint{0.000000in}{0.027778in}}%
\pgfusepath{stroke,fill}%
}%
\begin{pgfscope}%
\pgfsys@transformshift{1.466344in}{0.456901in}%
\pgfsys@useobject{currentmarker}{}%
\end{pgfscope}%
\end{pgfscope}%
\begin{pgfscope}%
\pgfpathrectangle{\pgfqpoint{0.781757in}{0.456901in}}{\pgfqpoint{5.078868in}{2.816432in}}%
\pgfusepath{clip}%
\pgfsetbuttcap%
\pgfsetroundjoin%
\pgfsetlinewidth{0.501875pt}%
\definecolor{currentstroke}{rgb}{0.501961,0.501961,0.501961}%
\pgfsetstrokecolor{currentstroke}%
\pgfsetstrokeopacity{0.400000}%
\pgfsetdash{{1.850000pt}{0.800000pt}}{0.000000pt}%
\pgfpathmoveto{\pgfqpoint{1.808638in}{0.456901in}}%
\pgfpathlineto{\pgfqpoint{1.808638in}{3.273333in}}%
\pgfusepath{stroke}%
\end{pgfscope}%
\begin{pgfscope}%
\pgfsetbuttcap%
\pgfsetroundjoin%
\definecolor{currentfill}{rgb}{1.000000,1.000000,1.000000}%
\pgfsetfillcolor{currentfill}%
\pgfsetlinewidth{0.602250pt}%
\definecolor{currentstroke}{rgb}{0.000000,0.000000,0.000000}%
\pgfsetstrokecolor{currentstroke}%
\pgfsetdash{}{0pt}%
\pgfsys@defobject{currentmarker}{\pgfqpoint{0.000000in}{0.000000in}}{\pgfqpoint{0.000000in}{0.027778in}}{%
\pgfpathmoveto{\pgfqpoint{0.000000in}{0.000000in}}%
\pgfpathlineto{\pgfqpoint{0.000000in}{0.027778in}}%
\pgfusepath{stroke,fill}%
}%
\begin{pgfscope}%
\pgfsys@transformshift{1.808638in}{0.456901in}%
\pgfsys@useobject{currentmarker}{}%
\end{pgfscope}%
\end{pgfscope}%
\begin{pgfscope}%
\pgfpathrectangle{\pgfqpoint{0.781757in}{0.456901in}}{\pgfqpoint{5.078868in}{2.816432in}}%
\pgfusepath{clip}%
\pgfsetbuttcap%
\pgfsetroundjoin%
\pgfsetlinewidth{0.501875pt}%
\definecolor{currentstroke}{rgb}{0.501961,0.501961,0.501961}%
\pgfsetstrokecolor{currentstroke}%
\pgfsetstrokeopacity{0.400000}%
\pgfsetdash{{1.850000pt}{0.800000pt}}{0.000000pt}%
\pgfpathmoveto{\pgfqpoint{1.979785in}{0.456901in}}%
\pgfpathlineto{\pgfqpoint{1.979785in}{3.273333in}}%
\pgfusepath{stroke}%
\end{pgfscope}%
\begin{pgfscope}%
\pgfsetbuttcap%
\pgfsetroundjoin%
\definecolor{currentfill}{rgb}{1.000000,1.000000,1.000000}%
\pgfsetfillcolor{currentfill}%
\pgfsetlinewidth{0.602250pt}%
\definecolor{currentstroke}{rgb}{0.000000,0.000000,0.000000}%
\pgfsetstrokecolor{currentstroke}%
\pgfsetdash{}{0pt}%
\pgfsys@defobject{currentmarker}{\pgfqpoint{0.000000in}{0.000000in}}{\pgfqpoint{0.000000in}{0.027778in}}{%
\pgfpathmoveto{\pgfqpoint{0.000000in}{0.000000in}}%
\pgfpathlineto{\pgfqpoint{0.000000in}{0.027778in}}%
\pgfusepath{stroke,fill}%
}%
\begin{pgfscope}%
\pgfsys@transformshift{1.979785in}{0.456901in}%
\pgfsys@useobject{currentmarker}{}%
\end{pgfscope}%
\end{pgfscope}%
\begin{pgfscope}%
\pgfpathrectangle{\pgfqpoint{0.781757in}{0.456901in}}{\pgfqpoint{5.078868in}{2.816432in}}%
\pgfusepath{clip}%
\pgfsetbuttcap%
\pgfsetroundjoin%
\pgfsetlinewidth{0.501875pt}%
\definecolor{currentstroke}{rgb}{0.501961,0.501961,0.501961}%
\pgfsetstrokecolor{currentstroke}%
\pgfsetstrokeopacity{0.400000}%
\pgfsetdash{{1.850000pt}{0.800000pt}}{0.000000pt}%
\pgfpathmoveto{\pgfqpoint{2.150932in}{0.456901in}}%
\pgfpathlineto{\pgfqpoint{2.150932in}{3.273333in}}%
\pgfusepath{stroke}%
\end{pgfscope}%
\begin{pgfscope}%
\pgfsetbuttcap%
\pgfsetroundjoin%
\definecolor{currentfill}{rgb}{1.000000,1.000000,1.000000}%
\pgfsetfillcolor{currentfill}%
\pgfsetlinewidth{0.602250pt}%
\definecolor{currentstroke}{rgb}{0.000000,0.000000,0.000000}%
\pgfsetstrokecolor{currentstroke}%
\pgfsetdash{}{0pt}%
\pgfsys@defobject{currentmarker}{\pgfqpoint{0.000000in}{0.000000in}}{\pgfqpoint{0.000000in}{0.027778in}}{%
\pgfpathmoveto{\pgfqpoint{0.000000in}{0.000000in}}%
\pgfpathlineto{\pgfqpoint{0.000000in}{0.027778in}}%
\pgfusepath{stroke,fill}%
}%
\begin{pgfscope}%
\pgfsys@transformshift{2.150932in}{0.456901in}%
\pgfsys@useobject{currentmarker}{}%
\end{pgfscope}%
\end{pgfscope}%
\begin{pgfscope}%
\pgfpathrectangle{\pgfqpoint{0.781757in}{0.456901in}}{\pgfqpoint{5.078868in}{2.816432in}}%
\pgfusepath{clip}%
\pgfsetbuttcap%
\pgfsetroundjoin%
\pgfsetlinewidth{0.501875pt}%
\definecolor{currentstroke}{rgb}{0.501961,0.501961,0.501961}%
\pgfsetstrokecolor{currentstroke}%
\pgfsetstrokeopacity{0.400000}%
\pgfsetdash{{1.850000pt}{0.800000pt}}{0.000000pt}%
\pgfpathmoveto{\pgfqpoint{2.322079in}{0.456901in}}%
\pgfpathlineto{\pgfqpoint{2.322079in}{3.273333in}}%
\pgfusepath{stroke}%
\end{pgfscope}%
\begin{pgfscope}%
\pgfsetbuttcap%
\pgfsetroundjoin%
\definecolor{currentfill}{rgb}{1.000000,1.000000,1.000000}%
\pgfsetfillcolor{currentfill}%
\pgfsetlinewidth{0.602250pt}%
\definecolor{currentstroke}{rgb}{0.000000,0.000000,0.000000}%
\pgfsetstrokecolor{currentstroke}%
\pgfsetdash{}{0pt}%
\pgfsys@defobject{currentmarker}{\pgfqpoint{0.000000in}{0.000000in}}{\pgfqpoint{0.000000in}{0.027778in}}{%
\pgfpathmoveto{\pgfqpoint{0.000000in}{0.000000in}}%
\pgfpathlineto{\pgfqpoint{0.000000in}{0.027778in}}%
\pgfusepath{stroke,fill}%
}%
\begin{pgfscope}%
\pgfsys@transformshift{2.322079in}{0.456901in}%
\pgfsys@useobject{currentmarker}{}%
\end{pgfscope}%
\end{pgfscope}%
\begin{pgfscope}%
\pgfpathrectangle{\pgfqpoint{0.781757in}{0.456901in}}{\pgfqpoint{5.078868in}{2.816432in}}%
\pgfusepath{clip}%
\pgfsetbuttcap%
\pgfsetroundjoin%
\pgfsetlinewidth{0.501875pt}%
\definecolor{currentstroke}{rgb}{0.501961,0.501961,0.501961}%
\pgfsetstrokecolor{currentstroke}%
\pgfsetstrokeopacity{0.400000}%
\pgfsetdash{{1.850000pt}{0.800000pt}}{0.000000pt}%
\pgfpathmoveto{\pgfqpoint{2.664372in}{0.456901in}}%
\pgfpathlineto{\pgfqpoint{2.664372in}{3.273333in}}%
\pgfusepath{stroke}%
\end{pgfscope}%
\begin{pgfscope}%
\pgfsetbuttcap%
\pgfsetroundjoin%
\definecolor{currentfill}{rgb}{1.000000,1.000000,1.000000}%
\pgfsetfillcolor{currentfill}%
\pgfsetlinewidth{0.602250pt}%
\definecolor{currentstroke}{rgb}{0.000000,0.000000,0.000000}%
\pgfsetstrokecolor{currentstroke}%
\pgfsetdash{}{0pt}%
\pgfsys@defobject{currentmarker}{\pgfqpoint{0.000000in}{0.000000in}}{\pgfqpoint{0.000000in}{0.027778in}}{%
\pgfpathmoveto{\pgfqpoint{0.000000in}{0.000000in}}%
\pgfpathlineto{\pgfqpoint{0.000000in}{0.027778in}}%
\pgfusepath{stroke,fill}%
}%
\begin{pgfscope}%
\pgfsys@transformshift{2.664372in}{0.456901in}%
\pgfsys@useobject{currentmarker}{}%
\end{pgfscope}%
\end{pgfscope}%
\begin{pgfscope}%
\pgfpathrectangle{\pgfqpoint{0.781757in}{0.456901in}}{\pgfqpoint{5.078868in}{2.816432in}}%
\pgfusepath{clip}%
\pgfsetbuttcap%
\pgfsetroundjoin%
\pgfsetlinewidth{0.501875pt}%
\definecolor{currentstroke}{rgb}{0.501961,0.501961,0.501961}%
\pgfsetstrokecolor{currentstroke}%
\pgfsetstrokeopacity{0.400000}%
\pgfsetdash{{1.850000pt}{0.800000pt}}{0.000000pt}%
\pgfpathmoveto{\pgfqpoint{2.835519in}{0.456901in}}%
\pgfpathlineto{\pgfqpoint{2.835519in}{3.273333in}}%
\pgfusepath{stroke}%
\end{pgfscope}%
\begin{pgfscope}%
\pgfsetbuttcap%
\pgfsetroundjoin%
\definecolor{currentfill}{rgb}{1.000000,1.000000,1.000000}%
\pgfsetfillcolor{currentfill}%
\pgfsetlinewidth{0.602250pt}%
\definecolor{currentstroke}{rgb}{0.000000,0.000000,0.000000}%
\pgfsetstrokecolor{currentstroke}%
\pgfsetdash{}{0pt}%
\pgfsys@defobject{currentmarker}{\pgfqpoint{0.000000in}{0.000000in}}{\pgfqpoint{0.000000in}{0.027778in}}{%
\pgfpathmoveto{\pgfqpoint{0.000000in}{0.000000in}}%
\pgfpathlineto{\pgfqpoint{0.000000in}{0.027778in}}%
\pgfusepath{stroke,fill}%
}%
\begin{pgfscope}%
\pgfsys@transformshift{2.835519in}{0.456901in}%
\pgfsys@useobject{currentmarker}{}%
\end{pgfscope}%
\end{pgfscope}%
\begin{pgfscope}%
\pgfpathrectangle{\pgfqpoint{0.781757in}{0.456901in}}{\pgfqpoint{5.078868in}{2.816432in}}%
\pgfusepath{clip}%
\pgfsetbuttcap%
\pgfsetroundjoin%
\pgfsetlinewidth{0.501875pt}%
\definecolor{currentstroke}{rgb}{0.501961,0.501961,0.501961}%
\pgfsetstrokecolor{currentstroke}%
\pgfsetstrokeopacity{0.400000}%
\pgfsetdash{{1.850000pt}{0.800000pt}}{0.000000pt}%
\pgfpathmoveto{\pgfqpoint{3.006666in}{0.456901in}}%
\pgfpathlineto{\pgfqpoint{3.006666in}{3.273333in}}%
\pgfusepath{stroke}%
\end{pgfscope}%
\begin{pgfscope}%
\pgfsetbuttcap%
\pgfsetroundjoin%
\definecolor{currentfill}{rgb}{1.000000,1.000000,1.000000}%
\pgfsetfillcolor{currentfill}%
\pgfsetlinewidth{0.602250pt}%
\definecolor{currentstroke}{rgb}{0.000000,0.000000,0.000000}%
\pgfsetstrokecolor{currentstroke}%
\pgfsetdash{}{0pt}%
\pgfsys@defobject{currentmarker}{\pgfqpoint{0.000000in}{0.000000in}}{\pgfqpoint{0.000000in}{0.027778in}}{%
\pgfpathmoveto{\pgfqpoint{0.000000in}{0.000000in}}%
\pgfpathlineto{\pgfqpoint{0.000000in}{0.027778in}}%
\pgfusepath{stroke,fill}%
}%
\begin{pgfscope}%
\pgfsys@transformshift{3.006666in}{0.456901in}%
\pgfsys@useobject{currentmarker}{}%
\end{pgfscope}%
\end{pgfscope}%
\begin{pgfscope}%
\pgfpathrectangle{\pgfqpoint{0.781757in}{0.456901in}}{\pgfqpoint{5.078868in}{2.816432in}}%
\pgfusepath{clip}%
\pgfsetbuttcap%
\pgfsetroundjoin%
\pgfsetlinewidth{0.501875pt}%
\definecolor{currentstroke}{rgb}{0.501961,0.501961,0.501961}%
\pgfsetstrokecolor{currentstroke}%
\pgfsetstrokeopacity{0.400000}%
\pgfsetdash{{1.850000pt}{0.800000pt}}{0.000000pt}%
\pgfpathmoveto{\pgfqpoint{3.177813in}{0.456901in}}%
\pgfpathlineto{\pgfqpoint{3.177813in}{3.273333in}}%
\pgfusepath{stroke}%
\end{pgfscope}%
\begin{pgfscope}%
\pgfsetbuttcap%
\pgfsetroundjoin%
\definecolor{currentfill}{rgb}{1.000000,1.000000,1.000000}%
\pgfsetfillcolor{currentfill}%
\pgfsetlinewidth{0.602250pt}%
\definecolor{currentstroke}{rgb}{0.000000,0.000000,0.000000}%
\pgfsetstrokecolor{currentstroke}%
\pgfsetdash{}{0pt}%
\pgfsys@defobject{currentmarker}{\pgfqpoint{0.000000in}{0.000000in}}{\pgfqpoint{0.000000in}{0.027778in}}{%
\pgfpathmoveto{\pgfqpoint{0.000000in}{0.000000in}}%
\pgfpathlineto{\pgfqpoint{0.000000in}{0.027778in}}%
\pgfusepath{stroke,fill}%
}%
\begin{pgfscope}%
\pgfsys@transformshift{3.177813in}{0.456901in}%
\pgfsys@useobject{currentmarker}{}%
\end{pgfscope}%
\end{pgfscope}%
\begin{pgfscope}%
\pgfpathrectangle{\pgfqpoint{0.781757in}{0.456901in}}{\pgfqpoint{5.078868in}{2.816432in}}%
\pgfusepath{clip}%
\pgfsetbuttcap%
\pgfsetroundjoin%
\pgfsetlinewidth{0.501875pt}%
\definecolor{currentstroke}{rgb}{0.501961,0.501961,0.501961}%
\pgfsetstrokecolor{currentstroke}%
\pgfsetstrokeopacity{0.400000}%
\pgfsetdash{{1.850000pt}{0.800000pt}}{0.000000pt}%
\pgfpathmoveto{\pgfqpoint{3.520106in}{0.456901in}}%
\pgfpathlineto{\pgfqpoint{3.520106in}{3.273333in}}%
\pgfusepath{stroke}%
\end{pgfscope}%
\begin{pgfscope}%
\pgfsetbuttcap%
\pgfsetroundjoin%
\definecolor{currentfill}{rgb}{1.000000,1.000000,1.000000}%
\pgfsetfillcolor{currentfill}%
\pgfsetlinewidth{0.602250pt}%
\definecolor{currentstroke}{rgb}{0.000000,0.000000,0.000000}%
\pgfsetstrokecolor{currentstroke}%
\pgfsetdash{}{0pt}%
\pgfsys@defobject{currentmarker}{\pgfqpoint{0.000000in}{0.000000in}}{\pgfqpoint{0.000000in}{0.027778in}}{%
\pgfpathmoveto{\pgfqpoint{0.000000in}{0.000000in}}%
\pgfpathlineto{\pgfqpoint{0.000000in}{0.027778in}}%
\pgfusepath{stroke,fill}%
}%
\begin{pgfscope}%
\pgfsys@transformshift{3.520106in}{0.456901in}%
\pgfsys@useobject{currentmarker}{}%
\end{pgfscope}%
\end{pgfscope}%
\begin{pgfscope}%
\pgfpathrectangle{\pgfqpoint{0.781757in}{0.456901in}}{\pgfqpoint{5.078868in}{2.816432in}}%
\pgfusepath{clip}%
\pgfsetbuttcap%
\pgfsetroundjoin%
\pgfsetlinewidth{0.501875pt}%
\definecolor{currentstroke}{rgb}{0.501961,0.501961,0.501961}%
\pgfsetstrokecolor{currentstroke}%
\pgfsetstrokeopacity{0.400000}%
\pgfsetdash{{1.850000pt}{0.800000pt}}{0.000000pt}%
\pgfpathmoveto{\pgfqpoint{3.691253in}{0.456901in}}%
\pgfpathlineto{\pgfqpoint{3.691253in}{3.273333in}}%
\pgfusepath{stroke}%
\end{pgfscope}%
\begin{pgfscope}%
\pgfsetbuttcap%
\pgfsetroundjoin%
\definecolor{currentfill}{rgb}{1.000000,1.000000,1.000000}%
\pgfsetfillcolor{currentfill}%
\pgfsetlinewidth{0.602250pt}%
\definecolor{currentstroke}{rgb}{0.000000,0.000000,0.000000}%
\pgfsetstrokecolor{currentstroke}%
\pgfsetdash{}{0pt}%
\pgfsys@defobject{currentmarker}{\pgfqpoint{0.000000in}{0.000000in}}{\pgfqpoint{0.000000in}{0.027778in}}{%
\pgfpathmoveto{\pgfqpoint{0.000000in}{0.000000in}}%
\pgfpathlineto{\pgfqpoint{0.000000in}{0.027778in}}%
\pgfusepath{stroke,fill}%
}%
\begin{pgfscope}%
\pgfsys@transformshift{3.691253in}{0.456901in}%
\pgfsys@useobject{currentmarker}{}%
\end{pgfscope}%
\end{pgfscope}%
\begin{pgfscope}%
\pgfpathrectangle{\pgfqpoint{0.781757in}{0.456901in}}{\pgfqpoint{5.078868in}{2.816432in}}%
\pgfusepath{clip}%
\pgfsetbuttcap%
\pgfsetroundjoin%
\pgfsetlinewidth{0.501875pt}%
\definecolor{currentstroke}{rgb}{0.501961,0.501961,0.501961}%
\pgfsetstrokecolor{currentstroke}%
\pgfsetstrokeopacity{0.400000}%
\pgfsetdash{{1.850000pt}{0.800000pt}}{0.000000pt}%
\pgfpathmoveto{\pgfqpoint{3.862400in}{0.456901in}}%
\pgfpathlineto{\pgfqpoint{3.862400in}{3.273333in}}%
\pgfusepath{stroke}%
\end{pgfscope}%
\begin{pgfscope}%
\pgfsetbuttcap%
\pgfsetroundjoin%
\definecolor{currentfill}{rgb}{1.000000,1.000000,1.000000}%
\pgfsetfillcolor{currentfill}%
\pgfsetlinewidth{0.602250pt}%
\definecolor{currentstroke}{rgb}{0.000000,0.000000,0.000000}%
\pgfsetstrokecolor{currentstroke}%
\pgfsetdash{}{0pt}%
\pgfsys@defobject{currentmarker}{\pgfqpoint{0.000000in}{0.000000in}}{\pgfqpoint{0.000000in}{0.027778in}}{%
\pgfpathmoveto{\pgfqpoint{0.000000in}{0.000000in}}%
\pgfpathlineto{\pgfqpoint{0.000000in}{0.027778in}}%
\pgfusepath{stroke,fill}%
}%
\begin{pgfscope}%
\pgfsys@transformshift{3.862400in}{0.456901in}%
\pgfsys@useobject{currentmarker}{}%
\end{pgfscope}%
\end{pgfscope}%
\begin{pgfscope}%
\pgfpathrectangle{\pgfqpoint{0.781757in}{0.456901in}}{\pgfqpoint{5.078868in}{2.816432in}}%
\pgfusepath{clip}%
\pgfsetbuttcap%
\pgfsetroundjoin%
\pgfsetlinewidth{0.501875pt}%
\definecolor{currentstroke}{rgb}{0.501961,0.501961,0.501961}%
\pgfsetstrokecolor{currentstroke}%
\pgfsetstrokeopacity{0.400000}%
\pgfsetdash{{1.850000pt}{0.800000pt}}{0.000000pt}%
\pgfpathmoveto{\pgfqpoint{4.033547in}{0.456901in}}%
\pgfpathlineto{\pgfqpoint{4.033547in}{3.273333in}}%
\pgfusepath{stroke}%
\end{pgfscope}%
\begin{pgfscope}%
\pgfsetbuttcap%
\pgfsetroundjoin%
\definecolor{currentfill}{rgb}{1.000000,1.000000,1.000000}%
\pgfsetfillcolor{currentfill}%
\pgfsetlinewidth{0.602250pt}%
\definecolor{currentstroke}{rgb}{0.000000,0.000000,0.000000}%
\pgfsetstrokecolor{currentstroke}%
\pgfsetdash{}{0pt}%
\pgfsys@defobject{currentmarker}{\pgfqpoint{0.000000in}{0.000000in}}{\pgfqpoint{0.000000in}{0.027778in}}{%
\pgfpathmoveto{\pgfqpoint{0.000000in}{0.000000in}}%
\pgfpathlineto{\pgfqpoint{0.000000in}{0.027778in}}%
\pgfusepath{stroke,fill}%
}%
\begin{pgfscope}%
\pgfsys@transformshift{4.033547in}{0.456901in}%
\pgfsys@useobject{currentmarker}{}%
\end{pgfscope}%
\end{pgfscope}%
\begin{pgfscope}%
\pgfpathrectangle{\pgfqpoint{0.781757in}{0.456901in}}{\pgfqpoint{5.078868in}{2.816432in}}%
\pgfusepath{clip}%
\pgfsetbuttcap%
\pgfsetroundjoin%
\pgfsetlinewidth{0.501875pt}%
\definecolor{currentstroke}{rgb}{0.501961,0.501961,0.501961}%
\pgfsetstrokecolor{currentstroke}%
\pgfsetstrokeopacity{0.400000}%
\pgfsetdash{{1.850000pt}{0.800000pt}}{0.000000pt}%
\pgfpathmoveto{\pgfqpoint{4.375841in}{0.456901in}}%
\pgfpathlineto{\pgfqpoint{4.375841in}{3.273333in}}%
\pgfusepath{stroke}%
\end{pgfscope}%
\begin{pgfscope}%
\pgfsetbuttcap%
\pgfsetroundjoin%
\definecolor{currentfill}{rgb}{1.000000,1.000000,1.000000}%
\pgfsetfillcolor{currentfill}%
\pgfsetlinewidth{0.602250pt}%
\definecolor{currentstroke}{rgb}{0.000000,0.000000,0.000000}%
\pgfsetstrokecolor{currentstroke}%
\pgfsetdash{}{0pt}%
\pgfsys@defobject{currentmarker}{\pgfqpoint{0.000000in}{0.000000in}}{\pgfqpoint{0.000000in}{0.027778in}}{%
\pgfpathmoveto{\pgfqpoint{0.000000in}{0.000000in}}%
\pgfpathlineto{\pgfqpoint{0.000000in}{0.027778in}}%
\pgfusepath{stroke,fill}%
}%
\begin{pgfscope}%
\pgfsys@transformshift{4.375841in}{0.456901in}%
\pgfsys@useobject{currentmarker}{}%
\end{pgfscope}%
\end{pgfscope}%
\begin{pgfscope}%
\pgfpathrectangle{\pgfqpoint{0.781757in}{0.456901in}}{\pgfqpoint{5.078868in}{2.816432in}}%
\pgfusepath{clip}%
\pgfsetbuttcap%
\pgfsetroundjoin%
\pgfsetlinewidth{0.501875pt}%
\definecolor{currentstroke}{rgb}{0.501961,0.501961,0.501961}%
\pgfsetstrokecolor{currentstroke}%
\pgfsetstrokeopacity{0.400000}%
\pgfsetdash{{1.850000pt}{0.800000pt}}{0.000000pt}%
\pgfpathmoveto{\pgfqpoint{4.546987in}{0.456901in}}%
\pgfpathlineto{\pgfqpoint{4.546987in}{3.273333in}}%
\pgfusepath{stroke}%
\end{pgfscope}%
\begin{pgfscope}%
\pgfsetbuttcap%
\pgfsetroundjoin%
\definecolor{currentfill}{rgb}{1.000000,1.000000,1.000000}%
\pgfsetfillcolor{currentfill}%
\pgfsetlinewidth{0.602250pt}%
\definecolor{currentstroke}{rgb}{0.000000,0.000000,0.000000}%
\pgfsetstrokecolor{currentstroke}%
\pgfsetdash{}{0pt}%
\pgfsys@defobject{currentmarker}{\pgfqpoint{0.000000in}{0.000000in}}{\pgfqpoint{0.000000in}{0.027778in}}{%
\pgfpathmoveto{\pgfqpoint{0.000000in}{0.000000in}}%
\pgfpathlineto{\pgfqpoint{0.000000in}{0.027778in}}%
\pgfusepath{stroke,fill}%
}%
\begin{pgfscope}%
\pgfsys@transformshift{4.546987in}{0.456901in}%
\pgfsys@useobject{currentmarker}{}%
\end{pgfscope}%
\end{pgfscope}%
\begin{pgfscope}%
\pgfpathrectangle{\pgfqpoint{0.781757in}{0.456901in}}{\pgfqpoint{5.078868in}{2.816432in}}%
\pgfusepath{clip}%
\pgfsetbuttcap%
\pgfsetroundjoin%
\pgfsetlinewidth{0.501875pt}%
\definecolor{currentstroke}{rgb}{0.501961,0.501961,0.501961}%
\pgfsetstrokecolor{currentstroke}%
\pgfsetstrokeopacity{0.400000}%
\pgfsetdash{{1.850000pt}{0.800000pt}}{0.000000pt}%
\pgfpathmoveto{\pgfqpoint{4.718134in}{0.456901in}}%
\pgfpathlineto{\pgfqpoint{4.718134in}{3.273333in}}%
\pgfusepath{stroke}%
\end{pgfscope}%
\begin{pgfscope}%
\pgfsetbuttcap%
\pgfsetroundjoin%
\definecolor{currentfill}{rgb}{1.000000,1.000000,1.000000}%
\pgfsetfillcolor{currentfill}%
\pgfsetlinewidth{0.602250pt}%
\definecolor{currentstroke}{rgb}{0.000000,0.000000,0.000000}%
\pgfsetstrokecolor{currentstroke}%
\pgfsetdash{}{0pt}%
\pgfsys@defobject{currentmarker}{\pgfqpoint{0.000000in}{0.000000in}}{\pgfqpoint{0.000000in}{0.027778in}}{%
\pgfpathmoveto{\pgfqpoint{0.000000in}{0.000000in}}%
\pgfpathlineto{\pgfqpoint{0.000000in}{0.027778in}}%
\pgfusepath{stroke,fill}%
}%
\begin{pgfscope}%
\pgfsys@transformshift{4.718134in}{0.456901in}%
\pgfsys@useobject{currentmarker}{}%
\end{pgfscope}%
\end{pgfscope}%
\begin{pgfscope}%
\pgfpathrectangle{\pgfqpoint{0.781757in}{0.456901in}}{\pgfqpoint{5.078868in}{2.816432in}}%
\pgfusepath{clip}%
\pgfsetbuttcap%
\pgfsetroundjoin%
\pgfsetlinewidth{0.501875pt}%
\definecolor{currentstroke}{rgb}{0.501961,0.501961,0.501961}%
\pgfsetstrokecolor{currentstroke}%
\pgfsetstrokeopacity{0.400000}%
\pgfsetdash{{1.850000pt}{0.800000pt}}{0.000000pt}%
\pgfpathmoveto{\pgfqpoint{4.889281in}{0.456901in}}%
\pgfpathlineto{\pgfqpoint{4.889281in}{3.273333in}}%
\pgfusepath{stroke}%
\end{pgfscope}%
\begin{pgfscope}%
\pgfsetbuttcap%
\pgfsetroundjoin%
\definecolor{currentfill}{rgb}{1.000000,1.000000,1.000000}%
\pgfsetfillcolor{currentfill}%
\pgfsetlinewidth{0.602250pt}%
\definecolor{currentstroke}{rgb}{0.000000,0.000000,0.000000}%
\pgfsetstrokecolor{currentstroke}%
\pgfsetdash{}{0pt}%
\pgfsys@defobject{currentmarker}{\pgfqpoint{0.000000in}{0.000000in}}{\pgfqpoint{0.000000in}{0.027778in}}{%
\pgfpathmoveto{\pgfqpoint{0.000000in}{0.000000in}}%
\pgfpathlineto{\pgfqpoint{0.000000in}{0.027778in}}%
\pgfusepath{stroke,fill}%
}%
\begin{pgfscope}%
\pgfsys@transformshift{4.889281in}{0.456901in}%
\pgfsys@useobject{currentmarker}{}%
\end{pgfscope}%
\end{pgfscope}%
\begin{pgfscope}%
\pgfpathrectangle{\pgfqpoint{0.781757in}{0.456901in}}{\pgfqpoint{5.078868in}{2.816432in}}%
\pgfusepath{clip}%
\pgfsetbuttcap%
\pgfsetroundjoin%
\pgfsetlinewidth{0.501875pt}%
\definecolor{currentstroke}{rgb}{0.501961,0.501961,0.501961}%
\pgfsetstrokecolor{currentstroke}%
\pgfsetstrokeopacity{0.400000}%
\pgfsetdash{{1.850000pt}{0.800000pt}}{0.000000pt}%
\pgfpathmoveto{\pgfqpoint{5.231575in}{0.456901in}}%
\pgfpathlineto{\pgfqpoint{5.231575in}{3.273333in}}%
\pgfusepath{stroke}%
\end{pgfscope}%
\begin{pgfscope}%
\pgfsetbuttcap%
\pgfsetroundjoin%
\definecolor{currentfill}{rgb}{1.000000,1.000000,1.000000}%
\pgfsetfillcolor{currentfill}%
\pgfsetlinewidth{0.602250pt}%
\definecolor{currentstroke}{rgb}{0.000000,0.000000,0.000000}%
\pgfsetstrokecolor{currentstroke}%
\pgfsetdash{}{0pt}%
\pgfsys@defobject{currentmarker}{\pgfqpoint{0.000000in}{0.000000in}}{\pgfqpoint{0.000000in}{0.027778in}}{%
\pgfpathmoveto{\pgfqpoint{0.000000in}{0.000000in}}%
\pgfpathlineto{\pgfqpoint{0.000000in}{0.027778in}}%
\pgfusepath{stroke,fill}%
}%
\begin{pgfscope}%
\pgfsys@transformshift{5.231575in}{0.456901in}%
\pgfsys@useobject{currentmarker}{}%
\end{pgfscope}%
\end{pgfscope}%
\begin{pgfscope}%
\pgfpathrectangle{\pgfqpoint{0.781757in}{0.456901in}}{\pgfqpoint{5.078868in}{2.816432in}}%
\pgfusepath{clip}%
\pgfsetbuttcap%
\pgfsetroundjoin%
\pgfsetlinewidth{0.501875pt}%
\definecolor{currentstroke}{rgb}{0.501961,0.501961,0.501961}%
\pgfsetstrokecolor{currentstroke}%
\pgfsetstrokeopacity{0.400000}%
\pgfsetdash{{1.850000pt}{0.800000pt}}{0.000000pt}%
\pgfpathmoveto{\pgfqpoint{5.402722in}{0.456901in}}%
\pgfpathlineto{\pgfqpoint{5.402722in}{3.273333in}}%
\pgfusepath{stroke}%
\end{pgfscope}%
\begin{pgfscope}%
\pgfsetbuttcap%
\pgfsetroundjoin%
\definecolor{currentfill}{rgb}{1.000000,1.000000,1.000000}%
\pgfsetfillcolor{currentfill}%
\pgfsetlinewidth{0.602250pt}%
\definecolor{currentstroke}{rgb}{0.000000,0.000000,0.000000}%
\pgfsetstrokecolor{currentstroke}%
\pgfsetdash{}{0pt}%
\pgfsys@defobject{currentmarker}{\pgfqpoint{0.000000in}{0.000000in}}{\pgfqpoint{0.000000in}{0.027778in}}{%
\pgfpathmoveto{\pgfqpoint{0.000000in}{0.000000in}}%
\pgfpathlineto{\pgfqpoint{0.000000in}{0.027778in}}%
\pgfusepath{stroke,fill}%
}%
\begin{pgfscope}%
\pgfsys@transformshift{5.402722in}{0.456901in}%
\pgfsys@useobject{currentmarker}{}%
\end{pgfscope}%
\end{pgfscope}%
\begin{pgfscope}%
\pgfpathrectangle{\pgfqpoint{0.781757in}{0.456901in}}{\pgfqpoint{5.078868in}{2.816432in}}%
\pgfusepath{clip}%
\pgfsetbuttcap%
\pgfsetroundjoin%
\pgfsetlinewidth{0.501875pt}%
\definecolor{currentstroke}{rgb}{0.501961,0.501961,0.501961}%
\pgfsetstrokecolor{currentstroke}%
\pgfsetstrokeopacity{0.400000}%
\pgfsetdash{{1.850000pt}{0.800000pt}}{0.000000pt}%
\pgfpathmoveto{\pgfqpoint{5.573868in}{0.456901in}}%
\pgfpathlineto{\pgfqpoint{5.573868in}{3.273333in}}%
\pgfusepath{stroke}%
\end{pgfscope}%
\begin{pgfscope}%
\pgfsetbuttcap%
\pgfsetroundjoin%
\definecolor{currentfill}{rgb}{1.000000,1.000000,1.000000}%
\pgfsetfillcolor{currentfill}%
\pgfsetlinewidth{0.602250pt}%
\definecolor{currentstroke}{rgb}{0.000000,0.000000,0.000000}%
\pgfsetstrokecolor{currentstroke}%
\pgfsetdash{}{0pt}%
\pgfsys@defobject{currentmarker}{\pgfqpoint{0.000000in}{0.000000in}}{\pgfqpoint{0.000000in}{0.027778in}}{%
\pgfpathmoveto{\pgfqpoint{0.000000in}{0.000000in}}%
\pgfpathlineto{\pgfqpoint{0.000000in}{0.027778in}}%
\pgfusepath{stroke,fill}%
}%
\begin{pgfscope}%
\pgfsys@transformshift{5.573868in}{0.456901in}%
\pgfsys@useobject{currentmarker}{}%
\end{pgfscope}%
\end{pgfscope}%
\begin{pgfscope}%
\pgfpathrectangle{\pgfqpoint{0.781757in}{0.456901in}}{\pgfqpoint{5.078868in}{2.816432in}}%
\pgfusepath{clip}%
\pgfsetbuttcap%
\pgfsetroundjoin%
\pgfsetlinewidth{0.501875pt}%
\definecolor{currentstroke}{rgb}{0.501961,0.501961,0.501961}%
\pgfsetstrokecolor{currentstroke}%
\pgfsetstrokeopacity{0.400000}%
\pgfsetdash{{1.850000pt}{0.800000pt}}{0.000000pt}%
\pgfpathmoveto{\pgfqpoint{5.745015in}{0.456901in}}%
\pgfpathlineto{\pgfqpoint{5.745015in}{3.273333in}}%
\pgfusepath{stroke}%
\end{pgfscope}%
\begin{pgfscope}%
\pgfsetbuttcap%
\pgfsetroundjoin%
\definecolor{currentfill}{rgb}{1.000000,1.000000,1.000000}%
\pgfsetfillcolor{currentfill}%
\pgfsetlinewidth{0.602250pt}%
\definecolor{currentstroke}{rgb}{0.000000,0.000000,0.000000}%
\pgfsetstrokecolor{currentstroke}%
\pgfsetdash{}{0pt}%
\pgfsys@defobject{currentmarker}{\pgfqpoint{0.000000in}{0.000000in}}{\pgfqpoint{0.000000in}{0.027778in}}{%
\pgfpathmoveto{\pgfqpoint{0.000000in}{0.000000in}}%
\pgfpathlineto{\pgfqpoint{0.000000in}{0.027778in}}%
\pgfusepath{stroke,fill}%
}%
\begin{pgfscope}%
\pgfsys@transformshift{5.745015in}{0.456901in}%
\pgfsys@useobject{currentmarker}{}%
\end{pgfscope}%
\end{pgfscope}%
\begin{pgfscope}%
\definecolor{textcolor}{rgb}{0.000000,0.000000,0.000000}%
\pgfsetstrokecolor{textcolor}%
\pgfsetfillcolor{textcolor}%
\pgftext[x=3.321191in,y=0.201367in,,top]{\color{textcolor}{\rmfamily\fontsize{12.000000}{14.400000}\selectfont\catcode`\^=\active\def^{\ifmmode\sp\else\^{}\fi}\catcode`\%=\active\def%{\%}$t$}}%
\end{pgfscope}%
\begin{pgfscope}%
\pgfpathrectangle{\pgfqpoint{0.781757in}{0.456901in}}{\pgfqpoint{5.078868in}{2.816432in}}%
\pgfusepath{clip}%
\pgfsetbuttcap%
\pgfsetroundjoin%
\pgfsetlinewidth{0.501875pt}%
\definecolor{currentstroke}{rgb}{0.501961,0.501961,0.501961}%
\pgfsetstrokecolor{currentstroke}%
\pgfsetstrokeopacity{0.400000}%
\pgfsetdash{{1.850000pt}{0.800000pt}}{0.000000pt}%
\pgfpathmoveto{\pgfqpoint{0.781757in}{0.880996in}}%
\pgfpathlineto{\pgfqpoint{5.860625in}{0.880996in}}%
\pgfusepath{stroke}%
\end{pgfscope}%
\begin{pgfscope}%
\pgfsetbuttcap%
\pgfsetroundjoin%
\definecolor{currentfill}{rgb}{1.000000,1.000000,1.000000}%
\pgfsetfillcolor{currentfill}%
\pgfsetlinewidth{0.803000pt}%
\definecolor{currentstroke}{rgb}{0.000000,0.000000,0.000000}%
\pgfsetstrokecolor{currentstroke}%
\pgfsetdash{}{0pt}%
\pgfsys@defobject{currentmarker}{\pgfqpoint{0.000000in}{0.000000in}}{\pgfqpoint{0.048611in}{0.000000in}}{%
\pgfpathmoveto{\pgfqpoint{0.000000in}{0.000000in}}%
\pgfpathlineto{\pgfqpoint{0.048611in}{0.000000in}}%
\pgfusepath{stroke,fill}%
}%
\begin{pgfscope}%
\pgfsys@transformshift{0.781757in}{0.880996in}%
\pgfsys@useobject{currentmarker}{}%
\end{pgfscope}%
\end{pgfscope}%
\begin{pgfscope}%
\definecolor{textcolor}{rgb}{0.000000,0.000000,0.000000}%
\pgfsetstrokecolor{textcolor}%
\pgfsetfillcolor{textcolor}%
\pgftext[x=0.353325in, y=0.823135in, left, base]{\color{textcolor}{\rmfamily\fontsize{12.000000}{14.400000}\selectfont\catcode`\^=\active\def^{\ifmmode\sp\else\^{}\fi}\catcode`\%=\active\def%{\%}$\mathdefault{\ensuremath{-}1.0}$}}%
\end{pgfscope}%
\begin{pgfscope}%
\pgfpathrectangle{\pgfqpoint{0.781757in}{0.456901in}}{\pgfqpoint{5.078868in}{2.816432in}}%
\pgfusepath{clip}%
\pgfsetbuttcap%
\pgfsetroundjoin%
\pgfsetlinewidth{0.501875pt}%
\definecolor{currentstroke}{rgb}{0.501961,0.501961,0.501961}%
\pgfsetstrokecolor{currentstroke}%
\pgfsetstrokeopacity{0.400000}%
\pgfsetdash{{1.850000pt}{0.800000pt}}{0.000000pt}%
\pgfpathmoveto{\pgfqpoint{0.781757in}{1.410242in}}%
\pgfpathlineto{\pgfqpoint{5.860625in}{1.410242in}}%
\pgfusepath{stroke}%
\end{pgfscope}%
\begin{pgfscope}%
\pgfsetbuttcap%
\pgfsetroundjoin%
\definecolor{currentfill}{rgb}{1.000000,1.000000,1.000000}%
\pgfsetfillcolor{currentfill}%
\pgfsetlinewidth{0.803000pt}%
\definecolor{currentstroke}{rgb}{0.000000,0.000000,0.000000}%
\pgfsetstrokecolor{currentstroke}%
\pgfsetdash{}{0pt}%
\pgfsys@defobject{currentmarker}{\pgfqpoint{0.000000in}{0.000000in}}{\pgfqpoint{0.048611in}{0.000000in}}{%
\pgfpathmoveto{\pgfqpoint{0.000000in}{0.000000in}}%
\pgfpathlineto{\pgfqpoint{0.048611in}{0.000000in}}%
\pgfusepath{stroke,fill}%
}%
\begin{pgfscope}%
\pgfsys@transformshift{0.781757in}{1.410242in}%
\pgfsys@useobject{currentmarker}{}%
\end{pgfscope}%
\end{pgfscope}%
\begin{pgfscope}%
\definecolor{textcolor}{rgb}{0.000000,0.000000,0.000000}%
\pgfsetstrokecolor{textcolor}%
\pgfsetfillcolor{textcolor}%
\pgftext[x=0.353325in, y=1.352381in, left, base]{\color{textcolor}{\rmfamily\fontsize{12.000000}{14.400000}\selectfont\catcode`\^=\active\def^{\ifmmode\sp\else\^{}\fi}\catcode`\%=\active\def%{\%}$\mathdefault{\ensuremath{-}0.5}$}}%
\end{pgfscope}%
\begin{pgfscope}%
\pgfpathrectangle{\pgfqpoint{0.781757in}{0.456901in}}{\pgfqpoint{5.078868in}{2.816432in}}%
\pgfusepath{clip}%
\pgfsetbuttcap%
\pgfsetroundjoin%
\pgfsetlinewidth{0.501875pt}%
\definecolor{currentstroke}{rgb}{0.501961,0.501961,0.501961}%
\pgfsetstrokecolor{currentstroke}%
\pgfsetstrokeopacity{0.400000}%
\pgfsetdash{{1.850000pt}{0.800000pt}}{0.000000pt}%
\pgfpathmoveto{\pgfqpoint{0.781757in}{1.939488in}}%
\pgfpathlineto{\pgfqpoint{5.860625in}{1.939488in}}%
\pgfusepath{stroke}%
\end{pgfscope}%
\begin{pgfscope}%
\pgfsetbuttcap%
\pgfsetroundjoin%
\definecolor{currentfill}{rgb}{1.000000,1.000000,1.000000}%
\pgfsetfillcolor{currentfill}%
\pgfsetlinewidth{0.803000pt}%
\definecolor{currentstroke}{rgb}{0.000000,0.000000,0.000000}%
\pgfsetstrokecolor{currentstroke}%
\pgfsetdash{}{0pt}%
\pgfsys@defobject{currentmarker}{\pgfqpoint{0.000000in}{0.000000in}}{\pgfqpoint{0.048611in}{0.000000in}}{%
\pgfpathmoveto{\pgfqpoint{0.000000in}{0.000000in}}%
\pgfpathlineto{\pgfqpoint{0.048611in}{0.000000in}}%
\pgfusepath{stroke,fill}%
}%
\begin{pgfscope}%
\pgfsys@transformshift{0.781757in}{1.939488in}%
\pgfsys@useobject{currentmarker}{}%
\end{pgfscope}%
\end{pgfscope}%
\begin{pgfscope}%
\definecolor{textcolor}{rgb}{0.000000,0.000000,0.000000}%
\pgfsetstrokecolor{textcolor}%
\pgfsetfillcolor{textcolor}%
\pgftext[x=0.482955in, y=1.881627in, left, base]{\color{textcolor}{\rmfamily\fontsize{12.000000}{14.400000}\selectfont\catcode`\^=\active\def^{\ifmmode\sp\else\^{}\fi}\catcode`\%=\active\def%{\%}$\mathdefault{0.0}$}}%
\end{pgfscope}%
\begin{pgfscope}%
\pgfpathrectangle{\pgfqpoint{0.781757in}{0.456901in}}{\pgfqpoint{5.078868in}{2.816432in}}%
\pgfusepath{clip}%
\pgfsetbuttcap%
\pgfsetroundjoin%
\pgfsetlinewidth{0.501875pt}%
\definecolor{currentstroke}{rgb}{0.501961,0.501961,0.501961}%
\pgfsetstrokecolor{currentstroke}%
\pgfsetstrokeopacity{0.400000}%
\pgfsetdash{{1.850000pt}{0.800000pt}}{0.000000pt}%
\pgfpathmoveto{\pgfqpoint{0.781757in}{2.468734in}}%
\pgfpathlineto{\pgfqpoint{5.860625in}{2.468734in}}%
\pgfusepath{stroke}%
\end{pgfscope}%
\begin{pgfscope}%
\pgfsetbuttcap%
\pgfsetroundjoin%
\definecolor{currentfill}{rgb}{1.000000,1.000000,1.000000}%
\pgfsetfillcolor{currentfill}%
\pgfsetlinewidth{0.803000pt}%
\definecolor{currentstroke}{rgb}{0.000000,0.000000,0.000000}%
\pgfsetstrokecolor{currentstroke}%
\pgfsetdash{}{0pt}%
\pgfsys@defobject{currentmarker}{\pgfqpoint{0.000000in}{0.000000in}}{\pgfqpoint{0.048611in}{0.000000in}}{%
\pgfpathmoveto{\pgfqpoint{0.000000in}{0.000000in}}%
\pgfpathlineto{\pgfqpoint{0.048611in}{0.000000in}}%
\pgfusepath{stroke,fill}%
}%
\begin{pgfscope}%
\pgfsys@transformshift{0.781757in}{2.468734in}%
\pgfsys@useobject{currentmarker}{}%
\end{pgfscope}%
\end{pgfscope}%
\begin{pgfscope}%
\definecolor{textcolor}{rgb}{0.000000,0.000000,0.000000}%
\pgfsetstrokecolor{textcolor}%
\pgfsetfillcolor{textcolor}%
\pgftext[x=0.482955in, y=2.410873in, left, base]{\color{textcolor}{\rmfamily\fontsize{12.000000}{14.400000}\selectfont\catcode`\^=\active\def^{\ifmmode\sp\else\^{}\fi}\catcode`\%=\active\def%{\%}$\mathdefault{0.5}$}}%
\end{pgfscope}%
\begin{pgfscope}%
\pgfpathrectangle{\pgfqpoint{0.781757in}{0.456901in}}{\pgfqpoint{5.078868in}{2.816432in}}%
\pgfusepath{clip}%
\pgfsetbuttcap%
\pgfsetroundjoin%
\pgfsetlinewidth{0.501875pt}%
\definecolor{currentstroke}{rgb}{0.501961,0.501961,0.501961}%
\pgfsetstrokecolor{currentstroke}%
\pgfsetstrokeopacity{0.400000}%
\pgfsetdash{{1.850000pt}{0.800000pt}}{0.000000pt}%
\pgfpathmoveto{\pgfqpoint{0.781757in}{2.997980in}}%
\pgfpathlineto{\pgfqpoint{5.860625in}{2.997980in}}%
\pgfusepath{stroke}%
\end{pgfscope}%
\begin{pgfscope}%
\pgfsetbuttcap%
\pgfsetroundjoin%
\definecolor{currentfill}{rgb}{1.000000,1.000000,1.000000}%
\pgfsetfillcolor{currentfill}%
\pgfsetlinewidth{0.803000pt}%
\definecolor{currentstroke}{rgb}{0.000000,0.000000,0.000000}%
\pgfsetstrokecolor{currentstroke}%
\pgfsetdash{}{0pt}%
\pgfsys@defobject{currentmarker}{\pgfqpoint{0.000000in}{0.000000in}}{\pgfqpoint{0.048611in}{0.000000in}}{%
\pgfpathmoveto{\pgfqpoint{0.000000in}{0.000000in}}%
\pgfpathlineto{\pgfqpoint{0.048611in}{0.000000in}}%
\pgfusepath{stroke,fill}%
}%
\begin{pgfscope}%
\pgfsys@transformshift{0.781757in}{2.997980in}%
\pgfsys@useobject{currentmarker}{}%
\end{pgfscope}%
\end{pgfscope}%
\begin{pgfscope}%
\definecolor{textcolor}{rgb}{0.000000,0.000000,0.000000}%
\pgfsetstrokecolor{textcolor}%
\pgfsetfillcolor{textcolor}%
\pgftext[x=0.482955in, y=2.940119in, left, base]{\color{textcolor}{\rmfamily\fontsize{12.000000}{14.400000}\selectfont\catcode`\^=\active\def^{\ifmmode\sp\else\^{}\fi}\catcode`\%=\active\def%{\%}$\mathdefault{1.0}$}}%
\end{pgfscope}%
\begin{pgfscope}%
\definecolor{textcolor}{rgb}{0.000000,0.000000,0.000000}%
\pgfsetstrokecolor{textcolor}%
\pgfsetfillcolor{textcolor}%
\pgftext[x=0.297770in,y=1.865117in,,bottom,rotate=90.000000]{\color{textcolor}{\rmfamily\fontsize{12.000000}{14.400000}\selectfont\catcode`\^=\active\def^{\ifmmode\sp\else\^{}\fi}\catcode`\%=\active\def%{\%}$i_2(t)$}}%
\end{pgfscope}%
\begin{pgfscope}%
\pgfpathrectangle{\pgfqpoint{0.781757in}{0.456901in}}{\pgfqpoint{5.078868in}{2.816432in}}%
\pgfusepath{clip}%
\pgfsetrectcap%
\pgfsetroundjoin%
\pgfsetlinewidth{2.007500pt}%
\definecolor{currentstroke}{rgb}{0.000000,0.000000,0.000000}%
\pgfsetstrokecolor{currentstroke}%
\pgfsetdash{}{0pt}%
\pgfpathmoveto{\pgfqpoint{0.773200in}{2.025786in}}%
\pgfpathlineto{\pgfqpoint{0.792619in}{2.098926in}}%
\pgfpathlineto{\pgfqpoint{0.850878in}{2.338715in}}%
\pgfpathlineto{\pgfqpoint{0.860587in}{2.370677in}}%
\pgfpathlineto{\pgfqpoint{0.870297in}{2.398237in}}%
\pgfpathlineto{\pgfqpoint{0.880007in}{2.421065in}}%
\pgfpathlineto{\pgfqpoint{0.889717in}{2.439093in}}%
\pgfpathlineto{\pgfqpoint{0.899426in}{2.452508in}}%
\pgfpathlineto{\pgfqpoint{0.909136in}{2.461739in}}%
\pgfpathlineto{\pgfqpoint{0.918846in}{2.467425in}}%
\pgfpathlineto{\pgfqpoint{0.928556in}{2.470372in}}%
\pgfpathlineto{\pgfqpoint{0.947975in}{2.471781in}}%
\pgfpathlineto{\pgfqpoint{0.967395in}{2.473620in}}%
\pgfpathlineto{\pgfqpoint{0.977104in}{2.476874in}}%
\pgfpathlineto{\pgfqpoint{0.986814in}{2.482579in}}%
\pgfpathlineto{\pgfqpoint{0.996524in}{2.491179in}}%
\pgfpathlineto{\pgfqpoint{1.006234in}{2.502912in}}%
\pgfpathlineto{\pgfqpoint{1.015943in}{2.517811in}}%
\pgfpathlineto{\pgfqpoint{1.025653in}{2.535709in}}%
\pgfpathlineto{\pgfqpoint{1.045073in}{2.579024in}}%
\pgfpathlineto{\pgfqpoint{1.074202in}{2.654592in}}%
\pgfpathlineto{\pgfqpoint{1.103331in}{2.729161in}}%
\pgfpathlineto{\pgfqpoint{1.122751in}{2.773183in}}%
\pgfpathlineto{\pgfqpoint{1.142170in}{2.811960in}}%
\pgfpathlineto{\pgfqpoint{1.171299in}{2.863799in}}%
\pgfpathlineto{\pgfqpoint{1.210139in}{2.933216in}}%
\pgfpathlineto{\pgfqpoint{1.278107in}{3.063058in}}%
\pgfpathlineto{\pgfqpoint{1.297526in}{3.092024in}}%
\pgfpathlineto{\pgfqpoint{1.307236in}{3.103635in}}%
\pgfpathlineto{\pgfqpoint{1.316946in}{3.113168in}}%
\pgfpathlineto{\pgfqpoint{1.326656in}{3.120630in}}%
\pgfpathlineto{\pgfqpoint{1.336365in}{3.126135in}}%
\pgfpathlineto{\pgfqpoint{1.346075in}{3.129890in}}%
\pgfpathlineto{\pgfqpoint{1.365495in}{3.133324in}}%
\pgfpathlineto{\pgfqpoint{1.394624in}{3.133418in}}%
\pgfpathlineto{\pgfqpoint{1.423753in}{3.134081in}}%
\pgfpathlineto{\pgfqpoint{1.472302in}{3.138467in}}%
\pgfpathlineto{\pgfqpoint{1.491721in}{3.136664in}}%
\pgfpathlineto{\pgfqpoint{1.501431in}{3.134013in}}%
\pgfpathlineto{\pgfqpoint{1.511141in}{3.130002in}}%
\pgfpathlineto{\pgfqpoint{1.520851in}{3.124594in}}%
\pgfpathlineto{\pgfqpoint{1.540270in}{3.109918in}}%
\pgfpathlineto{\pgfqpoint{1.559690in}{3.091524in}}%
\pgfpathlineto{\pgfqpoint{1.588819in}{3.062642in}}%
\pgfpathlineto{\pgfqpoint{1.608238in}{3.046390in}}%
\pgfpathlineto{\pgfqpoint{1.627658in}{3.034476in}}%
\pgfpathlineto{\pgfqpoint{1.647077in}{3.026914in}}%
\pgfpathlineto{\pgfqpoint{1.676207in}{3.020294in}}%
\pgfpathlineto{\pgfqpoint{1.695626in}{3.015530in}}%
\pgfpathlineto{\pgfqpoint{1.715046in}{3.008010in}}%
\pgfpathlineto{\pgfqpoint{1.734465in}{2.996902in}}%
\pgfpathlineto{\pgfqpoint{1.763594in}{2.975566in}}%
\pgfpathlineto{\pgfqpoint{1.783014in}{2.961671in}}%
\pgfpathlineto{\pgfqpoint{1.802433in}{2.951321in}}%
\pgfpathlineto{\pgfqpoint{1.812143in}{2.948168in}}%
\pgfpathlineto{\pgfqpoint{1.821853in}{2.946558in}}%
\pgfpathlineto{\pgfqpoint{1.831563in}{2.946498in}}%
\pgfpathlineto{\pgfqpoint{1.841272in}{2.947875in}}%
\pgfpathlineto{\pgfqpoint{1.860692in}{2.953958in}}%
\pgfpathlineto{\pgfqpoint{1.899531in}{2.968875in}}%
\pgfpathlineto{\pgfqpoint{1.909241in}{2.970910in}}%
\pgfpathlineto{\pgfqpoint{1.918951in}{2.971717in}}%
\pgfpathlineto{\pgfqpoint{1.928660in}{2.971213in}}%
\pgfpathlineto{\pgfqpoint{1.948080in}{2.966605in}}%
\pgfpathlineto{\pgfqpoint{1.996629in}{2.949342in}}%
\pgfpathlineto{\pgfqpoint{2.006338in}{2.948391in}}%
\pgfpathlineto{\pgfqpoint{2.016048in}{2.949127in}}%
\pgfpathlineto{\pgfqpoint{2.025758in}{2.951699in}}%
\pgfpathlineto{\pgfqpoint{2.035468in}{2.956100in}}%
\pgfpathlineto{\pgfqpoint{2.045177in}{2.962156in}}%
\pgfpathlineto{\pgfqpoint{2.064597in}{2.977788in}}%
\pgfpathlineto{\pgfqpoint{2.084016in}{2.994562in}}%
\pgfpathlineto{\pgfqpoint{2.093726in}{3.001844in}}%
\pgfpathlineto{\pgfqpoint{2.103436in}{3.007609in}}%
\pgfpathlineto{\pgfqpoint{2.113146in}{3.011392in}}%
\pgfpathlineto{\pgfqpoint{2.122855in}{3.012880in}}%
\pgfpathlineto{\pgfqpoint{2.132565in}{3.011950in}}%
\pgfpathlineto{\pgfqpoint{2.142275in}{3.008692in}}%
\pgfpathlineto{\pgfqpoint{2.151985in}{3.003413in}}%
\pgfpathlineto{\pgfqpoint{2.171404in}{2.989015in}}%
\pgfpathlineto{\pgfqpoint{2.190824in}{2.974626in}}%
\pgfpathlineto{\pgfqpoint{2.200533in}{2.969579in}}%
\pgfpathlineto{\pgfqpoint{2.210243in}{2.967024in}}%
\pgfpathlineto{\pgfqpoint{2.219953in}{2.967571in}}%
\pgfpathlineto{\pgfqpoint{2.229663in}{2.971597in}}%
\pgfpathlineto{\pgfqpoint{2.239372in}{2.979190in}}%
\pgfpathlineto{\pgfqpoint{2.249082in}{2.990112in}}%
\pgfpathlineto{\pgfqpoint{2.268502in}{3.019285in}}%
\pgfpathlineto{\pgfqpoint{2.287921in}{3.050655in}}%
\pgfpathlineto{\pgfqpoint{2.297631in}{3.063417in}}%
\pgfpathlineto{\pgfqpoint{2.307341in}{3.071969in}}%
\pgfpathlineto{\pgfqpoint{2.317050in}{3.074621in}}%
\pgfpathlineto{\pgfqpoint{2.326760in}{3.069817in}}%
\pgfpathlineto{\pgfqpoint{2.336470in}{3.056239in}}%
\pgfpathlineto{\pgfqpoint{2.346180in}{3.032899in}}%
\pgfpathlineto{\pgfqpoint{2.355889in}{2.999218in}}%
\pgfpathlineto{\pgfqpoint{2.365599in}{2.955077in}}%
\pgfpathlineto{\pgfqpoint{2.375309in}{2.900845in}}%
\pgfpathlineto{\pgfqpoint{2.385019in}{2.837375in}}%
\pgfpathlineto{\pgfqpoint{2.404438in}{2.688325in}}%
\pgfpathlineto{\pgfqpoint{2.452987in}{2.280842in}}%
\pgfpathlineto{\pgfqpoint{2.472406in}{2.147806in}}%
\pgfpathlineto{\pgfqpoint{2.482116in}{2.093882in}}%
\pgfpathlineto{\pgfqpoint{2.491826in}{2.049174in}}%
\pgfpathlineto{\pgfqpoint{2.501536in}{2.013678in}}%
\pgfpathlineto{\pgfqpoint{2.511245in}{1.986904in}}%
\pgfpathlineto{\pgfqpoint{2.520955in}{1.967926in}}%
\pgfpathlineto{\pgfqpoint{2.530665in}{1.955447in}}%
\pgfpathlineto{\pgfqpoint{2.540375in}{1.947895in}}%
\pgfpathlineto{\pgfqpoint{2.550085in}{1.943518in}}%
\pgfpathlineto{\pgfqpoint{2.569504in}{1.937078in}}%
\pgfpathlineto{\pgfqpoint{2.579214in}{1.931621in}}%
\pgfpathlineto{\pgfqpoint{2.588924in}{1.922738in}}%
\pgfpathlineto{\pgfqpoint{2.598633in}{1.909340in}}%
\pgfpathlineto{\pgfqpoint{2.608343in}{1.890682in}}%
\pgfpathlineto{\pgfqpoint{2.618053in}{1.866389in}}%
\pgfpathlineto{\pgfqpoint{2.627763in}{1.836442in}}%
\pgfpathlineto{\pgfqpoint{2.637472in}{1.801158in}}%
\pgfpathlineto{\pgfqpoint{2.656892in}{1.717198in}}%
\pgfpathlineto{\pgfqpoint{2.686021in}{1.571777in}}%
\pgfpathlineto{\pgfqpoint{2.715150in}{1.425832in}}%
\pgfpathlineto{\pgfqpoint{2.734570in}{1.336772in}}%
\pgfpathlineto{\pgfqpoint{2.753989in}{1.255702in}}%
\pgfpathlineto{\pgfqpoint{2.783119in}{1.144541in}}%
\pgfpathlineto{\pgfqpoint{2.851087in}{0.889243in}}%
\pgfpathlineto{\pgfqpoint{2.880216in}{0.784440in}}%
\pgfpathlineto{\pgfqpoint{2.899636in}{0.724435in}}%
\pgfpathlineto{\pgfqpoint{2.909345in}{0.698733in}}%
\pgfpathlineto{\pgfqpoint{2.919055in}{0.676275in}}%
\pgfpathlineto{\pgfqpoint{2.928765in}{0.657180in}}%
\pgfpathlineto{\pgfqpoint{2.938475in}{0.641425in}}%
\pgfpathlineto{\pgfqpoint{2.948184in}{0.628851in}}%
\pgfpathlineto{\pgfqpoint{2.957894in}{0.619182in}}%
\pgfpathlineto{\pgfqpoint{2.967604in}{0.612047in}}%
\pgfpathlineto{\pgfqpoint{2.977314in}{0.607017in}}%
\pgfpathlineto{\pgfqpoint{2.987023in}{0.603639in}}%
\pgfpathlineto{\pgfqpoint{3.006443in}{0.600127in}}%
\pgfpathlineto{\pgfqpoint{3.035572in}{0.598381in}}%
\pgfpathlineto{\pgfqpoint{3.064701in}{0.598961in}}%
\pgfpathlineto{\pgfqpoint{3.084121in}{0.602710in}}%
\pgfpathlineto{\pgfqpoint{3.093831in}{0.606397in}}%
\pgfpathlineto{\pgfqpoint{3.103540in}{0.611625in}}%
\pgfpathlineto{\pgfqpoint{3.113250in}{0.618576in}}%
\pgfpathlineto{\pgfqpoint{3.122960in}{0.627343in}}%
\pgfpathlineto{\pgfqpoint{3.132670in}{0.637921in}}%
\pgfpathlineto{\pgfqpoint{3.152089in}{0.663975in}}%
\pgfpathlineto{\pgfqpoint{3.171509in}{0.694756in}}%
\pgfpathlineto{\pgfqpoint{3.210348in}{0.758367in}}%
\pgfpathlineto{\pgfqpoint{3.229767in}{0.785673in}}%
\pgfpathlineto{\pgfqpoint{3.249187in}{0.808207in}}%
\pgfpathlineto{\pgfqpoint{3.268606in}{0.826549in}}%
\pgfpathlineto{\pgfqpoint{3.326865in}{0.875499in}}%
\pgfpathlineto{\pgfqpoint{3.355994in}{0.905352in}}%
\pgfpathlineto{\pgfqpoint{3.385123in}{0.936113in}}%
\pgfpathlineto{\pgfqpoint{3.404543in}{0.953233in}}%
\pgfpathlineto{\pgfqpoint{3.414253in}{0.959900in}}%
\pgfpathlineto{\pgfqpoint{3.423962in}{0.965042in}}%
\pgfpathlineto{\pgfqpoint{3.433672in}{0.968567in}}%
\pgfpathlineto{\pgfqpoint{3.443382in}{0.970482in}}%
\pgfpathlineto{\pgfqpoint{3.453092in}{0.970897in}}%
\pgfpathlineto{\pgfqpoint{3.472511in}{0.968091in}}%
\pgfpathlineto{\pgfqpoint{3.530770in}{0.953358in}}%
\pgfpathlineto{\pgfqpoint{3.550189in}{0.952947in}}%
\pgfpathlineto{\pgfqpoint{3.579318in}{0.956339in}}%
\pgfpathlineto{\pgfqpoint{3.598738in}{0.958123in}}%
\pgfpathlineto{\pgfqpoint{3.618157in}{0.956568in}}%
\pgfpathlineto{\pgfqpoint{3.627867in}{0.953945in}}%
\pgfpathlineto{\pgfqpoint{3.637577in}{0.949949in}}%
\pgfpathlineto{\pgfqpoint{3.656996in}{0.938069in}}%
\pgfpathlineto{\pgfqpoint{3.676416in}{0.922477in}}%
\pgfpathlineto{\pgfqpoint{3.705545in}{0.898634in}}%
\pgfpathlineto{\pgfqpoint{3.715255in}{0.892228in}}%
\pgfpathlineto{\pgfqpoint{3.724965in}{0.887166in}}%
\pgfpathlineto{\pgfqpoint{3.734674in}{0.883651in}}%
\pgfpathlineto{\pgfqpoint{3.744384in}{0.881759in}}%
\pgfpathlineto{\pgfqpoint{3.754094in}{0.881428in}}%
\pgfpathlineto{\pgfqpoint{3.773513in}{0.884524in}}%
\pgfpathlineto{\pgfqpoint{3.802643in}{0.892318in}}%
\pgfpathlineto{\pgfqpoint{3.812352in}{0.893727in}}%
\pgfpathlineto{\pgfqpoint{3.822062in}{0.893741in}}%
\pgfpathlineto{\pgfqpoint{3.831772in}{0.892012in}}%
\pgfpathlineto{\pgfqpoint{3.841482in}{0.888336in}}%
\pgfpathlineto{\pgfqpoint{3.851191in}{0.882682in}}%
\pgfpathlineto{\pgfqpoint{3.860901in}{0.875207in}}%
\pgfpathlineto{\pgfqpoint{3.880321in}{0.856397in}}%
\pgfpathlineto{\pgfqpoint{3.899740in}{0.836806in}}%
\pgfpathlineto{\pgfqpoint{3.909450in}{0.828812in}}%
\pgfpathlineto{\pgfqpoint{3.919160in}{0.823249in}}%
\pgfpathlineto{\pgfqpoint{3.928870in}{0.821012in}}%
\pgfpathlineto{\pgfqpoint{3.938579in}{0.822901in}}%
\pgfpathlineto{\pgfqpoint{3.948289in}{0.829574in}}%
\pgfpathlineto{\pgfqpoint{3.957999in}{0.841493in}}%
\pgfpathlineto{\pgfqpoint{3.967709in}{0.858894in}}%
\pgfpathlineto{\pgfqpoint{3.977418in}{0.881763in}}%
\pgfpathlineto{\pgfqpoint{3.987128in}{0.909831in}}%
\pgfpathlineto{\pgfqpoint{3.996838in}{0.942581in}}%
\pgfpathlineto{\pgfqpoint{4.016257in}{1.018989in}}%
\pgfpathlineto{\pgfqpoint{4.064806in}{1.224484in}}%
\pgfpathlineto{\pgfqpoint{4.084226in}{1.290780in}}%
\pgfpathlineto{\pgfqpoint{4.093935in}{1.317672in}}%
\pgfpathlineto{\pgfqpoint{4.103645in}{1.340083in}}%
\pgfpathlineto{\pgfqpoint{4.113355in}{1.358083in}}%
\pgfpathlineto{\pgfqpoint{4.123065in}{1.371966in}}%
\pgfpathlineto{\pgfqpoint{4.132774in}{1.382224in}}%
\pgfpathlineto{\pgfqpoint{4.142484in}{1.389512in}}%
\pgfpathlineto{\pgfqpoint{4.152194in}{1.394596in}}%
\pgfpathlineto{\pgfqpoint{4.191033in}{1.409416in}}%
\pgfpathlineto{\pgfqpoint{4.200743in}{1.415491in}}%
\pgfpathlineto{\pgfqpoint{4.210452in}{1.423643in}}%
\pgfpathlineto{\pgfqpoint{4.220162in}{1.434176in}}%
\pgfpathlineto{\pgfqpoint{4.229872in}{1.447236in}}%
\pgfpathlineto{\pgfqpoint{4.239582in}{1.462807in}}%
\pgfpathlineto{\pgfqpoint{4.259001in}{1.500756in}}%
\pgfpathlineto{\pgfqpoint{4.278421in}{1.545576in}}%
\pgfpathlineto{\pgfqpoint{4.346389in}{1.711836in}}%
\pgfpathlineto{\pgfqpoint{4.375518in}{1.774298in}}%
\pgfpathlineto{\pgfqpoint{4.414357in}{1.850875in}}%
\pgfpathlineto{\pgfqpoint{4.462906in}{1.942335in}}%
\pgfpathlineto{\pgfqpoint{4.492035in}{1.991831in}}%
\pgfpathlineto{\pgfqpoint{4.511455in}{2.019938in}}%
\pgfpathlineto{\pgfqpoint{4.530874in}{2.042886in}}%
\pgfpathlineto{\pgfqpoint{4.550294in}{2.060266in}}%
\pgfpathlineto{\pgfqpoint{4.569713in}{2.072350in}}%
\pgfpathlineto{\pgfqpoint{4.589133in}{2.079951in}}%
\pgfpathlineto{\pgfqpoint{4.608552in}{2.084132in}}%
\pgfpathlineto{\pgfqpoint{4.627972in}{2.085882in}}%
\pgfpathlineto{\pgfqpoint{4.647391in}{2.085859in}}%
\pgfpathlineto{\pgfqpoint{4.666811in}{2.084281in}}%
\pgfpathlineto{\pgfqpoint{4.686230in}{2.080983in}}%
\pgfpathlineto{\pgfqpoint{4.705650in}{2.075609in}}%
\pgfpathlineto{\pgfqpoint{4.725069in}{2.067852in}}%
\pgfpathlineto{\pgfqpoint{4.744489in}{2.057663in}}%
\pgfpathlineto{\pgfqpoint{4.773618in}{2.038606in}}%
\pgfpathlineto{\pgfqpoint{4.851296in}{1.983332in}}%
\pgfpathlineto{\pgfqpoint{4.880425in}{1.966228in}}%
\pgfpathlineto{\pgfqpoint{4.919264in}{1.946286in}}%
\pgfpathlineto{\pgfqpoint{4.967813in}{1.923552in}}%
\pgfpathlineto{\pgfqpoint{4.996942in}{1.911734in}}%
\pgfpathlineto{\pgfqpoint{5.026072in}{1.902779in}}%
\pgfpathlineto{\pgfqpoint{5.045491in}{1.898860in}}%
\pgfpathlineto{\pgfqpoint{5.064911in}{1.896604in}}%
\pgfpathlineto{\pgfqpoint{5.094040in}{1.895792in}}%
\pgfpathlineto{\pgfqpoint{5.132879in}{1.897486in}}%
\pgfpathlineto{\pgfqpoint{5.191137in}{1.902055in}}%
\pgfpathlineto{\pgfqpoint{5.229977in}{1.907079in}}%
\pgfpathlineto{\pgfqpoint{5.268816in}{1.914667in}}%
\pgfpathlineto{\pgfqpoint{5.346494in}{1.931618in}}%
\pgfpathlineto{\pgfqpoint{5.385333in}{1.937491in}}%
\pgfpathlineto{\pgfqpoint{5.482430in}{1.948405in}}%
\pgfpathlineto{\pgfqpoint{5.530979in}{1.952282in}}%
\pgfpathlineto{\pgfqpoint{5.569818in}{1.952943in}}%
\pgfpathlineto{\pgfqpoint{5.618367in}{1.951428in}}%
\pgfpathlineto{\pgfqpoint{5.618367in}{1.951428in}}%
\pgfusepath{stroke}%
\end{pgfscope}%
\begin{pgfscope}%
\pgfpathrectangle{\pgfqpoint{0.781757in}{0.456901in}}{\pgfqpoint{5.078868in}{2.816432in}}%
\pgfusepath{clip}%
\pgfsetbuttcap%
\pgfsetroundjoin%
\pgfsetlinewidth{2.007500pt}%
\definecolor{currentstroke}{rgb}{0.835294,0.368627,0.000000}%
\pgfsetstrokecolor{currentstroke}%
\pgfsetdash{{7.400000pt}{3.200000pt}}{0.000000pt}%
\pgfpathmoveto{\pgfqpoint{0.773200in}{1.939488in}}%
\pgfpathlineto{\pgfqpoint{0.782909in}{1.955088in}}%
\pgfpathlineto{\pgfqpoint{0.792619in}{2.070136in}}%
\pgfpathlineto{\pgfqpoint{0.802329in}{2.159925in}}%
\pgfpathlineto{\pgfqpoint{0.812039in}{2.229399in}}%
\pgfpathlineto{\pgfqpoint{0.821748in}{2.282647in}}%
\pgfpathlineto{\pgfqpoint{0.831458in}{2.323042in}}%
\pgfpathlineto{\pgfqpoint{0.841168in}{2.353367in}}%
\pgfpathlineto{\pgfqpoint{0.850878in}{2.375914in}}%
\pgfpathlineto{\pgfqpoint{0.860587in}{2.392564in}}%
\pgfpathlineto{\pgfqpoint{0.870297in}{2.404857in}}%
\pgfpathlineto{\pgfqpoint{0.880007in}{2.414049in}}%
\pgfpathlineto{\pgfqpoint{0.899426in}{2.427012in}}%
\pgfpathlineto{\pgfqpoint{0.938265in}{2.449084in}}%
\pgfpathlineto{\pgfqpoint{0.957685in}{2.464157in}}%
\pgfpathlineto{\pgfqpoint{0.977104in}{2.483790in}}%
\pgfpathlineto{\pgfqpoint{0.996524in}{2.508377in}}%
\pgfpathlineto{\pgfqpoint{1.015943in}{2.537806in}}%
\pgfpathlineto{\pgfqpoint{1.035363in}{2.571632in}}%
\pgfpathlineto{\pgfqpoint{1.054782in}{2.609193in}}%
\pgfpathlineto{\pgfqpoint{1.083912in}{2.670803in}}%
\pgfpathlineto{\pgfqpoint{1.142170in}{2.802456in}}%
\pgfpathlineto{\pgfqpoint{1.181009in}{2.887777in}}%
\pgfpathlineto{\pgfqpoint{1.210139in}{2.946642in}}%
\pgfpathlineto{\pgfqpoint{1.239268in}{2.999215in}}%
\pgfpathlineto{\pgfqpoint{1.258687in}{3.030156in}}%
\pgfpathlineto{\pgfqpoint{1.278107in}{3.057515in}}%
\pgfpathlineto{\pgfqpoint{1.297526in}{3.081138in}}%
\pgfpathlineto{\pgfqpoint{1.316946in}{3.100953in}}%
\pgfpathlineto{\pgfqpoint{1.336365in}{3.116966in}}%
\pgfpathlineto{\pgfqpoint{1.355785in}{3.129248in}}%
\pgfpathlineto{\pgfqpoint{1.375204in}{3.137934in}}%
\pgfpathlineto{\pgfqpoint{1.394624in}{3.143212in}}%
\pgfpathlineto{\pgfqpoint{1.414043in}{3.145314in}}%
\pgfpathlineto{\pgfqpoint{1.433463in}{3.144506in}}%
\pgfpathlineto{\pgfqpoint{1.452882in}{3.141084in}}%
\pgfpathlineto{\pgfqpoint{1.472302in}{3.135362in}}%
\pgfpathlineto{\pgfqpoint{1.491721in}{3.127668in}}%
\pgfpathlineto{\pgfqpoint{1.520851in}{3.113153in}}%
\pgfpathlineto{\pgfqpoint{1.549980in}{3.096049in}}%
\pgfpathlineto{\pgfqpoint{1.608238in}{3.058200in}}%
\pgfpathlineto{\pgfqpoint{1.656787in}{3.027177in}}%
\pgfpathlineto{\pgfqpoint{1.695626in}{3.005019in}}%
\pgfpathlineto{\pgfqpoint{1.724755in}{2.990703in}}%
\pgfpathlineto{\pgfqpoint{1.753885in}{2.978697in}}%
\pgfpathlineto{\pgfqpoint{1.783014in}{2.969159in}}%
\pgfpathlineto{\pgfqpoint{1.812143in}{2.962124in}}%
\pgfpathlineto{\pgfqpoint{1.841272in}{2.957524in}}%
\pgfpathlineto{\pgfqpoint{1.870402in}{2.955199in}}%
\pgfpathlineto{\pgfqpoint{1.899531in}{2.954924in}}%
\pgfpathlineto{\pgfqpoint{1.938370in}{2.957259in}}%
\pgfpathlineto{\pgfqpoint{1.977209in}{2.962001in}}%
\pgfpathlineto{\pgfqpoint{2.035468in}{2.971943in}}%
\pgfpathlineto{\pgfqpoint{2.171404in}{2.996659in}}%
\pgfpathlineto{\pgfqpoint{2.229663in}{3.004144in}}%
\pgfpathlineto{\pgfqpoint{2.278211in}{3.008189in}}%
\pgfpathlineto{\pgfqpoint{2.326760in}{3.010251in}}%
\pgfpathlineto{\pgfqpoint{2.385019in}{3.010433in}}%
\pgfpathlineto{\pgfqpoint{2.394728in}{2.989362in}}%
\pgfpathlineto{\pgfqpoint{2.404438in}{2.756814in}}%
\pgfpathlineto{\pgfqpoint{2.414148in}{2.575169in}}%
\pgfpathlineto{\pgfqpoint{2.423858in}{2.434468in}}%
\pgfpathlineto{\pgfqpoint{2.433567in}{2.326477in}}%
\pgfpathlineto{\pgfqpoint{2.443277in}{2.244399in}}%
\pgfpathlineto{\pgfqpoint{2.452987in}{2.182627in}}%
\pgfpathlineto{\pgfqpoint{2.462697in}{2.136549in}}%
\pgfpathlineto{\pgfqpoint{2.472406in}{2.102377in}}%
\pgfpathlineto{\pgfqpoint{2.482116in}{2.077008in}}%
\pgfpathlineto{\pgfqpoint{2.491826in}{2.057912in}}%
\pgfpathlineto{\pgfqpoint{2.501536in}{2.043037in}}%
\pgfpathlineto{\pgfqpoint{2.520955in}{2.019652in}}%
\pgfpathlineto{\pgfqpoint{2.540375in}{1.997250in}}%
\pgfpathlineto{\pgfqpoint{2.559794in}{1.969970in}}%
\pgfpathlineto{\pgfqpoint{2.569504in}{1.953403in}}%
\pgfpathlineto{\pgfqpoint{2.588924in}{1.913269in}}%
\pgfpathlineto{\pgfqpoint{2.608343in}{1.863280in}}%
\pgfpathlineto{\pgfqpoint{2.627763in}{1.803654in}}%
\pgfpathlineto{\pgfqpoint{2.647182in}{1.735285in}}%
\pgfpathlineto{\pgfqpoint{2.666602in}{1.659496in}}%
\pgfpathlineto{\pgfqpoint{2.695731in}{1.535384in}}%
\pgfpathlineto{\pgfqpoint{2.753989in}{1.270706in}}%
\pgfpathlineto{\pgfqpoint{2.792828in}{1.099467in}}%
\pgfpathlineto{\pgfqpoint{2.821958in}{0.981453in}}%
\pgfpathlineto{\pgfqpoint{2.841377in}{0.909674in}}%
\pgfpathlineto{\pgfqpoint{2.860797in}{0.844319in}}%
\pgfpathlineto{\pgfqpoint{2.880216in}{0.785958in}}%
\pgfpathlineto{\pgfqpoint{2.899636in}{0.734981in}}%
\pgfpathlineto{\pgfqpoint{2.919055in}{0.691609in}}%
\pgfpathlineto{\pgfqpoint{2.938475in}{0.655900in}}%
\pgfpathlineto{\pgfqpoint{2.957894in}{0.627772in}}%
\pgfpathlineto{\pgfqpoint{2.967604in}{0.616487in}}%
\pgfpathlineto{\pgfqpoint{2.977314in}{0.607007in}}%
\pgfpathlineto{\pgfqpoint{2.987023in}{0.599288in}}%
\pgfpathlineto{\pgfqpoint{2.996733in}{0.593278in}}%
\pgfpathlineto{\pgfqpoint{3.006443in}{0.588922in}}%
\pgfpathlineto{\pgfqpoint{3.016153in}{0.586158in}}%
\pgfpathlineto{\pgfqpoint{3.025862in}{0.584920in}}%
\pgfpathlineto{\pgfqpoint{3.035572in}{0.585140in}}%
\pgfpathlineto{\pgfqpoint{3.045282in}{0.586744in}}%
\pgfpathlineto{\pgfqpoint{3.054992in}{0.589656in}}%
\pgfpathlineto{\pgfqpoint{3.074411in}{0.599090in}}%
\pgfpathlineto{\pgfqpoint{3.093831in}{0.612793in}}%
\pgfpathlineto{\pgfqpoint{3.113250in}{0.630104in}}%
\pgfpathlineto{\pgfqpoint{3.132670in}{0.650356in}}%
\pgfpathlineto{\pgfqpoint{3.161799in}{0.684822in}}%
\pgfpathlineto{\pgfqpoint{3.200638in}{0.735160in}}%
\pgfpathlineto{\pgfqpoint{3.258897in}{0.811193in}}%
\pgfpathlineto{\pgfqpoint{3.288026in}{0.846179in}}%
\pgfpathlineto{\pgfqpoint{3.317155in}{0.877637in}}%
\pgfpathlineto{\pgfqpoint{3.346284in}{0.904796in}}%
\pgfpathlineto{\pgfqpoint{3.365704in}{0.920263in}}%
\pgfpathlineto{\pgfqpoint{3.385123in}{0.933519in}}%
\pgfpathlineto{\pgfqpoint{3.404543in}{0.944534in}}%
\pgfpathlineto{\pgfqpoint{3.423962in}{0.953321in}}%
\pgfpathlineto{\pgfqpoint{3.443382in}{0.959933in}}%
\pgfpathlineto{\pgfqpoint{3.462801in}{0.964459in}}%
\pgfpathlineto{\pgfqpoint{3.482221in}{0.967015in}}%
\pgfpathlineto{\pgfqpoint{3.501640in}{0.967746in}}%
\pgfpathlineto{\pgfqpoint{3.521060in}{0.966812in}}%
\pgfpathlineto{\pgfqpoint{3.550189in}{0.962683in}}%
\pgfpathlineto{\pgfqpoint{3.579318in}{0.955849in}}%
\pgfpathlineto{\pgfqpoint{3.618157in}{0.943663in}}%
\pgfpathlineto{\pgfqpoint{3.666706in}{0.925604in}}%
\pgfpathlineto{\pgfqpoint{3.754094in}{0.892844in}}%
\pgfpathlineto{\pgfqpoint{3.792933in}{0.880567in}}%
\pgfpathlineto{\pgfqpoint{3.831772in}{0.870570in}}%
\pgfpathlineto{\pgfqpoint{3.870611in}{0.863132in}}%
\pgfpathlineto{\pgfqpoint{3.909450in}{0.858301in}}%
\pgfpathlineto{\pgfqpoint{3.948289in}{0.855935in}}%
\pgfpathlineto{\pgfqpoint{3.987128in}{0.855747in}}%
\pgfpathlineto{\pgfqpoint{3.996838in}{0.855997in}}%
\pgfpathlineto{\pgfqpoint{4.006548in}{0.861601in}}%
\pgfpathlineto{\pgfqpoint{4.016257in}{0.979342in}}%
\pgfpathlineto{\pgfqpoint{4.025967in}{1.071483in}}%
\pgfpathlineto{\pgfqpoint{4.035677in}{1.143035in}}%
\pgfpathlineto{\pgfqpoint{4.045387in}{1.198142in}}%
\pgfpathlineto{\pgfqpoint{4.055096in}{1.240223in}}%
\pgfpathlineto{\pgfqpoint{4.064806in}{1.272097in}}%
\pgfpathlineto{\pgfqpoint{4.074516in}{1.296083in}}%
\pgfpathlineto{\pgfqpoint{4.084226in}{1.314085in}}%
\pgfpathlineto{\pgfqpoint{4.093935in}{1.327661in}}%
\pgfpathlineto{\pgfqpoint{4.103645in}{1.338081in}}%
\pgfpathlineto{\pgfqpoint{4.123065in}{1.353377in}}%
\pgfpathlineto{\pgfqpoint{4.161904in}{1.379771in}}%
\pgfpathlineto{\pgfqpoint{4.181323in}{1.396868in}}%
\pgfpathlineto{\pgfqpoint{4.200743in}{1.418435in}}%
\pgfpathlineto{\pgfqpoint{4.220162in}{1.444856in}}%
\pgfpathlineto{\pgfqpoint{4.239582in}{1.476009in}}%
\pgfpathlineto{\pgfqpoint{4.259001in}{1.511434in}}%
\pgfpathlineto{\pgfqpoint{4.288130in}{1.571078in}}%
\pgfpathlineto{\pgfqpoint{4.326969in}{1.658421in}}%
\pgfpathlineto{\pgfqpoint{4.394938in}{1.813868in}}%
\pgfpathlineto{\pgfqpoint{4.424067in}{1.875003in}}%
\pgfpathlineto{\pgfqpoint{4.453196in}{1.930042in}}%
\pgfpathlineto{\pgfqpoint{4.472616in}{1.962675in}}%
\pgfpathlineto{\pgfqpoint{4.492035in}{1.991726in}}%
\pgfpathlineto{\pgfqpoint{4.511455in}{2.017011in}}%
\pgfpathlineto{\pgfqpoint{4.530874in}{2.038430in}}%
\pgfpathlineto{\pgfqpoint{4.550294in}{2.055965in}}%
\pgfpathlineto{\pgfqpoint{4.569713in}{2.069669in}}%
\pgfpathlineto{\pgfqpoint{4.589133in}{2.079659in}}%
\pgfpathlineto{\pgfqpoint{4.608552in}{2.086111in}}%
\pgfpathlineto{\pgfqpoint{4.627972in}{2.089248in}}%
\pgfpathlineto{\pgfqpoint{4.647391in}{2.089331in}}%
\pgfpathlineto{\pgfqpoint{4.666811in}{2.086653in}}%
\pgfpathlineto{\pgfqpoint{4.686230in}{2.081530in}}%
\pgfpathlineto{\pgfqpoint{4.705650in}{2.074291in}}%
\pgfpathlineto{\pgfqpoint{4.734779in}{2.060201in}}%
\pgfpathlineto{\pgfqpoint{4.763908in}{2.043240in}}%
\pgfpathlineto{\pgfqpoint{4.812457in}{2.011507in}}%
\pgfpathlineto{\pgfqpoint{4.880425in}{1.967035in}}%
\pgfpathlineto{\pgfqpoint{4.919264in}{1.944756in}}%
\pgfpathlineto{\pgfqpoint{4.948394in}{1.930445in}}%
\pgfpathlineto{\pgfqpoint{4.977523in}{1.918515in}}%
\pgfpathlineto{\pgfqpoint{5.006652in}{1.909110in}}%
\pgfpathlineto{\pgfqpoint{5.035781in}{1.902253in}}%
\pgfpathlineto{\pgfqpoint{5.064911in}{1.897861in}}%
\pgfpathlineto{\pgfqpoint{5.094040in}{1.895762in}}%
\pgfpathlineto{\pgfqpoint{5.123169in}{1.895718in}}%
\pgfpathlineto{\pgfqpoint{5.162008in}{1.898355in}}%
\pgfpathlineto{\pgfqpoint{5.200847in}{1.903371in}}%
\pgfpathlineto{\pgfqpoint{5.259106in}{1.913646in}}%
\pgfpathlineto{\pgfqpoint{5.385333in}{1.937182in}}%
\pgfpathlineto{\pgfqpoint{5.443591in}{1.945098in}}%
\pgfpathlineto{\pgfqpoint{5.492140in}{1.949482in}}%
\pgfpathlineto{\pgfqpoint{5.540689in}{1.951829in}}%
\pgfpathlineto{\pgfqpoint{5.598947in}{1.952259in}}%
\pgfpathlineto{\pgfqpoint{5.618367in}{1.951921in}}%
\pgfpathlineto{\pgfqpoint{5.618367in}{1.951921in}}%
\pgfusepath{stroke}%
\end{pgfscope}%
\begin{pgfscope}%
\pgfsetrectcap%
\pgfsetmiterjoin%
\pgfsetlinewidth{0.602250pt}%
\definecolor{currentstroke}{rgb}{0.000000,0.000000,0.000000}%
\pgfsetstrokecolor{currentstroke}%
\pgfsetdash{}{0pt}%
\pgfpathmoveto{\pgfqpoint{0.781757in}{0.456901in}}%
\pgfpathlineto{\pgfqpoint{0.781757in}{3.273333in}}%
\pgfusepath{stroke}%
\end{pgfscope}%
\begin{pgfscope}%
\pgfsetrectcap%
\pgfsetmiterjoin%
\pgfsetlinewidth{0.602250pt}%
\definecolor{currentstroke}{rgb}{0.000000,0.000000,0.000000}%
\pgfsetstrokecolor{currentstroke}%
\pgfsetdash{}{0pt}%
\pgfpathmoveto{\pgfqpoint{0.781757in}{0.456901in}}%
\pgfpathlineto{\pgfqpoint{5.860625in}{0.456901in}}%
\pgfusepath{stroke}%
\end{pgfscope}%
\begin{pgfscope}%
\pgfsetrectcap%
\pgfsetroundjoin%
\pgfsetlinewidth{2.007500pt}%
\definecolor{currentstroke}{rgb}{0.000000,0.000000,0.000000}%
\pgfsetstrokecolor{currentstroke}%
\pgfsetdash{}{0pt}%
\pgfpathmoveto{\pgfqpoint{0.116667in}{3.565000in}}%
\pgfpathlineto{\pgfqpoint{0.366667in}{3.565000in}}%
\pgfpathlineto{\pgfqpoint{0.616667in}{3.565000in}}%
\pgfusepath{stroke}%
\end{pgfscope}%
\begin{pgfscope}%
\definecolor{textcolor}{rgb}{0.000000,0.000000,0.000000}%
\pgfsetstrokecolor{textcolor}%
\pgfsetfillcolor{textcolor}%
\pgftext[x=0.750000in,y=3.506667in,left,base]{\color{textcolor}{\rmfamily\fontsize{12.000000}{14.400000}\selectfont\catcode`\^=\active\def^{\ifmmode\sp\else\^{}\fi}\catcode`\%=\active\def%{\%}Приближенный расчет по спектру}}%
\end{pgfscope}%
\begin{pgfscope}%
\pgfsetbuttcap%
\pgfsetroundjoin%
\pgfsetlinewidth{2.007500pt}%
\definecolor{currentstroke}{rgb}{0.835294,0.368627,0.000000}%
\pgfsetstrokecolor{currentstroke}%
\pgfsetdash{{7.400000pt}{3.200000pt}}{0.000000pt}%
\pgfpathmoveto{\pgfqpoint{3.507080in}{3.565000in}}%
\pgfpathlineto{\pgfqpoint{3.757080in}{3.565000in}}%
\pgfpathlineto{\pgfqpoint{4.007080in}{3.565000in}}%
\pgfusepath{stroke}%
\end{pgfscope}%
\begin{pgfscope}%
\definecolor{textcolor}{rgb}{0.000000,0.000000,0.000000}%
\pgfsetstrokecolor{textcolor}%
\pgfsetfillcolor{textcolor}%
\pgftext[x=4.140413in,y=3.506667in,left,base]{\color{textcolor}{\rmfamily\fontsize{12.000000}{14.400000}\selectfont\catcode`\^=\active\def^{\ifmmode\sp\else\^{}\fi}\catcode`\%=\active\def%{\%}Точный расчет (разд. 6). $t_\mathrm{и}=9.42$}}%
\end{pgfscope}%
\end{pgfpicture}%
\makeatother%
\endgroup%

    \caption{Сравнение точного и приближенного расчета реакции цепи для $t_{\text{и}} = \pv[.2]{n['ti2']}$}
    \label{fig:plot_approx_reaction_2}
\end{figure}

Анализ графиков показывает практически полное совпадение кривых для обеих длительностей меандра. Это подтверждает корректность расчетов спектральных характеристик и адекватность применения спектрального метода для анализа реакции цепи на сложные воздействия.